%=======================================================================
\chapter{Balance Laws}
%=======================================================================

Everything so far has been apparatus. Chapter 1 built the algebra and
the calculus, Chapters 2 and 3 described what a deformation is, Chapter
4 differentiated it, and Chapter 5 answered the single question a
balance law must ask, namely how fast the content of a moving lump of
matter changes. None of it was physics. This chapter is where the
physics arrives, in the form of four axioms: mass is conserved, linear
and rotational momentum balance against force and torque, energy
balances against power and heating, and entropy is produced but never
destroyed.

The axioms are statements about \emph{integrals} over an arbitrary part
of the body. That is the only form in which they can be stated, since
force and heating are things one applies to a region and measures on
its boundary, not things one evaluates at a point. Turning them into
differential equations is therefore the business of the chapter, and
the machinery for doing it is already complete:

\begin{itemize}\itemsep2pt
\item Reynolds' transport formula \eqr{5.15}, and its extension
      \eqr{5.31} to fields that jump across a moving surface, converts
      the left-hand side of every balance law into an integral of a
      pointwise expression plus a surface term.
\item The divergence theorem \eqr{5.1}, and its extension \eqr{5.34},
      converts every boundary flux into an integral of a pointwise
      expression plus a surface term.
\item The localization theorem, Theorem~\ref{thm:localization}, turns
      the resulting statement, valid for \emph{every} part, into two
      pointwise statements: a field equation away from the singular
      surface, and a jump condition on it.
\end{itemize}

That is the entire method, and it is used eight times in what follows
without variation. Remark~\ref{rem:distrib} of Chapter 5 predicted this
shape, and Remark~\ref{rem:jumpsign} promised that the surface terms
would emerge as the Rankine--Hugoniot jump conditions. Section
\ref{sec:ch6generic} delivers both promises at once, for a balance law
of unspecified content, and every later section is an instance of it.

Two results in the chapter are not of that kind, and they are the
deepest things in the book so far. Cauchy's stress theorem, Section
\ref{sec:ch6cauchythm}, proves that the force per unit area transmitted
across a surface depends on the surface only through its unit normal
and depends on that normal \emph{linearly}. Nothing in Chapter 5
suggested such a thing; it comes from the balance of linear momentum
alone, applied to a shrinking tetrahedron, and it is what brings the
stress tensor into existence. Noll's theorem, Section
\ref{sec:ch6noll}, proves the first half of that statement separately
and shows it was never an assumption at all.

%=======================================================================
\section{Generic forms of the balance laws}\label{sec:ch6generic}
%=======================================================================

We will show that scalar balance laws are all of the form
\begin{equation}
   \frac{\dl}{\dl t}\int_{P_{t}}\rho\varphi\,\dl v
   =\int_{P_{t}}\rho\sigma\,\dl v
   +\int_{\partial P_{t}}\ip{\vek{\pi}}\bn\,\dl a ,
   \tag{6.1}\label{eq:6.1}
\end{equation}
where $P_{t}\subset\kappa_{t}$ is the region occupied by an arbitrary
part $S$ of the body $B$ and $\bn$ is the exterior unit normal to its
boundary $\partial P_{t}$. Here $\varphi=\tilde\varphi(\bx,t)$ is the
\emph{mass} density of the quantity considered, $\sigma=
\tilde\sigma(\bx,t)$ is the rate, per unit mass, at which that quantity
is supplied, the vector field $\vek{\pi}=\tilde{\vek{\pi}}(\bx,t)$ is
its \emph{flux}, the rate of supply per unit area through the boundary,
and $\rho$ is the spatial mass density of \eqr{2.56}. The time
derivative on the left pertains to a fixed set of material points: it
is the derivative of Remark~\ref{rem:twoclocks}, which weighs the same
piece of matter twice. The three balance laws of this kind are those of
mass, energy and entropy.

Vector balance laws will be shown, with the caveat of
Remark~\ref{rem:symbolic}, to have the symbolic form
\begin{equation}
   \frac{\dl}{\dl t}\int_{P_{t}}\rho\vek{\psi}\,\dl v
   =\int_{P_{t}}\rho\br\,\dl v
   +\int_{\partial P_{t}}\bA\bn\,\dl a ,
   \tag{6.2}\label{eq:6.2}
\end{equation}
where $\vek{\psi}=\tilde{\vek{\psi}}(\bx,t)$ is the mass density of the
vector-valued quantity, $\br=\tilde\br(\bx,t)$ is its supply rate per
unit mass, and $\bA=\tilde\bA(\bx,t)$ is its tensor-valued flux. The
relevant examples are the balances of linear and rotational momentum,
that is, the equations of motion.

\begin{remark}[why a balance law has exactly these three
terms]\label{rem:ch6three}
The form of \eqr{6.1} is not a convention chosen for convenience. It is
the statement that the quantity in question is \emph{conserved except
for what is added}, and there are only two ways of adding to the
contents of a region: through its interior, and through its boundary.

The first term on the right is production in the bulk: radioactive
heating, a chemical source, a body force doing work. It is written
$\rho\sigma$ rather than as a volume density because supplies of this
kind are almost always proportional to how much matter is present, so
the per-unit-mass rate $\sigma$ is the quantity a constitutive equation
will later have to specify. The second term is transport across the
boundary. Its integrand is a rate per unit area, and what makes it a
\emph{flux} rather than a mere boundary datum is that it is determined
by the orientation of the surface it crosses.

Two features of the boundary term deserve to be flagged now, because
each will be argued for later rather than assumed.

\emph{Linearity in $\bn$.} Equation \eqr{6.1} writes the flux as
$\ip{\vek{\pi}}\bn$, which is a linear function of the unit normal.
For a quantity like heat this is a hypothesis, and it is the exact
analogue of Cauchy's hypothesis for the traction. It will be
\emph{proved} for the traction in Section~\ref{sec:ch6cauchythm}, and
the proof will then be transplanted to the heat flux at \eqr{6.166}.

\emph{Dependence on the surface only through $\bn$.} Before one can ask
whether the flux is linear in the normal one must know that it depends
on the surface only through the normal, and not on, say, the curvature.
That is Noll's theorem, Section~\ref{sec:ch6noll}.

No term proportional to the surface area alone, and no term involving
second derivatives of the boundary, appears in \eqr{6.1}. Admitting
such terms is possible and leads to theories of a different kind; the
classical theory is the one that does not.
\end{remark}

Before turning to particular balance laws we pause to derive the local
forms of the generic balances \eqr{6.1} and \eqr{6.2}, holding at
points $\bx\in\kappa_{t}$ and on any singular surface that may be
present. Doing the work once, on an unspecified $\varphi$, saves doing
it five times.

\subsection*{The scalar case}

Start from the left-hand side of \eqr{6.1}. It is the derivative of a
volume integral over a material region, which is exactly what
Chapter 5 computed, so substitute $\rho\varphi$ in place of $\varphi$
in Reynolds' formula for a field that jumps across a singular surface,
\eqr{5.31}:
\[
   \frac{\dl}{\dl t}\int_{P_{t}}\rho\varphi\,\dl v
   =\int_{P_{t}}(\rho\varphi)'\,\dl v
   +\int_{\partial P_{t}}\rho\varphi\,\ip\bv\bn\,\dl a
   -\int_{s}\jump{\rho\varphi}\,u\,\dl a ,
\]
in which $s$ is the singular surface, $u$ its normal speed, and $\bn$
on $s$ is directed from $V^{-}$ into $V^{+}$ as in \eqr{5.28}.

The middle term is a boundary flux of the vector field $\rho\varphi\bv$
and is therefore the business of the divergence theorem for
discontinuous fields, \eqr{5.34}, with $\ba=\rho\varphi\bv$:
\[
   \int_{\partial P_{t}}\ip{\rho\varphi\bv}\bn\,\dl a
   =\int_{P_{t}}\dvg(\rho\varphi\bv)\,\dl v
   +\int_{s}\ip{\jump{\rho\varphi\bv}}\bn\,\dl a .
\]
Substituting, and collecting the two integrals over $s$,
\begin{equation}
\begin{aligned}
   \frac{\dl}{\dl t}\int_{P_{t}}\rho\varphi\,\dl v
   &=\int_{P_{t}}(\rho\varphi)'\,\dl v
    +\int_{\partial P_{t}}\rho\varphi\,\ip\bv\bn\,\dl a
    -\int_{s}\jump{\rho\varphi}u\,\dl a\\
   &=\int_{P_{t}}\big\{(\rho\varphi)'+\dvg(\rho\varphi\bv)\big\}\,\dl v
    +\int_{s}\ip{\jump{\rho\varphi(\bv-\bu)}}\bn\,\dl a ,
\end{aligned}
   \tag{6.3}\label{eq:6.3}
\end{equation}
where $\bu$ is the vector defined on $s$ by $\ip\bu\bn=u$, the speed of
the singular surface, exactly as in \eqr{5.27}.

The last step is the one place in the section where something is moved
across a jump bracket, so it is done in full. We have
\[
   \ip{\jump{\rho\varphi\bv}}\bn-\jump{\rho\varphi}u .
\]
Both $u$ and $\bn$ belong to the \emph{surface}, not to the material:
$s$ has one normal and advances at one speed at each of its points, and
neither of them takes two values there. A jump bracket subtracts the
two limiting values of a field approached from the two sides, so a
quantity with only one value passes through the bracket untouched:
$\jump{\rho\varphi}u=\jump{\rho\varphi u}$, and $\ip{\jump\ba}\bn$ may
be written $\jump{\ip\ba\bn}$. Writing $u=\ip\bu\bn$ and using linearity of the
inner product in its first slot,
\[
   \ip{\jump{\rho\varphi\bv}}\bn-\jump{\rho\varphi}\ip\bu\bn
   =\ip{\jump{\rho\varphi\bv}}\bn-\ip{\jump{\rho\varphi\bu}}\bn
   =\ip{\jump{\rho\varphi(\bv-\bu)}}\bn ,
\]
which is the integrand in \eqr{6.3}$_{2}$. The combination $\bv-\bu$ is
the velocity of the material relative to the surface, and only its
normal component survives the inner product:
\[
   \ip{(\bv-\bu)}\bn=\ip\bv\bn-\ip\bu\bn=\ip\bv\bn-u ,
\]
which is the quantity Problem 3 of Chapter 5 identified as the one with
physical content.

Now bring in the right-hand side of \eqr{6.1}, treating its boundary
term by \eqr{5.34} once more, this time with $\ba=\vek{\pi}$:
\[
   \int_{\partial P_{t}}\ip{\vek{\pi}}\bn\,\dl a
   =\int_{P_{t}}\dvg\vek{\pi}\,\dl v
   +\int_{s}\ip{\jump{\vek{\pi}}}\bn\,\dl a .
\]
Equating \eqr{6.3}$_{2}$ to $\int\rho\sigma\,\dl v+\int\dvg\vek{\pi}\,
\dl v+\int_{s}\ip{\jump{\vek{\pi}}}\bn\,\dl a$ and moving everything to
one side,
\begin{equation}
   \int_{P_{t}}\big\{(\rho\varphi)'+\dvg(\rho\varphi\bv-\vek{\pi})
   -\rho\sigma\big\}\,\dl v
   +\int_{s}\ip{\jump{\rho\varphi(\bv-\bu)-\vek{\pi}}}\bn\,\dl a=0 .
   \tag{6.4}\label{eq:6.4}
\end{equation}

\subsection*{Two localizations, not one}

Equation \eqr{6.4} holds for every part $S$ of the body. It contains
two integrals over two different sets, and the standard localization
argument extracts one pointwise statement from each. The extraction is
done in two stages, and the order matters.

\emph{Stage one: parts that miss the singular surface.} Apply \eqr{6.4}
to an arbitrary part $P_{t}$ that contains no piece of a singular
surface, so that $s=\emptyset$ and the second integral is absent
(Figure~\ref{fig:ch6parts}). What remains is
\[
   \int_{P_{t}}\big\{(\rho\varphi)'+\dvg(\rho\varphi\bv-\vek{\pi})
   -\rho\sigma\big\}\,\dl v=0
\]
for every such $P_{t}$. Assuming, as always, that the integrand is
continuous there, Theorem~\ref{thm:localization} applies and gives
\begin{equation}
   (\rho\varphi)'+\dvg(\rho\varphi\bv)=\dvg\vek{\pi}+\rho\sigma
   \tag{6.5}\label{eq:6.5}
\end{equation}
at all points $\bx\in\kappa_{t}$ that do not belong to a singular
surface. The argument is the one recalled in Section~\ref{sec:piolaid}:
if the integrand were non-zero at a point it would keep its sign
throughout a small ball around that point, and integrating over that
ball alone would contradict the hypothesis. Freedom to choose the
region is what carries it.

\emph{Stage two: parts that meet the singular surface.} Now let
$P_{t}$ contain a piece $s$ of a singular surface. It is cut by $s$
into subvolumes $V^{\pm}$, and in the interior of each of them the
fields are smooth, so \eqr{6.5} holds there by stage one. Therefore the
volume integral in \eqr{6.4} vanishes, whatever $P_{t}$ is, and
\eqr{6.4} collapses to
\[
   \int_{s}\ip{\jump{\rho\varphi(\bv-\bu)-\vek{\pi}}}\bn\,\dl a=0 .
\]
The part $P_{t}$ is otherwise arbitrary, so its intersection $s$ with
the singular surface is an arbitrary piece of that surface, and a
version of the localization theorem for surface integrals, together
with continuity of the integrand on $s$, yields the \emph{jump
condition}
\begin{equation}
   \ip{\jump{\rho\varphi(\bv-\bu)-\vek{\pi}}}\bn=0
   \tag{6.6}\label{eq:6.6}
\end{equation}
at every point of a singular surface.

\begin{remark}[why the two statements are independent]\label{rem:ch6two}
It is tempting to think that \eqr{6.6} is contained in \eqr{6.5}, or
follows from it by a limit. It does not, and the reason is the reason
Chapter 5 needed Section~\ref{sec:reydisc} at all. Equation \eqr{6.5}
is a statement about derivatives, and across $s$ the fields have no
derivatives: the classical divergence of a field that jumps is not
defined there. What \eqr{6.4} says is that the balance, read
distributionally, has a bulk part and a surface part, and that both
must vanish separately because a part of the body may be chosen to
contain the surface or to avoid it. That freedom of choice is exactly
what Remark~\ref{rem:distrib} anticipated.

The two-stage order is also not decorative. Had we tried to localize
\eqr{6.4} in one step we would have concluded nothing, since a volume
integral and a surface integral summing to zero does not by itself
force either to vanish: they might cancel. What forces the split is the
existence of parts with $s=\emptyset$, which kills the surface term
while leaving the volume term arbitrary.
\end{remark}

The jump condition \eqr{6.6} reduces to
\begin{equation}
   \ip{\jump{\vek{\pi}}}\bn=0
   \tag{6.7}\label{eq:6.7}
\end{equation}
if the singular surface is a \emph{material} surface, that is, if it
moves with the material, $u=\ip\bv\bn$ with $\jump\bv=\bze$. Then
$\bv-\bu$ has vanishing normal component on each side and the convected
term drops out. In fact \eqr{6.6} places no restriction whatever on the
tangential discontinuity $(\I-\bn\otimes\bn)\jump\bv$, since only the
normal component of $\bv-\bu$ enters, so \eqr{6.7} survives under the
weaker hypotheses $u=\ip\bv\bn$ with $\ip{\jump\bv}\bn=0$. A vortex
sheet, across which the tangential velocity jumps while no material
crosses the surface, is admitted.

\bookfig{1}\begin{figure}[htbp]\centering
\begin{tikzpicture}[scale=1.0,
   ar/.style={-{Stealth[length=2.2mm]},thick}]
% ---------------- the body, with a singular surface running through it
\draw[thick] plot[smooth cycle,tension=0.80] coordinates
  {(0,2.05) (2.30,1.45) (2.75,-0.55) (0.95,-2.15) (-1.55,-1.75)
   (-2.45,0.25)};
\node[font=\small,anchor=west] at (2.55,1.55) {$\kappa_{t}$};
% the singular surface
\draw[very thick,red!65!black]
   (-2.30,-0.70) .. controls (-0.60,-0.15) and (0.75,0.45) .. (2.62,0.20);
\node[font=\small,text=red!65!black,anchor=west] at (2.30,0.62) {$s$};
% part 1: does not meet s
\draw[thick,fill=blue!7] (-1.30,-1.45) circle (0.62);
\node[font=\footnotesize] at (-1.30,-1.45) {$P_{t}^{(1)}$};
% part 2: cut by s
\draw[thick,fill=blue!7] plot[smooth cycle,tension=0.85] coordinates
  {(0.75,0.95) (1.65,0.45) (1.35,-0.55) (0.35,-0.60) (0.05,0.35)};
\node[font=\footnotesize] at (1.42,0.78) {$P_{t}^{(2)}$};
\node[font=\footnotesize] at (0.72,0.62) {$V^{+}$};
\node[font=\footnotesize] at (0.86,-0.32) {$V^{-}$};
% normal and speed on s inside part 2
\coordinate (Q) at (0.98,0.16);
\draw[ar] (Q) -- ++(-0.12,0.60) node[above,font=\small] {$\bn$};
\fill (Q) circle (1.1pt);
\draw[ar,red!65!black] (Q) -- ++(-0.07,0.36)
   node[right=1pt,font=\footnotesize,text=red!65!black] {$u$};
\end{tikzpicture}
\caption{Parts of a body with and without a singular surface. Applying
\eqr{6.4} to $P_{t}^{(1)}$, which misses $s$, kills the surface term
and localizes to the field equation \eqr{6.5}. Applying it to
$P_{t}^{(2)}$, which is cut by $s$ into $V^{\pm}$, then kills the
volume term, because \eqr{6.5} already holds in each half, and what is
left localizes on $s$ to the jump condition \eqr{6.6}. Both parts are
arbitrary, which is why both conclusions follow. The normal $\bn$ on
$s$ points from $V^{-}$ into $V^{+}$, and $s$ advances at speed $u$
along it, the convention of \eqr{5.28}.}
\label{fig:ch6parts}
\end{figure}

\subsection*{The vector case}

Turning to \eqr{6.2}, we proceed exactly as above, and as explained in
Section~\ref{sec:reydisc}, to cast \eqr{6.2} in the symbolic form
\begin{equation}
   \int_{P_{t}}\big\{(\rho\vek\psi)'+\dvg(\rho\vek\psi\otimes\bv-\bA)
   -\rho\br\big\}\,\dl v
   +\int_{s}\jump{\rho\vek\psi\otimes(\bv-\bu)-\bA}\bn\,\dl a=\bze ,
   \tag{6.8}\label{eq:6.8}
\end{equation}
use having been made of the relevant extensions of \eqr{5.31} and
\eqr{5.34}, namely
\begin{equation}
\begin{aligned}
   \frac{\dl}{\dl t}\int_{P_{t}}\rho\vek{\psi}\,\dl v
   &=\int_{P_{t}}(\rho\vek{\psi})'\,\dl v
    +\int_{\partial P_{t}}(\rho\vek{\psi}\otimes\bv)\bn\,\dl a
    -\int_{s}\jump{\rho\vek{\psi}}u\,\dl a
   \qquad\text{and}\\[1mm]
   \int_{\partial P_{t}}\bA\bn\,\dl a
   &=\int_{P_{t}}\dvg\bA\,\dl v+\int_{s}\jump\bA\bn\,\dl a ,
\end{aligned}
   \tag{6.9}\label{eq:6.9}
\end{equation}
respectively.

Both lines of \eqr{6.9} are obtained in the way described after
\eqr{5.31} and again at \eqr{5.18}: apply the scalar statement to each
Cartesian component on a fixed orthonormal basis, multiply by $\be_{i}$
and sum. The caveat of Remark~\ref{rem:symbolic} is in force
throughout, so each display is shorthand for three component
statements. In the first of them the surface integrand appeared as
$\rho\psi_{i}\ip\bv\bn$ in components, and
$\rho\psi_{i}(\ip\bv\bn)\be_{i}=(\ip\bv\bn)\rho\vek\psi
=(\rho\vek\psi\otimes\bv)\bn$ by the definition of the tensor product,
Definition~\ref{def:tp}, which is the step already taken between
\eqr{5.17} and \eqr{5.18}.

The passage from $\jump{\rho\vek\psi\otimes\bv}\bn
-\jump{\rho\vek\psi}u$ to $\jump{\rho\vek\psi\otimes(\bv-\bu)}\bn$ in
\eqr{6.8} is the tensor-product version of the manipulation performed
after \eqr{6.3}: with $u=\ip\bu\bn$ single-valued on $s$,
\[
   \jump{\rho\vek\psi}\ip\bu\bn
   =\jump{\rho\vek\psi\,\ip\bu\bn}
   =\jump{(\rho\vek\psi\otimes\bu)\bn}
   =\jump{\rho\vek\psi\otimes\bu}\bn ,
\]
and the two tensor products then subtract by bilinearity of $\otimes$.

The two localizations run as before. The local equation holding at
points of $\kappa_{t}$ off a singular surface is
\begin{equation}
   (\rho\vek\psi)'+\dvg(\rho\vek\psi\otimes\bv)=\dvg\bA+\rho\br ,
   \tag{6.10}\label{eq:6.10}
\end{equation}
and the jump condition holding at every point of a singular surface is
\begin{equation}
   \jump{\rho\vek\psi\otimes(\bv-\bu)-\bA}\bn=\bze ,
   \tag{6.11}\label{eq:6.11}
\end{equation}
reducing to
\begin{equation}
   \jump\bA\bn=\bze
   \tag{6.12}\label{eq:6.12}
\end{equation}
under the conditions stated after \eqr{6.7}.

The foregoing statements are all expressed in the spatial description.
Deriving their referential counterparts is a straightforward matter and
is done in Problem 1, where the hint is the one Chapter 5 already
supplied: in $\kappa$ nothing moves, so $\dot\bX=\bze$ and every term
built from the material velocity disappears.

%=======================================================================
\section{Conservation of mass}\label{sec:ch6mass}
%=======================================================================

Section~\ref{sec:mass} already introduced the idea that the mass of a
fixed set $S\subset B$ of material points is conserved in the course of
any motion the body may undergo. As an axiom:
\begin{equation}
   \frac{\dl}{\dl t}\int_{P_{t}}\rho(\bx,t)\,\dl v=0 .
   \tag{6.13}\label{eq:6.13}
\end{equation}
This is of the form \eqr{6.1} with
\[
   \varphi=1,\qquad \sigma=0,\qquad \vek\pi=\bze .
\]
Each of the three identifications says something. Taking $\varphi=1$
says that the quantity being balanced is mass itself, whose mass
density is one unit of mass per unit mass. Taking $\sigma=0$ says that
mass is not created in the interior. Taking $\vek\pi=\bze$ says that
mass does not enter through the boundary except by being carried there
by the material, since the convective transport $\rho\ip\bv\bn$ is
already accounted for on the left-hand side by the fact that the region
$P_{t}$ moves with the material. A non-zero $\vek\pi$ would describe
diffusion, and it is precisely the term one must restore in a mixture.

The two local statements now follow by substitution into \eqr{6.5} and
\eqr{6.6}. Since $\dvg\vek\pi=0$ and $\rho\sigma=0$, the field
equation is
\begin{equation}
   \rho'+\dvg(\rho\bv)=0 ,
   \tag{6.14}\label{eq:6.14}
\end{equation}
holding wherever the fields are smooth, and the jump condition is
\begin{equation}
   \ip{\jump{\rho(\bv-\bu)}}\bn=0 ,
   \tag{6.15}\label{eq:6.15}
\end{equation}
operative on singular surfaces.

\begin{remark}[the first Rankine--Hugoniot condition]\label{rem:ch6rh1}
Equation \eqr{6.15} is the promise of Remark~\ref{rem:jumpsign} coming
due. Write it out. With $\bn$ and $u$ single-valued on $s$ it says
\[
   \rho^{+}\big(\ip{\bv^{+}}\bn-u\big)
   =\rho^{-}\big(\ip{\bv^{-}}\bn-u\big),
\]
so the scalar
\[
   m^{*}:=\rho\big(\ip\bv\bn-u\big)
\]
takes the same value on both sides of the surface. Its meaning is
direct: $\ip\bv\bn-u$ is the normal speed of the material relative to
the surface, so $m^{*}$ is the mass crossing unit area of $s$ per unit
time. Conservation of mass says that whatever flows in on one side
flows out on the other, which is the only thing it could say.

Two limiting cases are worth separating. If $m^{*}=0$, no material
crosses $s$: the surface is material, $u=\ip\bv\bn$ on both sides, and
$\rho$ may jump freely. That is a contact discontinuity, or a phase
boundary at rest relative to the material. If $m^{*}\neq0$, material is
being processed by the surface, which is then a shock: neither $\rho$
nor $\bv$ need be continuous, but their combination $m^{*}$ is. Every
further jump condition in this chapter will be written most naturally
in terms of $m^{*}$, because $m^{*}$ passes through jump brackets.
\end{remark}

Equation \eqr{6.13} must be modified in the presence of diffusion. In
that case $\sigma$ is the supply of the diffusing species per unit mass
of the mixture and $\vek\pi$ is the mass flux of that species through
the boundary. Diffusion is not treated in this book.

\subsection*{Three equivalent forms}

Equation \eqr{6.14} is stated in the spatial clock. Trading it for the
material one is a matter of the product rule. By Problem 24(c) of
Chapter 1,
\begin{equation}
   \dvg(\rho\bv)=\ip{\grad\rho}\bv+\rho\dvg\bv ,
   \tag{6.16}\label{eq:6.16}
\end{equation}
and $\dvg\bv=\tr\bL$ with $\bL=\grad\bv$ the spatial velocity gradient
of \eqr{4.1}, by Proposition~\ref{prop:divv} and
Definition~\ref{def:divv}. Substituting \eqr{6.16} into \eqr{6.14} and
recognising the first two terms as a material derivative through
\eqr{2.20}, $\dot\rho=\rho'+\ip{\grad\rho}\bv$,
\[
   0=\rho'+\ip{\grad\rho}\bv+\rho\dvg\bv=\dot\rho+\rho\dvg\bv .
\]
By Problem 2 of Chapter 4, $\dot J=J\tr\bL$, so $\dvg\bv=\dot J/J$ and
\begin{equation}
   0=\dot\rho+\rho\dvg\bv=\dot\rho+\rho\,(\dot J/J).
   \tag{6.17}\label{eq:6.17}
\end{equation}
Multiplying \eqr{6.17} through by the positive number $J$ assembles the
two terms into a single product rule,
\[
   0=J\dot\rho+\rho\dot J=(\rho J)^{\textstyle\cdot},
\]
so $\rho J$ is constant in time along each material point. With
reference to \eqr{2.59}, $\rho_{\kappa}=\rho J$, this is equivalent to
the vanishing of the material derivative of the referential mass
density:
\begin{equation}
   \dot\rho_{\kappa}=0 .
   \tag{6.18}\label{eq:6.18}
\end{equation}
Recalling that $\dot\rho_{\kappa}=\partial\rho_{\kappa}(\bX,t)/
\partial t$, since the material derivative in the referential
description is simply the partial time derivative at fixed $\bX$, it
follows that the referential density, expressed as a function of $\bX$
and $t$, is a function of $\bX$ alone. We indicate this by writing
$\rho_{\kappa}(\bX)$.

An alternative derivation of the same result is obtained by
differentiating \eqr{2.56}, assuming $P\subset\kappa$ to be a region in
which $\rho_{\kappa}$ and $\dot\rho_{\kappa}$ satisfy the continuity
conditions stated after \eqr{5.13}. Because $P$ does not depend on
time, the derivative passes under the integral sign at once, with no
transport formula needed:
\begin{equation}
   0=\frac{\dl}{\dl t}\int_{P}\rho_{\kappa}\,\dl V
    =\int_{P}\dot\rho_{\kappa}\,\dl V .
   \tag{6.19}\label{eq:6.19}
\end{equation}
Choosing $P$ to be otherwise arbitrary, we localize as usual by
Theorem~\ref{thm:localization} to arrive at \eqr{6.18} directly.

\begin{remark}[which form to use, and why the third is the
best]\label{rem:ch6massforms}
The three statements
\[
   \rho'+\dvg(\rho\bv)=0,\qquad
   \dot\rho+\rho\dvg\bv=0,\qquad
   \dot\rho_{\kappa}=0
\]
are the same axiom in the three clocks of Remark~\ref{rem:twoclocks},
and \eqr{6.19} shows how cheap the third one is: pull back to a region
that does not move, and conservation of mass becomes the observation
that a time-independent integral has a vanishing derivative.

The third form is the one to remember, and the reason is stated at
\eqr{6.136} below. In the referential description $\rho_{\kappa}(\bX)$
is \emph{given data}, part of the description of the body, and the mass
balance is then not an equation to be solved at all: it has already
been solved. In the spatial description $\rho(\bx,t)$ is an unknown
field obeying the partial differential equation \eqr{6.14}, coupled to
the velocity. One equation has been traded for none, and that is the
whole advantage of the referential formulation of the equations of
motion.
\end{remark}

\subsection*{What mass conservation does to the generic balances}

The results just derived effect a significant simplification of the
generic local balances \eqr{6.5} and \eqr{6.10}. Concerning the former,
work on its left-hand side. Expand the divergence by Problem 24(c) of
Chapter 1, applied to the scalar field $\rho\varphi$:
\begin{equation}
\begin{aligned}
   (\rho\varphi)'+\dvg(\rho\varphi\bv)
   &=(\rho\varphi)'+\ip{\grad(\rho\varphi)}\bv+\rho\varphi\dvg\bv\\
   &=(\rho\varphi)^{\textstyle\cdot}+\rho\varphi\dvg\bv\\
   &=\big(\dot\rho+\rho\dvg\bv\big)\varphi+\rho\dot\varphi\\
   &=\rho\dot\varphi ,
\end{aligned}
   \tag{6.20}\label{eq:6.20}
\end{equation}
where conservation of mass in the form \eqr{6.17}$_{1}$ has been
applied in the last step. The second line used \eqr{2.20} on the field
$\rho\varphi$; the third expanded
$(\rho\varphi)^{\textstyle\cdot}=\dot\rho\varphi+\rho\dot\varphi$ by
the product rule and then grouped the two terms carrying $\varphi$.
That grouping is the whole point of the calculation: it isolates
exactly the combination that mass conservation annihilates.
Accordingly, \eqr{6.5} simplifies to
\begin{equation}
   \rho\dot\varphi=\dvg\vek\pi+\rho\sigma .
   \tag{6.21}\label{eq:6.21}
\end{equation}

In the same way, for the vector case, we use Problem 24(g) of
Chapter 1, which is Proposition~\ref{prop:divprod} in the form
\begin{equation}
   \dvg(\rho\vek\psi\otimes\bv)
   =\big[\grad(\rho\vek\psi)\big]\bv+(\rho\dvg\bv)\vek\psi ,
   \tag{6.22}\label{eq:6.22}
\end{equation}
together with conservation of mass. The steps are the same three:
\[
\begin{aligned}
   (\rho\vek\psi)'+\dvg(\rho\vek\psi\otimes\bv)
   &=(\rho\vek\psi)'+\big[\grad(\rho\vek\psi)\big]\bv
     +\rho(\dvg\bv)\vek\psi\\
   &=(\rho\vek\psi)^{\textstyle\cdot}+\rho(\dvg\bv)\vek\psi
    =\big(\dot\rho+\rho\dvg\bv\big)\vek\psi+\rho\dot{\vek\psi}
    =\rho\dot{\vek\psi},
\end{aligned}
\]
the second equality being \eqr{2.20} applied componentwise to
$\rho\vek\psi$. So \eqr{6.10} simplifies to
\begin{equation}
   \rho\dot{\vek\psi}=\dvg\bA+\rho\br .
   \tag{6.23}\label{eq:6.23}
\end{equation}

Alternatively, these reduced equations may be derived directly by
localizing the global statements \eqr{6.1} and \eqr{6.2}, if we first
replace their left-hand sides by
\begin{equation}
\begin{aligned}
   \frac{\dl}{\dl t}\int_{P_{t}}\rho\varphi\,\dl v
   &=\frac{\dl}{\dl t}\int_{P}\rho_{\kappa}\varphi\,\dl V
    =\int_{P}\rho_{\kappa}\dot\varphi\,\dl V
    =\int_{P_{t}}\rho\dot\varphi\,\dl v
   \qquad\text{and}\\[1mm]
   \frac{\dl}{\dl t}\int_{P_{t}}\rho\vek\psi\,\dl v
   &=\frac{\dl}{\dl t}\int_{P}\rho_{\kappa}\vek\psi\,\dl V
    =\int_{P}\rho_{\kappa}\dot{\vek\psi}\,\dl V
    =\int_{P_{t}}\rho\dot{\vek\psi}\,\dl v ,
\end{aligned}
   \tag{6.24}\label{eq:6.24}
\end{equation}
respectively, where use has been made of \eqr{2.50}, $\dl v=J\,\dl V$,
of \eqr{2.59}, $\rho_{\kappa}=\rho J$, so that
$\rho\,\dl v=\rho_{\kappa}\,\dl V$, and of mass conservation in the
form \eqr{6.18}, which is what allows $\rho_{\kappa}$ to be treated as
a constant when the derivative is taken inside the integral over the
fixed region $P$.

\begin{remark}[the shortest route, and why both are
kept]\label{rem:ch6shortroute}
The route through \eqr{6.24} is three lines and the route through
\eqr{6.20} is eight. The longer one was taken first for two reasons.

The first is that \eqr{6.24} assumes smoothness throughout $P$: the
step $\frac{\dl}{\dl t}\int_{P}=\int_{P}\frac{\dl}{\dl t}$ is
Remark~\ref{rem:pullback} in action and has nothing to say when the
integrand jumps. Every jump condition in this chapter comes from the
long route, and none from the short one.

The second is that \eqr{6.20} shows \emph{where} mass conservation
enters: in exactly one place, as the coefficient of $\varphi$ in the
line
\[
   (\rho\varphi)'+\dvg(\rho\varphi\bv)
   =\underbrace{\big(\dot\rho+\rho\dvg\bv\big)}_{=\,0
     \text{ by }\eqr{6.17}}\varphi+\rho\dot\varphi .
\]
The reduced equation \eqr{6.21} is therefore not a new axiom but the
old one with mass conservation already spent. The compact form
\[
   (\rho\varphi)'+\dvg(\rho\varphi\bv)=\rho\dot\varphi
\]
is used at every balance law from here to the end of the book, each an
instance of $\rho\dot\varphi=\dvg\vek\pi+\rho\sigma$ with a particular
reading of the three letters.
\end{remark}

%=======================================================================
\section{Linear and rotational momenta; forces and
torques}\label{sec:ch6momenta}
%=======================================================================

The equations of motion for a continuum, due to Euler, are relations
connecting the linear and rotational momenta of a set $S\subset B$ of
material points to the forces and torques acting on it. This section
defines the four objects; the next asserts the two relations.

%-----------------------------------------------------------------------
\subsection{Linear and rotational momenta}\label{sec:ch6lrm}
%-----------------------------------------------------------------------

The linear momentum of an arbitrary sub-body is the vector
\begin{equation}
   \vek{l}(S,t)=\int_{P_{t}}\rho\bv\,\dl v
   =\int_{P}\rho_{\kappa}\bv\,\dl V ,
   \tag{6.25}\label{eq:6.25}
\end{equation}
and the rotational momentum relative to an origin located at position
$\bo$, also called the moment of momentum, is
\begin{equation}
   \bh_{\bo}(S,t)=\int_{P_{t}}\rho\,\bx\times\bv\,\dl v
   =\int_{P}\rho_{\kappa}\,\bx\times\bv\,\dl V ,
   \tag{6.26}\label{eq:6.26}
\end{equation}
where $\bx$ is the position vector of a material point $p\in S$
relative to $\bo$, in the sense of Proposition~\ref{prop:posvec} and
with the abuse of notation recorded in Remark~\ref{rem:abuse}.

The second equality in each of \eqr{6.25} and \eqr{6.26} is not a
change of notation. It is the substitution $\rho\,\dl v=
\rho_{\kappa}\,\dl V$, which combines the change of variables
\eqr{2.50}, $\dl v=J\,\dl V$, with the definition \eqr{2.59},
$\rho_{\kappa}=\rho J$. The integrands are unchanged because $\bv$ and
$\bx$ are the same objects read at the same material point; only the
measure and the domain have been transferred to $\kappa$. This is the
pull-back of Remark~\ref{rem:pullback} once more, and it is available
here precisely because mass conservation makes $\rho_{\kappa}$
independent of $t$.

\begin{remark}[rotational momentum depends on the origin, and that is
not a defect]\label{rem:ch6origin}
The subscript on $\bh_{\bo}$ is doing work. Replace $\bo$ by
$\bo'=\bo+\bc$ with $\bc$ fixed. Then the position of a material point
relative to the new origin is $\bx'=\bx-\bc$, and
\[
   \bh_{\bo'}=\int_{P_{t}}\rho\,(\bx-\bc)\times\bv\,\dl v
   =\bh_{\bo}-\bc\times\int_{P_{t}}\rho\bv\,\dl v
   =\bh_{\bo}-\bc\times\vek{l} .
\]
So rotational momentum is origin dependent unless the linear momentum
vanishes. The same computation applied to \eqr{6.31} and \eqr{6.32}
gives $\bm_{\bo'}=\bm_{\bo}-\bc\times\vek{f}$, so the torque shifts by
exactly the amount needed for the balance \eqr{6.35} to survive the
change, as it must: a physical law cannot depend on where one chooses
to put the origin. Section~\ref{sec:ch6com} makes the dependence
explicit by choosing the one origin that simplifies it, the center of
mass, and \eqr{6.47} is the result.
\end{remark}

%-----------------------------------------------------------------------
\subsection{Forces}\label{sec:ch6forces}
%-----------------------------------------------------------------------

We assume the forces acting on $S$ to be of two types. The first,
called \emph{body forces}, are expressible in the form
\begin{equation}
   \vek{f}_{b}(S,t)=\int_{P_{t}}\rho\bb\,\dl v
   =\int_{P}\rho_{\kappa}\bb\,\dl V ,
   \tag{6.27}\label{eq:6.27}
\end{equation}
where $\bb=\tilde\bb(\bx,t)=\hat\bb(\bX,t)$ is the body force per unit
mass, the carets and tildes carrying their usual meaning from
Remark~\ref{rem:decorations}. For example, near a point on the surface
of the Earth the body force due to gravity is $\bb=-g\bk$, where
$g\simeq9.81\ \mathrm{m/s^{2}}$ and $\bk$ is a fixed unit vector
directed vertically upward from the point in question.

The second type are the \emph{contact forces}, arising from contact of
$S$ with its surroundings. These act on the surface $\partial P_{t}$
and are expressed in the form
\begin{equation}
   \vek{f}_{c}(S,t)=\int_{\partial P_{t}}\bt\,\dl a ,
   \tag{6.28}\label{eq:6.28}
\end{equation}
where $\bt=\tilde\bt(\bx,t)$, the force per unit area, is the
\emph{Cauchy traction}, or simply the traction. It is often convenient
to express this same force as an integral over $\partial P$. Thus
\begin{equation}
   \vek{f}_{c}(S,t)=\int_{\partial P}\bp\,\dl A ,
   \tag{6.29}\label{eq:6.29}
\end{equation}
in which $\bp=\hat\bp(\bX,t)$, the \emph{Piola traction}, furnishes the
same force but measured per unit area of $\partial P$ rather than of
$\partial P_{t}$.

\begin{remark}[per unit mass, per unit area, and the fiction in
$\bp$]\label{rem:ch6perunit}
Three conventions are fixed in \eqr{6.27} to \eqr{6.29} and each is a
choice about what the quantity is proportional to.

\emph{Why $\bb$ is per unit mass.} Gravity pulls on matter, not on
volume, and the coefficient it pulls with is the mass. Writing the body
force density as $\rho\bb$ rather than as a single field therefore
isolates in $\bb$ the part that is a property of the external agency
and not of the body. In a uniform gravitational field $\bb$ is the same
vector for lead and for cork.

\emph{Why $\bt$ is per unit area.} A contact force is exerted by what
is on the other side of a surface, and how much of the other side is in
contact is measured by area. This is not a definition one is free to
make differently: \eqr{6.28} asserts that the contact force on a part
is an \emph{integral} over its boundary, that is, that it is absolutely
continuous with respect to area. Contact forces concentrated on curves
or points within the surface are thereby excluded.

\emph{What $\bp$ is, and what it is not.} While $\bp$ is defined on
$\partial P$, the actual contact forces act on $\partial P_{t}$.
Nothing is pushing on the reference configuration; $\kappa$ is a
labelling device, as Remark~\ref{rem:whyL} insisted. The Piola traction
is the same force with the bookkeeping changed, so that the total is
recovered by integrating over a surface that does not move. Equating
\eqr{6.28} and \eqr{6.29} for every part will yield $\bp\,\dl A
=\bt\,\dl a$ at \eqr{6.142}, and Piola--Nanson \eqr{2.54} will convert
that into the relation $\bP=\bT\bF^{*}$ between the two stresses.
Confusing $\bp$ with a traction that something exerts on $\kappa$ is
the standard way of getting the factors of $J$ wrong.
\end{remark}

The net force acting on $S$ is
\begin{equation}
   \vek{f}(S,t)=\vek{f}_{b}(S,t)+\vek{f}_{c}(S,t).
   \tag{6.30}\label{eq:6.30}
\end{equation}

%-----------------------------------------------------------------------
\subsection{Torques}\label{sec:ch6torques}
%-----------------------------------------------------------------------

The moment of the body forces, relative to the origin located at
position $\bo$, is
\begin{equation}
   \bm_{b}(S,t)=\int_{P_{t}}\rho\,\bx\times\bb\,\dl v
   =\int_{P}\rho_{\kappa}\,\bx\times\bb\,\dl V ,
   \tag{6.31}\label{eq:6.31}
\end{equation}
and that of the contact forces is
\begin{equation}
   \bm_{c}(S,t)=\int_{\partial P_{t}}\bx\times\bt\,\dl a
   =\int_{\partial P}\bx\times\bp\,\dl A .
   \tag{6.32}\label{eq:6.32}
\end{equation}

In this book we follow the development of classical continuum mechanics
as conceived by Cauchy, and thus assume that \emph{intrinsic} body and
contact torques are absent. Continua of this kind are called
\emph{non-polar} materials. The net torque acting on $S$, relative to
$\bo$, is then
\begin{equation}
   \bm_{\bo}(S,t)=\bm_{b}(S,t)+\bm_{c}(S,t).
   \tag{6.33}\label{eq:6.33}
\end{equation}

\begin{remark}[what ``non-polar'' assumes, and what it will
cost]\label{rem:ch6nonpolar}
Read \eqr{6.31} and \eqr{6.32} carefully. Every torque appearing there
is the moment of a \emph{force} about $\bo$: it is of the form
$\bx\times(\text{something})$, and it therefore vanishes at
$\bx=\bze$. Nothing in the definitions allows for a torque that acts at
a point without acting through a lever arm.

Such torques exist. Consider a material reinforced by a continuous
distribution of embedded fibres. If the fibres resist local bending and
torsion, then cutting the body across the fibres transmits not only a
force per unit area but also a moment per unit area, a \emph{couple
traction}, wherever the fibres cross the cut. A medium with couple
tractions is called \emph{polar}, and for it \eqr{6.32} must be
augmented by a term $\int_{\partial P_{t}}\bmm\,\dl a$ and \eqr{6.26}
by an intrinsic spin density, since the material points then carry
angular momentum of their own rather than merely orbiting the origin.

The assumption is not free. What it buys is the symmetry of the Cauchy
stress,
\eqr{6.117}: rotational momentum balance will turn out to add no
differential equation at all, only the algebraic statement
$\bT=\bT^{t}$, and that statement is exactly the statement that no
couple tractions are present. What it costs is that the theory has no
internal length scale coming from the stress. Whether a real material
is well described this way is decided by whether its response depends
on the size of the specimen at fixed shape, and for ordinary
engineering materials at ordinary scales it does not.
\end{remark}

%=======================================================================
\section{Euler's postulates}\label{sec:ch6euler}
%=======================================================================

Euler's postulates apply to an arbitrary part $S$ of a body. They
entail the fundamental assumption that the motion of the part is
governed by the same equations that govern the motion of a rigid body.
Thus the rate of change of the linear momentum of a fixed set of
material points is balanced by the net force acting on it,
\begin{equation}
   \frac{\dl}{\dl t}\vek{l}(S,t)=\vek{f}(S,t),
   \tag{6.34}\label{eq:6.34}
\end{equation}
and the rate of change of the rotational momentum is balanced by the
net torque,
\begin{equation}
   \frac{\dl}{\dl t}\bh_{\bo}(S,t)=\bm_{\bo}(S,t).
   \tag{6.35}\label{eq:6.35}
\end{equation}

A frame of reference in which these apply, consisting of point spaces
$\Ept$ and $\hEpt$, associated translation spaces $\Vt$ and $\hVt$
respectively, and an origin located at a fixed position $\bo$, is
called an \emph{inertial frame}. Euler's postulates are thus tantamount
to the assumption that an inertial frame exists. Chapter 7 will show
that there are infinitely many of them, and will obtain the equations
holding relative to arbitrary, generally non-inertial, frames.

\begin{remark}[which word in \eqr{6.34} is doing the
work]\label{rem:ch6everypart}
Neither \eqr{6.34} nor \eqr{6.35} is new as a statement about a whole
body: they are Newton's and Euler's laws for a rigid body, written
down. What is new, and what makes continuum mechanics possible, is the
phrase \emph{arbitrary part}.

A rigid body has six degrees of freedom, and \eqr{6.34} and \eqr{6.35}
are six equations. Applied once, to the body as a whole, they would
determine a rigid motion and nothing else; for a deformable body they
would be six equations for infinitely many unknowns and would determine
nothing. Applied to \emph{every} part, they become infinitely many
equations, one for each subregion, and that is exactly the hypothesis
the localization theorem needs. The step from \eqr{6.36} to \eqr{6.109}
is nothing but the conversion of ``for every part'' into ``at every
point''. Remark~\ref{rem:cancel} of Chapter 5 made the same observation
about the quantifier ``for all $\bc$''; here the quantifier ranges over
subregions, and Remark~\ref{rem:piolawhat} predicted that this would be
the pattern of the whole subject.

There is a genuine assumption hidden in the phrase as well. A part $S$
is not isolated: it is glued to the rest of the body. Postulating
\eqr{6.34} for $S$ requires that everything the rest of the body does
to it be representable as a body force plus a surface traction, with no
further term. Section~\ref{sec:ch6action} shows that this hypothesis is
strong enough to force action and reaction, which in particle mechanics
is a separate axiom.
\end{remark}

\begin{remark}[why ``postulates'', and why Euler's]%
\label{rem:ch6whypostulate}
The name records two decisions, one about logical status and one about
attribution. Both are worth unpacking, because a reader arrives here
with Newton in mind and expects \eqr{6.34} to be a law already proved.

\emph{Newton's laws are laws for a mass point.} A point has no
extension. The moment $\bx\times\vek{f}$ of a force acting on it is a
statement about the choice of origin, not about the particle, and it
carries nothing that $\vek{f}=m\ba$ does not already carry: cross
$m\ddot\bx=\vek{f}$ with $\bx$ on the left and \eqr{6.35} for a single
particle drops out, since $\bx\times m\ddot\bx
=(\bx\times m\dot\bx)^{\textstyle\cdot}$ by the cancellation displayed
after \eqr{6.37}. So nothing among Newton's three laws plays the part of
\eqr{6.35}.

\emph{For a finite system of particles the first balance is a theorem.}
Let $N$ particles carry masses $m_{a}$ and positions $\bx_{a}$, let
$\vek{f}^{e}_{a}$ be the force applied to the $a$th from outside the
system, let $\vek{f}_{ab}$ be the force applied to it by the $b$th, and
let each particle obey Newton's second law,
\[
   m_{a}\ddot\bx_{a}=\vek{f}^{e}_{a}+\sum_{b\neq a}\vek{f}_{ab} .
\]
Sum over $a$ and collect the interior forces in pairs:
\[
   \dot{\vek{l}}=\sum_{a}m_{a}\ddot\bx_{a}
   =\sum_{a}\vek{f}^{e}_{a}
   +\sum_{a<b}\big(\vek{f}_{ab}+\vek{f}_{ba}\big)
   =\sum_{a}\vek{f}^{e}_{a} ,
\]
the pair sum vanishing term by term because Newton's third law gives
$\vek{f}_{ba}=-\vek{f}_{ab}$. That is \eqr{6.34} for the system, proved
rather than assumed.

\emph{The second balance is not a theorem.} Repeat the computation with
the moment arm attached:
\[
   \dot{\bh}_{\bo}=\sum_{a}m_{a}\,\bx_{a}\times\ddot\bx_{a}
   =\sum_{a}\bx_{a}\times\vek{f}^{e}_{a}
   +\sum_{a<b}\big(\bx_{a}\times\vek{f}_{ab}
                  +\bx_{b}\times\vek{f}_{ba}\big),
\]
the left side again by $(\bx\times\dot\bx)^{\textstyle\cdot}
=\bx\times\ddot\bx$. The third law turns each pair term into
\[
   \bx_{a}\times\vek{f}_{ab}-\bx_{b}\times\vek{f}_{ab}
   =(\bx_{a}-\bx_{b})\times\vek{f}_{ab} ,
\]
which vanishes for every pair if and only if $\vek{f}_{ab}$ is parallel
to $\bx_{a}-\bx_{b}$ for every pair. Newton's third law as usually
stated says only that the pair forces are equal and opposite. That their
common line of action joins the two particles is a further assumption,
called the strong form of the law, and it is exactly what the moment
balance needs. Even in particle mechanics, then, \eqr{6.35} is not a
consequence of Newton's three laws.

\emph{For a continuum there are no pairs to appeal to.} The interior
action is not given as a sum over particle pairs but as the traction
field of \eqr{6.28}, and Sections~\ref{sec:ch6action}
and~\ref{sec:ch6noll} will show that the two properties one might have
hoped to use are themselves consequences of \eqr{6.34}: Cauchy's lemma
$\bt_{21}=-\bt_{12}$ at \eqr{6.83}, which is the continuum form of the
third law, and Cauchy's hypothesis \eqr{6.84} itself, by Noll's theorem.
The logical order is the reverse of the particle case. Nothing
independent of \eqr{6.34} is left from which \eqr{6.35} might be
deduced, and Proposition~\ref{prop:ch6indep} shows that the deduction is
in fact impossible.

\emph{Why the word ``postulate''.} Three things are assumed here and
none of them can be checked inside the theory. Both statements are
carried over from whole rigid bodies, where they are theorems of
particle mechanics with central interior forces, to \emph{every part} of
\emph{every} body, deformable or not; that extension is the subject of
Remark~\ref{rem:ch6everypart}. Both presuppose a frame and an origin in
which they hold, so together they assert that an inertial frame exists.
And the second asserts something the first does not. Truesdell, whose
account of this material shaped the modern presentation, called them
Euler's \emph{laws} of mechanics; the two names carry the same content,
and ``postulate'' merely keeps the assumed status in view.

\emph{Why Euler.} The \emph{Principia} of 1687 states three laws for
what Newton calls bodies, but the mathematics carried out with them is
the mathematics of mass points, and no moment of momentum principle
appears among them. Euler, in a memoir of 1750 whose title announces the
discovery of a new principle of mechanics, took $\vek{f}=m\ba$ and
applied it to every element of an extended body, which is the first
appearance of the quantifier that Remark~\ref{rem:ch6everypart}
identifies as the whole content of the phrase ``arbitrary part''. In
1775 he stated the moment of momentum principle as a second and
independent law, holding for every body and for every part of one. Both
\eqr{6.34} and \eqr{6.35}, in the form used here, are therefore his.

Newton's law did not disappear in the process. Equation \eqr{6.44} is
exactly $\vek{f}=m\ba$, and Section~\ref{sec:ch6com} obtains it as a
\emph{theorem}: it is the statement Euler's first postulate makes about
the center of mass. Newton's second law is thus a corollary of Euler's
first, recovered by applying the postulate to a whole part instead of to
every part of it. Naming the pair after Newton would hide both that
containment and the fact that the second of them is not his at all.
\end{remark}

Proposition~\ref{prop:ch6indep} makes the independence precise. It
anticipates two results of Section~\ref{sec:ch6local}: the local form
\eqr{6.109} of the linear momentum balance, and the symmetry
\eqr{6.117} of the Cauchy stress, which is everything the rotational
balance contributes at a point.

\begin{proposition}[the second postulate is independent of the
first]\label{prop:ch6indep}
The symmetry of the Cauchy stress is not a consequence of the local
linear momentum balance. There exist a motion and fields $\bT$ and $\bb$
satisfying
\[
   \rho\dot\bv=\rho\bb+\dvg\bT
\]
at every point of $\kappa_{t}$, with $\bT\neq\bT^{t}$ everywhere.
\end{proposition}

\begin{proof}
Take any motion whatever; it fixes $\rho$ and $\dot\bv$. For the stress
take the constant field
\[
   \bT=\be_{1}\otimes\be_{2} ,
\]
which is not symmetric, since $\bT^{t}=\be_{2}\otimes\be_{1}$ and the
two dyads differ, and which has $\dvg\bT=\bze$ because its Cartesian
components are constants. Now \emph{define} the body force by
\[
   \bb:=\dot\bv-\rho^{-1}\dvg\bT=\dot\bv .
\]
This is an admissible body force: \eqr{6.27} imposes nothing on $\bb$
beyond enough regularity for the integral to exist, and $\dot\bv$ is as
regular as the motion. With these choices the displayed balance holds
identically while $\bT$ is nowhere symmetric.
\end{proof}

The construction is deliberately trivial, and its triviality is the
point. Equation \eqr{6.109} is three scalar equations, and the stress it
constrains has nine components; the three components of $\Skw\bT$ never
enter it. Those three are precisely what the rotational balance
supplies, through \eqr{6.117}. Counting the same fact from the other
side: drop the non-polar assumption of Section~\ref{sec:ch6torques} and
admit couple tractions, and \eqr{6.109} survives unchanged while
\eqr{6.117} fails, as Remark~\ref{rem:ch6nonpolar} records. A statement
that can fail while \eqr{6.34} continues to hold is not a theorem of
\eqr{6.34}.

On invoking conservation of mass and making use of \eqr{6.24}, Euler's
postulates reduce to
\begin{equation}
   \int_{P_{t}}\rho\dot\bv\,\dl v=\int_{P_{t}}\rho\bb\,\dl v
   +\int_{\partial P_{t}}\bt\,\dl a
   \tag{6.36}\label{eq:6.36}
\end{equation}
and
\begin{equation}
   \int_{P_{t}}\rho\,\bx\times\dot\bv\,\dl v
   =\int_{P_{t}}\rho\,\bx\times\bb\,\dl v
   +\int_{\partial P_{t}}\bx\times\bt\,\dl a ,
   \tag{6.37}\label{eq:6.37}
\end{equation}
respectively. For \eqr{6.36} take $\vek\psi=\bv$ in \eqr{6.24}$_{2}$;
for \eqr{6.37} take $\vek\psi=\bx\times\bv$, and then use
\[
   (\bx\times\bv)^{\textstyle\cdot}
   =\dot\bx\times\bv+\bx\times\dot\bv
   =\bv\times\bv+\bx\times\dot\bv
   =\bx\times\dot\bv ,
\]
since $\dot\bx=\bv$ by Definition~\ref{def:va} and $\bv\times\bv=\bze$
because the cross product of a vector with itself vanishes, by the
antisymmetry recorded in Definition~\ref{def:cross}. The moment of
momentum therefore changes only through the acceleration, exactly as in
particle mechanics, and the same cancellation will be met again in
components at \eqr{6.116}.

Equivalently, the linear momentum balance may be expressed in the form
\begin{equation}
   \int_{P}\rho_{\kappa}\dot\bv\,\dl V=\int_{P}\rho_{\kappa}\bb\,\dl V
   +\int_{\partial P}\bp\,\dl A ,
   \tag{6.38}\label{eq:6.38}
\end{equation}
and the rotational momentum balance in the form
\begin{equation}
   \int_{P}\rho_{\kappa}\,\bx\times\dot\bv\,\dl V
   =\int_{P}\rho_{\kappa}\,\bx\times\bb\,\dl V
   +\int_{\partial P}\bx\times\bp\,\dl A ,
   \tag{6.39}\label{eq:6.39}
\end{equation}
with $\bx=\bchi_{\kappa}(\bX,t)$. These are \eqr{6.36} and \eqr{6.37}
with every integral pulled back to $\kappa$, using $\rho\,\dl v
=\rho_{\kappa}\,\dl V$ in the volume integrals and \eqr{6.29} in place
of \eqr{6.28} on the boundary.

%-----------------------------------------------------------------------
\subsection{Center of mass}\label{sec:ch6com}
%-----------------------------------------------------------------------

The \emph{center of mass} of $S$, occupying position $\bar{\bx}$
relative to $\bo$, is located at the mass-weighted average of position,
\begin{equation}
   m(S)\,\bar{\bx}(t)=\int_{P_{t}}\rho\bx\,\dl v
   =\int_{P}\rho_{\kappa}\bx\,\dl V ,
   \tag{6.40}\label{eq:6.40}
\end{equation}
where $m(S)$ is the fixed mass of $S$; that $m(S)$ is fixed is
\eqr{6.13}. The position of a material point $p\in S$ relative to the
center of mass is
\begin{equation}
   \vek\pi=\bx-\bar{\bx} .
   \tag{6.41}\label{eq:6.41}
\end{equation}
The letter clashes with the flux vector of \eqr{6.1}; the two never
appear in the same computation, and from here to the end of
Section~\ref{sec:ch6rigid} $\vek\pi$ means relative position.

Accordingly the mass-weighted average of the relative position vanishes
identically:
\begin{equation}
\begin{aligned}
   \int_{P_{t}}\rho\vek\pi\,\dl v
   &=\int_{P_{t}}\rho(\bx-\bar{\bx})\,\dl v\\
   &=\Big[m(S)-\int_{P_{t}}\rho\,\dl v\Big]\bar{\bx}=\bze ,
\end{aligned}
   \tag{6.42}\label{eq:6.42}
\end{equation}
the first term having been evaluated by \eqr{6.40} and the second by
taking the constant vector $\bar{\bx}$ outside the integral, after
which $\int_{P_{t}}\rho\,\dl v=m(S)$ makes the bracket vanish. The
statement is a tautology: the average of the deviation from the average
is zero. Its usefulness is that it holds at every instant, so it may be
differentiated.

Differentiating, and invoking conservation of mass in the form
\eqr{6.24}$_{2}$ with $\vek\psi=\vek\pi$,
\begin{equation}
\begin{aligned}
   \bze=\frac{\dl}{\dl t}\int_{P_{t}}\rho\vek\pi\,\dl v
   &=\int_{P_{t}}\rho\dot{\vek\pi}\,\dl v\\
   &=\int_{P_{t}}\rho\bv\,\dl v-m(S)\,\dot{\bar{\bx}} ,
\end{aligned}
   \tag{6.43}\label{eq:6.43}
\end{equation}
where the last line used $\dot{\vek\pi}=\dot\bx-\dot{\bar{\bx}}
=\bv-\dot{\bar{\bx}}$, legitimate because $\bar{\bx}$ is a function of
$t$ alone and is therefore the same vector at every material point, so
that $\int_{P_{t}}\rho\dot{\bar{\bx}}\,\dl v=\dot{\bar{\bx}}
\int_{P_{t}}\rho\,\dl v=m(S)\dot{\bar{\bx}}$.

Comparing with \eqr{6.25}, this says $\vek{l}(S,t)=m(S)\dot{\bar{\bx}}$:
the linear momentum of $S$ is identical to that of a single particle of
the same mass located at the center of mass. Euler's postulate
\eqr{6.34} is thus equivalent, granted conservation of mass, to
\begin{equation}
   m(S)\,\ddot{\bar{\bx}}=\vek{f}(S,t),
   \tag{6.44}\label{eq:6.44}
\end{equation}
where $\ddot{\bar{\bx}}$ is the acceleration of the center of mass.

Using \eqr{6.41}, we also have
\begin{equation}
\begin{aligned}
   \bh_{\bo}(S,t)
   &=\int_{P_{t}}\rho\,\bx\times\dot\bx\,\dl v
    =\int_{P_{t}}\rho\,(\bar{\bx}+\vek\pi)\times
      \big[\dot{\bar{\bx}}+\dot{\vek\pi}\big]\,\dl v\\
   &=m(S)\,\bar{\bx}\times\dot{\bar{\bx}}
    +\bar{\bx}\times\int_{P_{t}}\rho\dot{\vek\pi}\,\dl v
    +\Big[\int_{P_{t}}\rho\vek\pi\,\dl v\Big]\times\dot{\bar{\bx}}
    +\bh(S,t),
\end{aligned}
   \tag{6.45}\label{eq:6.45}
\end{equation}
where
\begin{equation}
   \bh(S,t)=\int_{P_{t}}\rho\,\vek\pi\times\dot{\vek\pi}\,\dl v
   \tag{6.46}\label{eq:6.46}
\end{equation}
is the rotational momentum relative to the center of mass. The
expansion in \eqr{6.45} used bilinearity of the cross product to
produce four terms, and then moved the factors that do not depend on
the material point outside their integrals: $\bar{\bx}$ and
$\dot{\bar{\bx}}$ are functions of $t$ only, so
$\int\rho\,\bar{\bx}\times\dot{\bar{\bx}}\,\dl v
=m(S)\bar{\bx}\times\dot{\bar{\bx}}$, and in the second and third terms
one of the two slots is likewise constant over $P_{t}$.

Noting that the two integrals appearing in \eqr{6.45} vanish, the
second by \eqr{6.42} and the first by \eqr{6.43}$_{2}$, and using
$\vek{l}=m(S)\dot{\bar{\bx}}$, we conclude that
\begin{equation}
   \bh_{\bo}(S,t)=\bar{\bx}(t)\times\vek{l}(S,t)+\bh(S,t).
   \tag{6.47}\label{eq:6.47}
\end{equation}

We also have
\begin{equation}
\begin{aligned}
   \bm_{\bo}(S,t)
   &=\int_{P_{t}}\rho\,(\bar{\bx}+\vek\pi)\times\bb\,\dl v
    +\int_{\partial P_{t}}(\bar{\bx}+\vek\pi)\times\bt\,\dl a\\
   &=\bar{\bx}(t)\times\vek{f}(S,t)+\bm(S,t),
\end{aligned}
   \tag{6.48}\label{eq:6.48}
\end{equation}
where
\begin{equation}
   \bm(S,t)=\int_{P_{t}}\rho\,\vek\pi\times\bb\,\dl v
   +\int_{\partial P_{t}}\vek\pi\times\bt\,\dl a
   \tag{6.49}\label{eq:6.49}
\end{equation}
is the net torque about the center of mass. Here the terms carrying
$\bar{\bx}$ were collected as
$\bar{\bx}\times\big[\int\rho\bb\,\dl v+\int\bt\,\dl a\big]
=\bar{\bx}\times\vek{f}$ by \eqr{6.27}, \eqr{6.28} and \eqr{6.30}.

\begin{remark}[a misprint in the printed \eqr{6.48} and
\eqr{6.49}]\label{rem:ch6misprint48}
The printed text carries a factor $\rho$ in the surface integrals of
both displays, writing $\int_{\partial P_{t}}\rho(\bar{\bx}+\vek\pi)
\times\bt\,\dl a$ and $\int_{\partial P_{t}}\rho\vek\pi\times\bt\,\dl
a$. That factor should not be there. The contact torque is defined in
\eqr{6.32} as $\int_{\partial P_{t}}\bx\times\bt\,\dl a$ with no
density: $\bt$ is already a force per unit area, whereas $\bb$ is a
force per unit \emph{mass} and needs the $\rho$ to convert it. With the
stray $\rho$ the last equality of \eqr{6.48} would fail, since the
bracket would no longer be $\vek{f}$, and \eqr{6.51} would not follow.
The displays above are corrected.
\end{remark}

We thus recast Euler's postulate \eqr{6.35}. Differentiate \eqr{6.47}
and set the result equal to \eqr{6.48}:
\begin{equation}
\begin{aligned}
   \bar{\bx}(t)\times\vek{f}(S,t)+\bm(S,t)
   =\frac{\dl}{\dl t}\bh(S,t)
    +\dot{\bar{\bx}}(t)\times\vek{l}(S,t)
    +\bar{\bx}(t)\times\frac{\dl}{\dl t}\vek{l}(S,t),
\end{aligned}
   \tag{6.50}\label{eq:6.50}
\end{equation}
the product rule for the cross product being legitimate because the
cross product is bilinear. Now dispose of two of the three terms on the
right. By \eqr{6.34}, $\frac{\dl}{\dl t}\vek{l}=\vek{f}$, so the last
term is $\bar{\bx}\times\vek{f}$ and cancels the first term on the
left. By \eqr{6.43} in the form $\vek{l}=m(S)\dot{\bar{\bx}}$,
\[
   \dot{\bar{\bx}}\times\vek{l}
   =\dot{\bar{\bx}}\times m(S)\dot{\bar{\bx}}
   =m(S)\big(\dot{\bar{\bx}}\times\dot{\bar{\bx}}\big)=\bze ,
\]
a vector crossed with itself. What survives is
\begin{equation}
   \bm(S,t)=\frac{\dl}{\dl t}\bh(S,t).
   \tag{6.51}\label{eq:6.51}
\end{equation}

Thus Euler's postulate also holds relative to the center of mass, its
motion in the inertial frame notwithstanding. This is a genuinely
useful fact and not a triviality: the center of mass is in general
accelerating, so the frame that follows it is not inertial, and one
would expect the balance law to acquire correction terms. It does not,
and \eqr{6.50} shows why: the two terms that could have produced
corrections are the two that vanished, one because $\vek{l}$ is
parallel to $\dot{\bar{\bx}}$ and the other because the force balance
\eqr{6.34} already holds.

%-----------------------------------------------------------------------
\subsection{Rigid bodies}\label{sec:ch6rigid}
%-----------------------------------------------------------------------

The foregoing development naturally subsumes the motions of rigid
bodies. These are necessarily of the form \eqr{3.121}, that is,
\begin{equation}
   \bx=\bQ(t)\bX+\bc(t),
   \tag{6.52}\label{eq:6.52}
\end{equation}
where $\bQ(t)$ is a spatially uniform two-point rotation tensor and
$\bc(t)$ is a vector-valued function. Accordingly
\begin{equation}
   \vek\pi=\bQ(t)\vek\Pi,\qquad\text{where}\quad
   \vek\Pi=\bX-\bar{\bX}
   \tag{6.53}\label{eq:6.53}
\end{equation}
is the referential position of $p\in S$ relative to the center of mass
of $P$, located at position $\bar{\bX}$. That $\bar{\bX}$ really is the
center of mass of $P\subset\kappa$, and hence that \eqr{6.53} is
consistent with \eqr{6.41}, is Problem 2. Granting it, \eqr{6.53}
follows from \eqr{6.52} by subtraction: $\bar{\bx}=\bQ\bar{\bX}+\bc$,
so $\vek\pi=\bx-\bar{\bx}=\bQ(\bX-\bar{\bX})=\bQ\vek\Pi$, the constant
$\bc$ cancelling.

Thus
\begin{equation}
   \dot{\vek\pi}=\dot\bQ\vek\Pi=\bOm(t)\vek\pi=\bom(t)\times\vek\pi ,
   \tag{6.54}\label{eq:6.54}
\end{equation}
where $\bom$ is the axial vector, in the sense of
Proposition~\ref{prop:axial}, of the skew tensor
\begin{equation}
   \bOm=\dot\bQ\bQ^{t},
   \tag{6.55}\label{eq:6.55}
\end{equation}
this being simply a special case of \eqr{4.33}. Two steps in \eqr{6.54}
need saying. First, $\vek\Pi$ does not depend on $t$: both $\bX$ and
$\bar{\bX}$ are referential positions, and $\bar{\bX}$ is fixed because
$\rho_{\kappa}$ is a function of $\bX$ alone by \eqr{6.18}. So
differentiating \eqr{6.53} differentiates only $\bQ$. Second,
$\dot\bQ\vek\Pi=\dot\bQ(\bQ^{t}\bQ)\vek\Pi=(\dot\bQ\bQ^{t})(\bQ
\vek\Pi)=\bOm\vek\pi$, using $\bQ^{t}\bQ=\hI$ and then \eqr{6.53}. That
$\bOm$ is skew follows by differentiating $\bQ\bQ^{t}=\I$:
\[
   \dot\bQ\bQ^{t}+\bQ\dot\bQ^{t}=\Ob ,\qquad\text{that is}\qquad
   \bOm+\bOm^{t}=\Ob .
\]

Consider a material curve with unit tangent $\bM$ at $\bX$. From
\eqr{2.67} and \eqr{3.120} we have that the unit tangent $\bm$ to the
image of this curve after deformation, and the associated stretch
$\mu$, satisfy
\begin{equation}
   \mu\bm=\bF\bM=\bQ\bM;\qquad\text{hence}\qquad
   \mu=1\quad\text{and}\quad\bm=\bQ\bM ,
   \tag{6.56}\label{eq:6.56}
\end{equation}
because $\bF=\bQ$ for the motion \eqr{6.52} and because an orthogonal
tensor preserves length, $|\bQ\bM|^{2}=\ip{\bQ\bM}{\bQ\bM}
=\ip\bM{\bQ^{t}\bQ\bM}=\ip\bM{\hI\bM}=1$, by
Proposition~\ref{prop:orth}. Differentiating $\bm=\bQ\bM$, and using
$\dot\bM=\bze$ because $\bM$ is a material vector in the sense of
Definition~\ref{def:matvec} and Remark~\ref{rem:matvec},
\begin{equation}
   \dot\bm=\bOm\bm=\bom\times\bm ,
   \tag{6.57}\label{eq:6.57}
\end{equation}
by the same insertion of $\bQ^{t}\bQ=\hI$ that produced \eqr{6.54}.

Let $\{\hbE_{i}\}$ be a fixed orthonormal basis for $\hVt$ and let
\begin{equation}
   \be^{*}_{i}(t)=\bQ(t)\hbE_{i}\in\Vt .
   \tag{6.58}\label{eq:6.58}
\end{equation}
Then $\{\be^{*}_{i}\}$ is fixed in the body and rotates with it: it is
orthonormal at every instant, again by
Proposition~\ref{prop:orth}, and it is the image under the motion of a
triad of material fibres. On identifying $\bm$ in \eqr{6.57} with an
element of this basis we have simply
\begin{equation}
   \dot{\be}^{*}_{i}=\bom\times\be^{*}_{i} .
   \tag{6.59}\label{eq:6.59}
\end{equation}

\subsection*{The same kinematics, read from Chapter 4}

Everything from \eqr{6.52} to \eqr{6.59} is Chapter 4 specialised to a
rigid motion, and reading it that way does three things. It identifies
the angular velocity of rigid-body dynamics with an object already
constructed there. It supplies proofs of $\mu=1$ in \eqr{6.56} and of
$\rho=\rho_{\kappa}$ that never mention $\det\bQ$. And it locates
\eqr{6.57} as the one case in which a warning issued in Chapter 4 does
not apply.

\emph{The angular velocity tensor is the spin tensor.} For the motion
\eqr{6.52} the deformation gradient is $\bF=\bQ(t)$, spatially uniform
because $\bQ$ is. So by \eqr{4.3} the spatial velocity gradient is
\[
   \bL=\dot\bF\bF^{-1}=\dot\bQ\bQ^{-1}=\dot\bQ\bQ^{t} ,
\]
the last step by $\bQ^{-1}=\bQ^{t}$. The right-hand side is \eqr{6.55}.
Hence
\[
   \bOm=\bL ,
\]
which is \eqr{4.29}, obtained in Section~\ref{sec:rigidrate} as the
velocity gradient of a rigid motion. Chapter 4 recorded there that it is
skew, by Problem 1(c) of that chapter, which is the computation repeated
below \eqr{6.55}. Splitting $\bL$ by \eqr{4.7} and invoking
Remark~\ref{rem:Drigid}, which states that $\bD=\Ob$ characterises
rigidity,
\[
   \bD=\Ob ,\qquad \bW=\bL=\bOm .
\]
The angular velocity tensor of \eqr{6.55} and the spin tensor of
\eqr{4.7} are one tensor under two names. Their axial vectors agree as
well, so $\bom$ is the vorticity $\bw$ of \eqr{4.22}, and \eqr{4.23}
then gives a formula for the angular velocity that
Section~\ref{sec:ch6rigid} does not otherwise produce:
\[
   \bom=\tfrac12\crl\bv .
\]
Section~\ref{sec:disc} carried out exactly this computation for a
spinning disc and obtained $\bom=\omega\bn$ at \eqr{4.26}.

\emph{\eqr{6.54} is \eqr{4.32}.} Chapter 4 integrated the velocity field
of a rigid motion at \eqr{4.30} and reached \eqr{4.32},
\[
   \bv(\by,t)-\bv(\bx,t)=\bw(t)\times(\by-\bx) .
\]
Take $\by$ to be the position of the material point $p\in S$ and $\bx$
the position $\bar{\bx}$ of the center of mass. The left side is
$\dot\by-\dot{\bar{\bx}}=\dot{\vek\pi}$ by \eqr{6.41}, and the right
side is $\bom\times\vek\pi$. That is \eqr{6.54}. The book cites
\eqr{4.33} here; \eqr{4.32} is the same statement with the arbitrary
spatially uniform vector $\bd$ already removed by subtracting the
velocity at one point from the velocity at another, and subtraction is
what \eqr{6.54} needs, since $\vek\pi$ is a difference of positions.

\emph{$\mu=1$, by a rate argument.} The derivation of \eqr{6.56} above
used $|\bQ\bM|=|\bM|$. Chapter 4 reaches the same conclusion without
naming $\bQ$. Equation \eqr{4.15} reads
\[
   (\ln\mu)^{\textstyle\cdot}=\ip\bm{\bD\bm} ,
\]
and $\bD=\Ob$, so $(\ln\mu)^{\textstyle\cdot}=0$ and $\mu$ does not
change in time. At the instant at which the body occupies $\kappa$ we
have $\bF=\shift$ and therefore $\mu=1$, so $\mu\equiv1$ for all time.
The rate argument says more than the algebraic one: it says that no
material curve in a rigid motion is stretching at any instant, which is
the differential form of the distance-preserving property \eqr{3.113}.

\emph{$\rho=\rho_{\kappa}$, by a rate argument.} Equation \eqr{6.66}
needed $J=1$ and took it from $\det\bQ=1$. Chapter 4 gets it from the
skewness of $\bOm$ alone. Problem 2 of Chapter 4 established
$\dot J=J\tr\bL$, so here
\[
   \dot J=J\tr\bOm=0 ,
\]
because the trace of a skew tensor vanishes: $\tr\bOm=\tr\bOm^{t}
=\tr(-\bOm)=-\tr\bOm$, so $\tr\bOm=0$. Hence $J$ is constant in time,
and it equals $1$ at the instant at which $\kappa$ is occupied. By
\eqr{2.59} this is $\rho=\rho_{\kappa}$, and by \eqr{6.17} it is also
$\dot\rho=0$: a rigid motion is isochoric, and carries its density along
unchanged.

\emph{\eqr{6.57} is \eqr{4.16} with $\bD=\Ob$.} The general law for the
rate of a material direction is \eqr{4.16},
\[
   \dot\bm=\bD\bm-(\ip\bm{\bD\bm})\bm+\bW\bm .
\]
Setting $\bD=\Ob$ removes the first two terms and leaves
$\dot\bm=\bW\bm=\bom\times\bm$, which is \eqr{6.57}. The contrast with
Chapter 4 is sharp. Equation \eqr{4.18} asserted $\dot\bm=\bW\bm$ only at an instant
$t_{*}$, and only for the special fibres momentarily aligned with a
principal direction of $\bD$; Remark~\ref{rem:Wtrap} warned that for a
general fibre in a general motion the identity is false, so that the
spin tensor is not the spin of the material. In a rigid motion it is,
for every fibre and at every instant, and \eqr{6.58} to \eqr{6.59} hang
an entire orthonormal frame on that fact. The frame $\{\be^{*}_{i}\}$
exists because rigidity is exactly the hypothesis under which all three
of its members obey the same rotation law.

Rigidity is not quite the right condition, and the sharp statement is
the following.

\begin{proposition}[when the spin tensor is the spin of the
material]\label{prop:ch6spinfibre}
At a point and an instant, $\dot\bm=\bW\bm$ for \emph{every} material
direction $\bm$ if and only if $\bD$ is spherical, that is
$\bD=\alpha\I$ for some scalar $\alpha$. The motion is rigid there if
and only if in addition $\alpha=0$.
\end{proposition}

\begin{proof}
By \eqr{4.16}, $\dot\bm-\bW\bm=\bD\bm-(\ip\bm{\bD\bm})\bm$ for every
unit $\bm$, so the stated condition is
\[
   \bD\bm=(\ip\bm{\bD\bm})\bm\qquad\text{for every unit }\bm ,
\]
that is, every direction is a principal direction of $\bD$.

Suppose this holds and let $\bD=\sum_{i}d_{i}\bv_{i}\otimes\bv_{i}$ be
the spectral representation supplied by Theorem~\ref{thm:spectral}, the
$\bv_{i}$ orthonormal. If two principal values differed, say
$d_{1}\neq d_{2}$, take $\bm=(\bv_{1}+\bv_{2})/\sqrt2$, a unit vector.
Then $\bD\bm=(d_{1}\bv_{1}+d_{2}\bv_{2})/\sqrt2$ and
$\ip\bm{\bD\bm}=\tfrac12(d_{1}+d_{2})$, so
\[
\begin{aligned}
   \bD\bm-(\ip\bm{\bD\bm})\bm
   &=\frac{1}{\sqrt2}\Big[\Big(d_{1}-\frac{d_{1}+d_{2}}{2}\Big)\bv_{1}
   +\Big(d_{2}-\frac{d_{1}+d_{2}}{2}\Big)\bv_{2}\Big]\\
   &=\frac{d_{1}-d_{2}}{2\sqrt2}\,(\bv_{1}-\bv_{2}),
\end{aligned}
\]
which is not $\bze$, because $d_{1}\neq d_{2}$ and $\bv_{1}\neq\bv_{2}$
are independent. This contradicts the condition, so all three principal
values coincide and $\bD=\alpha\I$ with $\alpha$ their common value.

Conversely, if $\bD=\alpha\I$ then $\bD\bm=\alpha\bm$ and
$\ip\bm{\bD\bm}=\alpha$ for every unit $\bm$, so
$\bD\bm-(\ip\bm{\bD\bm})\bm=\bze$ identically.

For the last statement, $\bD=\Ob$ characterises rigidity by
Remark~\ref{rem:Drigid}, and $\alpha\I=\Ob$ if and only if $\alpha=0$.
\end{proof}

The intermediate case $\alpha\neq0$ is not empty. By Problem 2 of
Chapter 4, $\bD=\alpha\I$ gives $\dot J/J=\tr\bD=3\alpha$, so the motion
is instantaneously a rigid rotation together with a uniform dilatation
at rate $3\alpha$. Every fibre turns with $\bW$ and every fibre stretches
at the same logarithmic rate $\alpha$, by \eqr{4.15}. That is the widest
class for which the sentence ``$\bW$ is the spin of the material'' is
true without qualification, and rigid motions are the members of it that
also preserve length.

\emph{One consequence carried forward.} Since $\bL=\bW$ in a rigid
motion and $\bT$ is symmetric by \eqr{6.117}, the stress power density
of \eqr{6.153} is
\[
   \bT\cdot\bL=\bT\cdot\bW=0
\]
by Lemma~\ref{lem:symskw}, a symmetric tensor contracted with a skew one
vanishing. Section~\ref{sec:ch6energy} uses precisely this to decouple
the thermal and the mechanical problem in a rigid body.

Let $\bu\in\Vt$ be an arbitrary vector, possibly time dependent. We
decompose it in the rotating basis, $\bu=u_{i}\be^{*}_{i}$ with
$u_{i}=\ip\bu{\be^{*}_{i}}$. Thus
\begin{equation}
\begin{aligned}
   \dot\bu&=\dot u_{i}\be^{*}_{i}+u_{i}\dot{\be}^{*}_{i}\\
          &=\mathring{\bu}+\bom\times\bu ,
\end{aligned}
   \tag{6.60}\label{eq:6.60}
\end{equation}
where
\begin{equation}
   \mathring{\bu}=\dot u_{i}\be^{*}_{i} ,
   \tag{6.61}\label{eq:6.61}
\end{equation}
the \emph{co-rotational derivative}, is the derivative relative to the
rotating basis. The second line of \eqr{6.60} used \eqr{6.59} and then
the linearity of the cross product: $u_{i}(\bom\times\be^{*}_{i})
=\bom\times(u_{i}\be^{*}_{i})=\bom\times\bu$. Thus $\mathring{\bu}$
vanishes if and only if $\bu$ is fixed in the rotating body, since it
vanishes exactly when all the components of $\bu$ on a basis that
rotates with the body are constant. From \eqr{6.54} it follows that
$\vek\pi$ is such a vector.

\begin{remark}[two rates of change, and why both are
needed]\label{rem:ch6corot}
Equation \eqr{6.60} splits $\dot\bu$ into the part due to the vector
changing and the part due to the frame turning. Neither piece is
frame-free by itself and the sum is: $\dot\bu$ is the honest material
derivative, computed with respect to the fixed basis of the inertial
frame.

Whether one wants $\dot\bu$ or $\mathring{\bu}$ depends on the
question. A vector painted on the body has $\mathring{\bu}=\bze$ and
$\dot\bu=\bom\times\bu\neq\bze$: it is not constant, it is merely not
changing \emph{relative to the body}. The inertia tensor of
\eqr{6.65} is an object of exactly this kind, and \eqr{6.68} is the
tensor version of the same split. The same construction reappears in
Problem 12(d), where the co-rotational rate of the left stretch tensor
turns out to be the deformation measure conjugate to the Bell stress.
\end{remark}

From \eqr{6.46} and \eqr{6.54} we have that
\begin{equation}
   \bh(S,t)=\int_{P_{t}}\rho\,\vek\pi\times(\bom\times\vek\pi)\,\dl v
   \tag{6.62}\label{eq:6.62}
\end{equation}
for rigid bodies. We would like to exploit the spatial uniformity of
the spin $\bom$ to factor it out of the integral, and we cannot do so
directly because $\bom$ sits inside a cross product with the
integration variable. The obstacle is removed by the result of Problem
11 of Chapter 1, $\ba\times(\bb\times\bc)=(\ip\ba\bc)\bb-(\ip\ba\bb)
\bc$, applied with $\ba=\bc=\vek\pi$ and $\bb=\bom$:
\begin{equation}
\begin{aligned}
   \vek\pi\times(\bom\times\vek\pi)
   &=(\ip{\vek\pi}{\vek\pi})\bom-(\ip{\vek\pi}\bom)\vek\pi\\
   &=\big[(\ip{\vek\pi}{\vek\pi})\I-\vek\pi\otimes\vek\pi\big]\bom ,
\end{aligned}
   \tag{6.63}\label{eq:6.63}
\end{equation}
the second line by $\I\bom=\bom$ and
$(\ip{\vek\pi}\bom)\vek\pi=(\vek\pi\otimes\vek\pi)\bom$, which is the
definition of the tensor product, Definition~\ref{def:tp}. The point of
the manoeuvre is that $\bom$ now stands to the right of a tensor that
does not involve it, so it can be taken outside the integral. This
yields
\begin{equation}
   \bh(S,t)=\vek{J}(S,t)\,\bom(t),
   \tag{6.64}\label{eq:6.64}
\end{equation}
where
\begin{equation}
   \vek{J}(S,t)=\int_{P_{t}}\rho
   \big[(\ip{\vek\pi}{\vek\pi})\I-\vek\pi\otimes\vek\pi\big]\,\dl v
   \tag{6.65}\label{eq:6.65}
\end{equation}
is the symmetric \emph{inertia tensor} of the rigid body. It is
symmetric because $\I$ and $\vek\pi\otimes\vek\pi$ both are, the latter
by $(\ba\otimes\ba)^{t}=\ba\otimes\ba$. Problem 4 shows that it is in
fact positive definite.

\begin{remark}[a collision of letters]\label{rem:ch6Jclash}
The book writes $J$ for the inertia tensor and $J$ for
$\det\bF$ in the same paragraph. Here the inertia tensor is
$\vek{J}$, carrying the under-tilde that marks every tensor in these
notes, and $J=\det\bF$ keeps its meaning from \eqr{2.42}. The next
display uses both.
\end{remark}

Equation \eqr{6.65} may be simplified on noting that $J=\det\bQ=1$ for
a rigid motion, and hence, by \eqr{2.59}, that $\rho=\rho_{\kappa}$ and
$\rho\,\dl v=\rho_{\kappa}\,\dl V$. Further, \eqr{6.53} implies
\[
   \ip{\vek\pi}{\vek\pi}=\ip{\bQ\vek\Pi}{\bQ\vek\Pi}
   =\ip{\vek\Pi}{\bQ^{t}\bQ\vek\Pi}=\ip{\vek\Pi}{\vek\Pi},
\]
so that
\begin{equation}
\begin{aligned}
   \vek{J}(S,t)
   &=\int_{P}\rho_{\kappa}\big[(\ip{\vek\Pi}{\vek\Pi})\bQ\bQ^{t}
     -\bQ\vek\Pi\otimes\bQ\vek\Pi\big]\,\dl V\\
   &=\bQ(t)\,\vek{J}_{\kappa}(S)\,\bQ(t)^{t},
\end{aligned}
   \tag{6.66}\label{eq:6.66}
\end{equation}
where $\vek{J}_{\kappa}(S)$ is the fixed referential tensor defined by
\begin{equation}
   \vek{J}_{\kappa}(S)=\int_{P}\rho_{\kappa}
   \big[(\ip{\vek\Pi}{\vek\Pi})\hI-\vek\Pi\otimes\vek\Pi\big]\,\dl V .
   \tag{6.67}\label{eq:6.67}
\end{equation}
The first line of \eqr{6.66} replaced the spatial identity by
$\I=\bQ\bQ^{t}$, which is orthogonality read in the direction
$\Vt\to\Vt$, and replaced $\vek\pi$ by $\bQ\vek\Pi$. The second line
used the dyad rule $\bQ(\vek\Pi\otimes\vek\Pi)\bQ^{t}
=\bQ\vek\Pi\otimes\bQ\vek\Pi$ of Lemma~\ref{lem:dyadprod} backwards,
and then took the constant tensors $\bQ$ and $\bQ^{t}$ outside the
integral, which is legitimate because they do not depend on $\bX$.

What \eqr{6.66} says is that the inertia tensor is fixed in the body:
its referential version $\vek{J}_{\kappa}$ does not depend on time at
all, and all the time dependence of $\vek{J}$ is the rotation carrying
it along. Stated as a rate, this is the vanishing of the co-rotational
derivative $\mathring{\vek{J}}$, defined by
\begin{equation}
   \mathring{\vek{J}}=\dot{\vek{J}}-\bOm\vek{J}+\vek{J}\bOm .
   \tag{6.68}\label{eq:6.68}
\end{equation}
This follows easily from \eqr{6.55} and \eqr{6.66}. Differentiating
\eqr{6.66}$_{2}$ and using $\dot\bQ=\bOm\bQ$ from \eqr{6.55},
together with $\dot{\bQ^{t}}=(\dot\bQ)^{t}=(\bOm\bQ)^{t}
=\bQ^{t}\bOm^{t}=-\bQ^{t}\bOm$,
\[
   \dot{\vek{J}}=\dot\bQ\vek{J}_{\kappa}\bQ^{t}
   +\bQ\vek{J}_{\kappa}\dot{\bQ^{t}}
   =\bOm\bQ\vek{J}_{\kappa}\bQ^{t}-\bQ\vek{J}_{\kappa}\bQ^{t}\bOm
   =\bOm\vek{J}-\vek{J}\bOm ,
\]
whence $\mathring{\vek{J}}=\Ob$. Problem 3 shows that \eqr{6.68} is the
tensor counterpart of \eqr{6.61}: $\mathring{\vek{J}}$ is the tensor
whose components on the rotating basis are the derivatives of the
components of $\vek{J}$.

The form of \eqr{6.68} is not a convention. Once the vector rate
\eqr{6.61} is fixed, the tensor rate is determined by the demand that it
annihilate tensors built from vectors fixed in the body, and both signs
follow from that demand.

\begin{proposition}[the tensor rate is forced by the vector rate]%
\label{prop:ch6jaumann}
Write $\mathring\bu=\dot\bu-\bOm\bu$, which is \eqr{6.60} rearranged. If
$\ba$ and $\bb$ satisfy $\mathring\ba=\mathring\bb=\bze$ and
$\bA=\ba\otimes\bb$, then
\[
   \dot\bA=\bOm\bA-\bA\bOm ,
\]
so that $\mathring\bA=\Ob$ with $\mathring\bA$ defined as in
\eqr{6.68}.
\end{proposition}

\begin{proof}
$\mathring\ba=\bze$ means $\dot\ba=\bOm\ba$, and likewise
$\dot\bb=\bOm\bb$. The product rule applied to the tensor product gives
\[
   \dot\bA=\dot\ba\otimes\bb+\ba\otimes\dot\bb
   =(\bOm\ba)\otimes\bb+\ba\otimes(\bOm\bb) .
\]
Rewrite each term by testing it on an arbitrary $\bc\in\Vt$. For the
first,
\[
   \big[(\bOm\ba)\otimes\bb\big]\bc=(\ip\bb\bc)\,\bOm\ba
   =\bOm\big[(\ip\bb\bc)\ba\big]=\bOm(\ba\otimes\bb)\bc=\bOm\bA\bc ,
\]
where the second equality is linearity of $\bOm$ and the third is
Definition~\ref{def:tp}. Hence $(\bOm\ba)\otimes\bb=\bOm\bA$. For the
second,
\[
   \big[\ba\otimes(\bOm\bb)\big]\bc=(\ip{\bOm\bb}\bc)\,\ba
   =(\ip\bb{\bOm^{t}\bc})\,\ba=(\ba\otimes\bb)\bOm^{t}\bc
   =-\bA\bOm\bc ,
\]
the second equality being the definition of the transpose and the last
using $\bOm^{t}=-\bOm$. Hence
$\ba\otimes(\bOm\bb)=-\bA\bOm$. Adding the two gives
$\dot\bA=\bOm\bA-\bA\bOm$, and substituting into \eqr{6.68} gives
$\mathring\bA=\Ob$.
\end{proof}

The asymmetry between the two signs in \eqr{6.68} is therefore nothing
but the transpose in the second computation meeting the skewness of
$\bOm$. The first leg of the dyad is acted on from the left and the
second from the right, and $\bOm^{t}=-\bOm$ turns the second action into
a subtraction.

This yields a second proof of $\mathring{\vek{J}}=\Ob$, one that never
passes through \eqr{6.66}. Read \eqr{6.65} ingredient by ingredient. The
measure $\rho\,\dl v=\rho_{\kappa}\,\dl V$ does not depend on $t$, by
\eqr{6.18}. The scalar $\ip{\vek\pi}{\vek\pi}$ does not either, since
\[
   \big(\ip{\vek\pi}{\vek\pi}\big)^{\textstyle\cdot}
   =2\ip{\vek\pi}{\dot{\vek\pi}}
   =2\ip{\vek\pi}{\bom\times\vek\pi}=0
\]
by \eqr{6.54}, a cross product $\bom\times\vek\pi$ being perpendicular
to $\vek\pi$. The identity satisfies
$\mathring\I=\dot\I-\bOm\I+\I\bOm=\Ob$, its material derivative being
zero and the two correction terms cancelling. And $\vek\pi$ is fixed in
the body, $\mathring{\vek\pi}=\bze$, so
Proposition~\ref{prop:ch6jaumann} applies to
$\vek\pi\otimes\vek\pi$. Every ingredient of \eqr{6.65} thus has
vanishing co-rotational rate, and so, term by term, does the integral.

Thus the rotational momentum balance \eqr{6.51} combines with
\eqr{6.64} to give
\begin{equation}
\begin{aligned}
   \bm(S,t)=(\vek{J}\bom)^{\textstyle\cdot}
   &=\dot{\vek{J}}\bom+\vek{J}\dot\bom\\
   &=\big(\bOm\vek{J}-\vek{J}\bOm\big)\bom+\vek{J}\dot\bom .
\end{aligned}
   \tag{6.69}\label{eq:6.69}
\end{equation}
This simplifies further on noting that $\bOm\bom=\bom\times\bom$
vanishes, so that the term $\vek{J}\bOm\bom$ disappears, while
$\bOm\vek{J}\bom=\bom\times(\vek{J}\bom)$ by the definition of the
axial vector. Thus
\begin{equation}
   \bm(S,t)=\bom\times\vek{J}\bom+\vek{J}\dot\bom ,
   \qquad\text{with}\quad\dot\bom=\mathring{\bom} ,
   \tag{6.70}\label{eq:6.70}
\end{equation}
the last identification being \eqr{6.60} applied to $\bu=\bom$, whose
correction term $\bom\times\bom$ vanishes.

Equations \eqr{6.70} and \eqr{6.44} are the classical equations of
rigid-body dynamics, determining the motion of the center of mass and
the rotation in terms of the net force and torque about the mass
center. Conversely, they yield the net force and torque required to
effect a prescribed motion of the center of mass and a prescribed
rotation.

\begin{remark}[where the gyroscopic term comes
from]\label{rem:ch6gyro}
The term $\bom\times\vek{J}\bom$ in \eqr{6.70} is the one that makes
rigid-body dynamics unlike particle dynamics, and \eqr{6.69} shows
exactly where it originates: in $\dot{\vek{J}}$, the fact that the
inertia tensor changes in the inertial frame even though the body is
rigid. Had $\vek{J}$ been a constant tensor the equation would read
$\bm=\vek{J}\dot\bom$, the rotational analogue of $\vek{f}=m\ddot
{\bar{\bx}}$, and a torque-free body would spin at constant $\bom$.

It does not, in general. Set $\bm=\bze$ in \eqr{6.70}:
\[
   \vek{J}\dot\bom=-\bom\times\vek{J}\bom ,
\]
which is non-zero whenever $\bom$ is not a principal direction of
$\vek{J}$. That is the tennis racket theorem and the precession of a
free top, and it is a consequence of \eqr{6.66} alone. The
co-rotational derivative was introduced for precisely this reason: in
the body frame the inertia tensor is constant, so \eqr{6.70} written
with $\mathring{\bom}$ becomes Euler's equations with constant
coefficients, at the price of the extra term.
\end{remark}

\begin{figure}[htbp]\centering
\renewcommand{\thefigure}{A}
\begin{tikzpicture}[scale=1.25,
   ar/.style={-{Stealth[length=2.5mm]},thick},
   ax/.style={-{Stealth[length=2.2mm]},semithick,gray!70}]
% the section of the inertia ellipsoid in the (e1*,e2*) plane, J1 < J2,
% so the ellipse is longer along the axis of smaller inertia
\draw[guide] (0,0) ellipse (2.0 and 1.35);
\draw[ax] (0,0)--(2.35,0); \node[font=\small,right] at (2.38,0)
   {$\be^{*}_{1}$};
\draw[ax] (0,0)--(0,1.70); \node[font=\small,above] at (0,1.73)
   {$\be^{*}_{2}$};
% omega at 35 degrees, J omega at 57 degrees for J1 = 1, J2 = 2.2
\draw[ar] (0,0)--(1.393,0.975);
\node[font=\small,anchor=north west] at (1.40,0.97) {$\bom$};
\draw[ar] (0,0)--(1.116,1.719);
\node[font=\small,anchor=south west] at (1.09,1.73) {$\vek{J}\bom$};
\draw[gray!75,thin] (35:0.95) arc (35:57:0.95);
\node[font=\footnotesize] at (46:1.20) {$\vartheta$};
% the gyroscopic term points out of the plane
\draw[semithick] (-1.30,-0.72) circle (0.135);
\fill (-1.30,-0.72) circle (0.038);
\node[font=\small,anchor=west] at (-1.12,-0.72)
   {$\bom\times\vek{J}\bom$};
\node[font=\footnotesize,text=gray!80,anchor=north] at (0,-1.55)
   {$J_{1}<J_{2}$};
\end{tikzpicture}
\caption{Why \eqr{6.70} carries a term that $\vek{f}=m\ddot{\bar{\bx}}$
does not. The rotational momentum $\vek{J}\bom$ is not parallel to
$\bom$ unless $\bom$ lies along a principal axis: the inertia tensor
weights the components of $\bom$ unequally, and tilts the result toward
the axis of larger moment of inertia. The angle $\vartheta$ between them
is what the gyroscopic term $\bom\times\vek{J}\bom$ measures, and that
term points out of the plane of the two vectors. The dashed curve is the
section of $\ip\bu{\vek{J}\bu}=\text{const}$, elongated along the axis
of \emph{smaller} inertia because the quadratic form is smaller there.}
\label{fig:ch6inertia}
\end{figure}

%-----------------------------------------------------------------------
\subsection*{Euler's equations in the principal basis}
%-----------------------------------------------------------------------

Equation \eqr{6.70} is a vector equation in the inertial frame whose
coefficient tensor $\vek{J}$ depends on time. Resolving it on a basis
fixed in the body removes that dependence completely, and the classical
form of rigid-body dynamics appears.

By \eqr{6.67} the referential inertia tensor $\vek{J}_{\kappa}$ is
symmetric, and Problem 4 shows it is positive definite. The spectral
theorem, Theorem~\ref{thm:spectral}, therefore supplies an orthonormal
basis $\{\hbE_{i}\}$ of $\hVt$ and positive numbers $J_{i}$ with
\[
   \vek{J}_{\kappa}=\sum_{i}J_{i}\,\hbE_{i}\otimes\hbE_{i} .
\]
The $J_{i}$ are the \emph{principal moments of inertia} and the
$\hbE_{i}$ the \emph{principal axes of inertia}. This basis is available
for use in \eqr{6.58}, and the choice pays at once. Carrying it through
\eqr{6.66} with the dyad rule of Lemma~\ref{lem:dyadprod},
\[
   \vek{J}=\bQ\vek{J}_{\kappa}\bQ^{t}
   =\sum_{i}J_{i}\,(\bQ\hbE_{i})\otimes(\bQ\hbE_{i})
   =\sum_{i}J_{i}\,\be^{*}_{i}\otimes\be^{*}_{i} ,
\]
by \eqr{6.58}. So the principal axes of the spatial inertia tensor are
the rotating basis vectors themselves, and its principal values are the
same constants $J_{i}$ at every instant. All the time dependence that
Remark~\ref{rem:ch6gyro} blamed for the gyroscopic term sits in the
dyads, none of it in the numbers.

Write $\bom=\omega_{i}\be^{*}_{i}$. Then
\[
   \vek{J}\bom=\sum_{i}J_{i}\omega_{i}\be^{*}_{i} ,
\]
and, because $\dot\bom=\mathring{\bom}$ by \eqr{6.70}$_{2}$ and
$\mathring{\bom}=\dot\omega_{i}\be^{*}_{i}$ by \eqr{6.61},
\[
   \vek{J}\dot\bom=\sum_{i}J_{i}\dot\omega_{i}\be^{*}_{i} .
\]
The gyroscopic term is a cross product of two vectors whose components
on the orthonormal triad $\{\be^{*}_{i}\}$ are now both known, so it may
be written out:
\[
   \bom\times\vek{J}\bom=
   \big[(J_{3}-J_{2})\omega_{2}\omega_{3}\big]\be^{*}_{1}
  +\big[(J_{1}-J_{3})\omega_{3}\omega_{1}\big]\be^{*}_{2}
  +\big[(J_{2}-J_{1})\omega_{1}\omega_{2}\big]\be^{*}_{3} ,
\]
the first component being $\omega_{2}(J_{3}\omega_{3})
-\omega_{3}(J_{2}\omega_{2})$ and the other two following by cyclic
permutation of $123$. Equating components in \eqr{6.70} gives
\emph{Euler's equations},
\[
\begin{aligned}
   J_{1}\dot\omega_{1}+(J_{3}-J_{2})\,\omega_{2}\omega_{3}&=m_{1},\\
   J_{2}\dot\omega_{2}+(J_{1}-J_{3})\,\omega_{3}\omega_{1}&=m_{2},\\
   J_{3}\dot\omega_{3}+(J_{2}-J_{1})\,\omega_{1}\omega_{2}&=m_{3},
\end{aligned}
\]
where $m_{i}=\ip{\bm(S,t)}{\be^{*}_{i}}$ are the components of the
torque about the center of mass on the same rotating triad. Every
coefficient is a constant. That is what the body basis was for, and it
is the reason the co-rotational derivative was introduced in \eqr{6.61}
rather than being avoided.

%-----------------------------------------------------------------------
\subsection*{The torque-free body}
%-----------------------------------------------------------------------

Set $\bm(S,t)=\bze$. Two scalars are then constant in time, and between
them they explain the assertion of Remark~\ref{rem:ch6gyro} that a free
body does not spin steadily.

The first is immediate from \eqr{6.51}: $\dot{\bh}=\bze$, so the
rotational momentum $\bh=\vek{J}\bom$ is a fixed vector of $\Vt$, and in
particular $|\bh|$ is constant.

The second is the kinetic energy of the rotation,
$T=\tfrac12\ip\bom{\vek{J}\bom}$. Differentiate, using the symmetry of
$\vek{J}$ to combine the first and third terms:
\[
   \dot T=\tfrac12\ip{\dot\bom}{\vek{J}\bom}
   +\tfrac12\ip{\bom}{\dot{\vek{J}}\bom}
   +\tfrac12\ip{\bom}{\vek{J}\dot\bom}
   =\ip{\bom}{\vek{J}\dot\bom}
   +\tfrac12\ip{\bom}{\dot{\vek{J}}\bom} .
\]
The last term vanishes. From the derivation of \eqr{6.68},
$\dot{\vek{J}}=\bOm\vek{J}-\vek{J}\bOm$, so
\[
   \ip{\bom}{\big(\bOm\vek{J}-\vek{J}\bOm\big)\bom}
   =\ip{\bom}{\bOm\vek{J}\bom}-\ip{\bom}{\vek{J}\bOm\bom}
   =-\ip{\bOm\bom}{\vek{J}\bom}-0=0 ,
\]
the first term because $\bOm$ is skew, so that moving it across the
inner product changes the sign, and the second because
$\bOm\bom=\bom\times\bom=\bze$. Hence
$\dot T=\ip\bom{\vek{J}\dot\bom}$. Now take the inner product of
\eqr{6.70} with $\bom$:
\[
   \ip\bom{\bm(S,t)}=\ip\bom{\bom\times\vek{J}\bom}
   +\ip\bom{\vek{J}\dot\bom}=\ip\bom{\vek{J}\dot\bom}=\dot T ,
\]
since $\bom\times\vek{J}\bom$ is perpendicular to $\bom$. So in general
\[
   \dot T=\ip{\bom}{\bm(S,t)} ,
\]
the rotational power of the torque about the mass center, and $T$ is
constant when that torque vanishes. This is the rotational counterpart
of the mechanical energy balance \eqr{6.150}, and it is exact rather
than approximate because a rigid body stores no strain energy: its
stress power vanished at the end of Section~\ref{sec:ch6rigid}.

On the body basis the two constants read
\[
   |\bh|^{2}=J_{1}^{2}\omega_{1}^{2}+J_{2}^{2}\omega_{2}^{2}
   +J_{3}^{2}\omega_{3}^{2} ,
   \qquad
   2T=J_{1}\omega_{1}^{2}+J_{2}\omega_{2}^{2}+J_{3}\omega_{3}^{2} ,
\]
two ellipsoids in the space of the components $\omega_{i}$. The motion
of $\bom$ relative to the body is confined to their intersection, which
is a closed curve. Whether that curve stays near an axis or wanders far
from it is the question of stability, and Euler's equations answer it in
three lines.

Order the moments $J_{1}<J_{2}<J_{3}$ and set $\bm=\bze$. Steady
rotation about the first axis is $\bom=\Omega\be^{*}_{1}$, and it does
solve the equations, since then $\omega_{2}=\omega_{3}=0$ kills every
product. Perturb it: put $\omega_{1}=\Omega$ and treat $\omega_{2}$,
$\omega_{3}$ as small, so that products of two small quantities are
discarded. The second and third equations become
\[
   J_{2}\dot\omega_{2}=(J_{3}-J_{1})\Omega\,\omega_{3} ,
   \qquad
   J_{3}\dot\omega_{3}=(J_{1}-J_{2})\Omega\,\omega_{2} ,
\]
and differentiating the first and substituting the second,
\[
   \ddot\omega_{2}
   =\frac{(J_{3}-J_{1})(J_{1}-J_{2})}{J_{2}J_{3}}\,\Omega^{2}\,
    \omega_{2} .
\]
With $J_{1}$ the smallest, $J_{3}-J_{1}>0$ and $J_{1}-J_{2}<0$, so the
coefficient is negative and $\omega_{2}$ oscillates: the perturbation
stays small. The same computation about the third axis gives
$(J_{1}-J_{3})(J_{3}-J_{2})<0$ again, so that axis is stable too. About
the intermediate axis, however, $\bom=\Omega\be^{*}_{2}$ gives
\[
   \ddot\omega_{1}
   =\frac{(J_{2}-J_{3})(J_{1}-J_{2})}{J_{1}J_{3}}\,\Omega^{2}\,
    \omega_{1} ,
\]
in which both factors in the numerator are negative and the coefficient
is positive. The perturbation grows exponentially. That is the tennis
racket theorem, promised in Remark~\ref{rem:ch6gyro}, and every step of
it came from \eqr{6.66}: the inertia tensor is fixed in the body, not in
space.

%=======================================================================
\section{Forces of mutual interaction in the body and the local
equations of motion}\label{sec:ch6local}
%=======================================================================

So far the traction $\bt$ has been an undetermined field on the
boundary of whatever part is under consideration. This section
determines its structure completely, in three steps. First, the
traction that one part of a body exerts on an adjoining part is equal
and opposite to the traction exerted back; that is a theorem, not an
axiom. Second, the traction at a point depends on the cutting surface
only through its unit normal there; that is Noll's theorem. Third, that
dependence is linear, so there is a tensor field $\bT$ with
$\bt=\bT\bn$; that is Cauchy's theorem. Only then can the linear
momentum balance be localized.

%-----------------------------------------------------------------------
\subsection{Action and reaction}\label{sec:ch6action}
%-----------------------------------------------------------------------

Consider an arbitrary part $P_{t}$ of the region $\kappa_{t}$ occupied
by the body at time $t$. Let $\sigma$ be a surface dividing it into two
sub-parts $P_{t1}$ and $P_{t2}$ (Figure~\ref{fig:ch6cut}). The letter
$\sigma$ is the book's, and in this subsection it denotes that surface
and not the supply rate of \eqr{6.1}. Let
\[
   \partial P'_{t}=\partial P_{t1}\cap\partial P_{t},\qquad
   \partial P''_{t}=\partial P_{t2}\cap\partial P_{t},
\]
with $\partial P_{t1}=\partial P'_{t}\cup\sigma$ and
$\partial P_{t2}=\partial P''_{t}\cup\sigma$. Then
$P_{t}=P_{t1}\cup P_{t2}$ and
$\partial P_{t}=\partial P'_{t}\cup\partial P''_{t}$. The last
statement is the one that will matter: the dividing surface is part of
the boundary of each half and of the boundary of neither the whole.

The forces acting on $P_{t1}$ are the body force and the contact force.
We express the former as
\begin{equation}
   \vek{f}_{b1}=\int_{P_{t1}}\rho\big(\bar{\bb}+\bb_{12}\big)\,\dl v ,
   \tag{6.71}\label{eq:6.71}
\end{equation}
where $\bar{\bb}$ is due to sources external to $P_{t}$ and
$\bb_{12}$ is the body force per unit mass exerted on $P_{t1}$ by
$P_{t2}$. The latter may be due to the gravitational attraction of the
two parts for one another, for example. The contact force is
\begin{equation}
   \vek{f}_{c1}=\int_{\partial P_{t1}}\bt\,\dl a
   =\int_{\partial P'_{t}}\bt\,\dl a+\int_{\sigma}\bt_{12}\,\dl a ,
   \tag{6.72}\label{eq:6.72}
\end{equation}
where $\bt_{12}$ is the traction transmitted to $P_{t1}$ by $P_{t2}$
across $\sigma$. The net force applied to $P_{t1}$ is
\begin{equation}
   \vek{f}_{1}=\vek{f}_{b1}+\vek{f}_{c1}.
   \tag{6.73}\label{eq:6.73}
\end{equation}
Similarly, the net force acting on $P_{t2}$ is
\begin{equation}
   \vek{f}_{2}=\vek{f}_{b2}+\vek{f}_{c2},
   \tag{6.74}\label{eq:6.74}
\end{equation}
where
\begin{equation}
   \vek{f}_{b2}=\int_{P_{t2}}\rho\big(\bar{\bb}+\bb_{21}\big)\,\dl v
   \qquad\text{and}\qquad
   \vek{f}_{c2}=\int_{\partial P''_{t}}\bt\,\dl a
   +\int_{\sigma}\bt_{21}\,\dl a ,
   \tag{6.75}\label{eq:6.75}
\end{equation}
wherein $\bb_{21}$ is due to $P_{t1}$ and $\bt_{21}$ is the traction
exerted by $P_{t1}$ on $P_{t2}$.

The linear momenta of the parts $P_{t1}$ and $P_{t2}$ are
\begin{equation}
   \vek{l}_{1}=\int_{P_{t1}}\rho\bv\,\dl v
   \qquad\text{and}\qquad
   \vek{l}_{2}=\int_{P_{t2}}\rho\bv\,\dl v ,
   \tag{6.76}\label{eq:6.76}
\end{equation}
respectively, and the momentum of $P_{t}=P_{t1}\cup P_{t2}$ is
$\vek{l}=\vek{l}_{1}+\vek{l}_{2}$, the two pieces overlapping only on
$\sigma$, a set of zero volume. Euler's postulate \eqr{6.34}, applied
to $P_{t1}$ and $P_{t2}$ separately, yields
\begin{equation}
\begin{aligned}
   \int_{P_{t1}}\rho\big(\bar{\bb}+\bb_{12}\big)\,\dl v
   +\int_{\partial P'_{t}}\bt\,\dl a+\int_{\sigma}\bt_{12}\,\dl a
   &=\frac{\dl}{\dl t}\vek{l}_{1}
   \qquad\text{and}\\[1mm]
   \int_{P_{t2}}\rho\big(\bar{\bb}+\bb_{21}\big)\,\dl v
   +\int_{\partial P''_{t}}\bt\,\dl a+\int_{\sigma}\bt_{21}\,\dl a
   &=\frac{\dl}{\dl t}\vek{l}_{2}.
\end{aligned}
   \tag{6.77}\label{eq:6.77}
\end{equation}
On adding these and recalling that
$\partial P_{t}=\partial P'_{t}\cup\partial P''_{t}$, so that the two
integrals of $\bt$ over the outer boundary combine into one, we obtain
\begin{equation}
\begin{aligned}
   \int_{P_{t}}\rho\bar{\bb}\,\dl v
   +\int_{P_{t1}}\rho\bb_{12}\,\dl v
   +\int_{P_{t2}}\rho\bb_{21}\,\dl v
   +\int_{\partial P_{t}}\bt\,\dl a
   +\int_{\sigma}\big(\bt_{12}+\bt_{21}\big)\,\dl a
   =\frac{\dl}{\dl t}\vek{l} .
\end{aligned}
   \tag{6.78}\label{eq:6.78}
\end{equation}
On the other hand, Euler's postulate applied to the whole of $P_{t}$
reads
\begin{equation}
   \int_{P_{t}}\rho\bar{\bb}\,\dl v+\int_{\partial P_{t}}\bt\,\dl a
   =\frac{\dl}{\dl t}\vek{l} ,
   \tag{6.79}\label{eq:6.79}
\end{equation}
since the only body force acting on $P_{t}$ from outside is
$\bar{\bb}$, the mutual attractions $\bb_{12}$ and $\bb_{21}$ being
internal to it. Subtracting \eqr{6.79} from \eqr{6.78}, the two terms
they have in common cancel and
\begin{equation}
   \int_{P_{t2}}\rho\bb_{21}\,\dl v+\int_{\sigma}\bt_{21}\,\dl a
   =-\left[\int_{P_{t1}}\rho\bb_{12}\,\dl v
   +\int_{\sigma}\bt_{12}\,\dl a\right].
   \tag{6.80}\label{eq:6.80}
\end{equation}
The force applied to $P_{t2}$ by $P_{t1}$ is thus opposite to the force
applied to $P_{t1}$ by $P_{t2}$.

\begin{remark}[action and reaction is a theorem
here]\label{rem:ch6actionreaction}
In particle mechanics the equality of action and reaction is Newton's
third law, an independent axiom. For continua it is a
\emph{consequence} of the balance of linear momentum, and \eqr{6.79} is
where the consequence comes from.

The substance of the argument is that Euler's postulate is asserted for
\emph{every} part, and that the momentum, the external body force and
the external traction of the whole are the sums of those of the halves.
Given that, the internal interactions must cancel, because they appear
in the two half-statements and are absent from the whole-statement. Any
theory in which they did not cancel would be inconsistent: applying the
same axiom at two different levels of subdivision would give two
different answers.

The word \emph{external} carries the weight. Equation \eqr{6.71}
separates $\bar{\bb}$ from $\bb_{12}$ by fiat, and \eqr{6.79} then
asserts that only $\bar{\bb}$ acts on the whole. That is the modelling
assumption: interactions come in pairs internal to the body, plus
whatever the outside world applies. It is the continuum version of the
statement that a system cannot lift itself by its own bootstraps.
\end{remark}

\bookfig{4}\begin{figure}[htbp]\centering
\begin{tikzpicture}[scale=1.05,
   ar/.style={-{Stealth[length=2.2mm]},thick}]
% the part P_t
\draw[thick] plot[smooth cycle,tension=0.82] coordinates
  {(0,1.75) (1.95,1.25) (2.35,-0.45) (0.85,-1.85) (-1.35,-1.55)
   (-2.05,0.15)};
% the dividing surface
\draw[very thick,red!65!black]
   (-1.92,-0.62) .. controls (-0.55,-0.05) and (0.60,0.35) .. (2.28,0.02);
\node[font=\small,text=red!65!black,anchor=west] at (2.05,0.45) {$\sigma$};
% labels for the two sub-parts
\node[font=\small] at (-0.20,0.85) {$P_{t1}$};
\node[font=\small] at (0.30,-1.00) {$P_{t2}$};
% outer boundary pieces
\node[font=\footnotesize,text=gray!85,anchor=west] at (1.95,1.55)
   {$\partial P'_{t}$};
\draw[gray!70,-{Stealth[length=1.6mm]}] (1.93,1.50) -- (1.18,1.42);
\node[font=\footnotesize,text=gray!85,anchor=west] at (1.95,-1.55)
   {$\partial P''_{t}$};
\draw[gray!70,-{Stealth[length=1.6mm]}] (1.93,-1.50) -- (1.20,-1.20);
% the two opposed tractions on sigma
\coordinate (Q) at (0.05,0.10);
\fill (Q) circle (1.2pt);
\draw[ar,blue!55!black] (Q) -- ++(0.30,0.78)
   node[above,font=\small,text=blue!55!black] {$\bt_{12}$};
\draw[ar,blue!55!black] (Q) -- ++(-0.30,-0.78)
   node[below,font=\small,text=blue!55!black] {$\bt_{21}$};
\end{tikzpicture}
\caption{An arbitrary part of a body divided into two sub-parts by an
interior surface $\sigma$. Euler's postulate holds for $P_{t1}$, for
$P_{t2}$ and for $P_{t}$; the first two together must reproduce the
third, and the difference is exactly \eqr{6.80}. Shrinking $P_{t}$ onto
$\sigma$ kills the volume terms and leaves Cauchy's lemma
$\bt_{21}=-\bt_{12}$. Note that $\sigma$ belongs to
$\partial P_{t1}$ and to $\partial P_{t2}$ but not to $\partial P_{t}$:
that is the whole of the bookkeeping.}
\label{fig:ch6cut}
\end{figure}

\subsection*{Collapsing the parts onto the surface}

Equation \eqr{6.80} is a statement about volumes and a surface
together. To isolate the surface part, consider a sequence of parts
$P_{t}$, all with the same interior surface $\sigma$ and collapsing
onto it in the limit. Each of these parts has associated sub-parts
$P_{t1}$ and $P_{t2}$ whose volumes tend to zero, while $\sigma$ is
held fixed throughout. Assuming $|\rho\bb_{21}|$ and $|\rho\bb_{12}|$
to be bounded, say by a constant $k$, the volume integrals in
\eqr{6.80} are bounded in magnitude by $k\,V(P_{t1})$ and
$k\,V(P_{t2})$ and therefore vanish in the limit, yielding
\begin{equation}
   \int_{\sigma}\big(\bt_{21}+\bt_{12}\big)\,\dl a=\bze .
   \tag{6.81}\label{eq:6.81}
\end{equation}
Substituting this back into \eqr{6.80}, the two surface integrals there
cancel and we conclude that
\begin{equation}
   \int_{P_{t2}}\rho\bb_{21}\,\dl v=-\int_{P_{t1}}\rho\bb_{12}\,\dl v .
   \tag{6.82}\label{eq:6.82}
\end{equation}
So the two statements separate: the contact interactions balance among
themselves, and so do the body-force interactions.

Because $P_{t}$ is arbitrary, $\sigma$ is an arbitrary surface in
$\kappa_{t}$, and \eqr{6.81} holds for every one of them. Assuming the
integrand to be continuous, we may invoke a version of the localization
theorem for surface integrals, in the manner already used to pass from
\eqr{6.4} to \eqr{6.6}, to arrive at \emph{Cauchy's lemma}:
\begin{equation}
   \boxed{\ \bt_{21}=-\bt_{12}.\ }
   \tag{6.83}\label{eq:6.83}
\end{equation}

\begin{remark}[why the body forces had to be assumed
bounded]\label{rem:ch6bounded}
The hypothesis that $|\rho\bb_{12}|$ is bounded is the only analytic
input in the argument, and it is doing exactly one job: making the
volume terms of \eqr{6.80} vanish faster than the surface terms as the
parts collapse. If the mutual body force were allowed to become
unbounded as the two parts approach one another, as happens for point
masses in Newtonian gravitation at zero separation, then the volume
integrals need not vanish and \eqr{6.81} fails.

The same competition between a volume and a surface decides Noll's
theorem and Cauchy's theorem below. It is the structural fact of the
whole section: \emph{a surface integral over a shrinking region is of
one order higher than a volume integral over it}, so in any limit that
divides by the surface measure, the bulk terms disappear and only the
surface terms survive. That is why all three theorems are proved by
shrinking something.
\end{remark}

%-----------------------------------------------------------------------
\subsection{Cauchy's hypothesis and Noll's theorem}\label{sec:ch6noll}
%-----------------------------------------------------------------------

Let $\bt_{s}(\bx,t)$ be the traction acting at $\bx\in s$, a surface in
$\kappa_{t}$. The subscript conveys the notion that the traction
depends on the surface $s$ and not merely on the point $\bx$. That it
must depend on something beyond $\bx$ is clear from a simple thought
experiment. Take a bar carrying an axial force $F$ and cut it at a
point $\bx$, once by a plane $s_{1}$ perpendicular to the axis and once
by a plane $s_{2}$ inclined to it. The same force $F$ is transmitted in
both cases, but it is spread over different areas, so the force per
unit area differs. The traction at $\bx$ therefore cannot be a function
of $\bx$ alone.

Since the traction is defined locally at points of the surface,
logically it should depend on a local measure of the surface
orientation, namely a unit normal $\bn$ to $s$ at $\bx$. This
assumption is the content of the famous \emph{Cauchy hypothesis}: there
exists a function $\bt(\bx,t;\cdot)$ such that
\begin{equation}
   \bt_{s}(\bx,t)=\bt(\bx,t;\bn).
   \tag{6.84}\label{eq:6.84}
\end{equation}
Cauchy's lemma \eqr{6.83} then assumes the form
\begin{equation}
   \bt(\bx,t;-\bn)=-\bt(\bx,t;\bn),
   \tag{6.85}\label{eq:6.85}
\end{equation}
since two parts adjoining along a surface see the same surface with
opposite exterior normals.

\begin{remark}[what \eqr{6.84} rules out]\label{rem:ch6cauchyhyp}
Equation \eqr{6.84} looks weak and is not. What it forbids is the
following. The traction might have depended on the
curvature of $s$ at $\bx$, so that a cut along a sphere transmitted a
different force from a cut along its tangent plane. It might have
depended on the diameter of the part being cut off, or on the shape of
$s$ far from $\bx$. Each of these is a coherent physical possibility,
and each is excluded by \eqr{6.84}.

That is why Noll's theorem matters. Cauchy's hypothesis has occupied a
central place in continuum mechanics since it was proposed in the early
nineteenth century, and in the mid-twentieth century Noll strengthened
it to the status of a proved theorem: it is a consequence of Euler's
linear momentum balance, given only continuity and boundedness of the
fields involved. Nothing has to be assumed about how the traction
responds to the shape of the cut, because the balance law already
decides.
\end{remark}

\subsection*{Sketch of Noll's theorem}

With reference to Figure~\ref{fig:ch6noll}, let $s_{1}$ and $s_{2}$ be
two surfaces in $\kappa_{t}$ intersecting at a point with position
$\bx_{0}$, and suppose they share a common tangent plane $T$, with unit
normal $\bn$, at this point.

Let $P_{t1}$ be the part of the body bounded by
$\partial P_{t1}=d_{1}\cup f_{1}\cup e$, where $d_{1}$ is a subset of
$s_{1}$, $f_{1}$ is a piece of the lateral surface of a circular
cylinder of radius $R$ and axis $\bn$, and $e$ is the part of the
surface of that cylinder common to both $\partial P_{t1}$ and
$\partial P_{t2}$. The part $P_{t2}$, not shown separately, is bounded
above by $s_{2}$, with $\partial P_{t2}=d_{2}\cup f_{2}\cup e$ and
corresponding definitions of $d_{2}$ and $f_{2}$.

\subsection*{The construction in coordinates}

The estimates that follow are the whole content of the argument, and
they depend on how deep the shared face $e$ is placed, so the geometry
is fixed precisely first.

Put the origin at $\bx_{0}$, let $T$ be the plane through it with
normal $\bn$, and let $\vek{\xi}$ denote position in $T$ measured from
$\bx_{0}$. Near $\bx_{0}$ each surface is a graph over $T$,
\[
   s_{\alpha}:\quad z=g_{\alpha}(\vek{\xi}),\qquad
   g_{\alpha}(\bze)=0,\qquad \nabla g_{\alpha}(\bze)=\bze ,
\]
the first condition saying that $s_{\alpha}$ passes through $\bx_{0}$
and the second that its tangent plane there is $T$. Differentiability
of $g_{\alpha}$ at $\bze$ then says precisely that
\begin{equation}
   g_{\alpha}(\vek{\xi})=o\big(|\vek{\xi}|\big)
   \qquad\text{as }\vek{\xi}\to\bze ,
   \tag{$\ast$}\label{eq:ch6nolltan}
\end{equation}
and that single statement is the hypothesis the theorem runs on.

Now build the pillbox. Let the cylinder be $|\vek{\xi}|\le R$ about the
axis through $\bx_{0}$ along $\bn$, and put the shared face at the depth
\begin{equation}
   h(R)=\max_{\alpha\in\{1,2\}}\ \Big(
   \max_{|\vek{\xi}|\le R}\big|g_{\alpha}(\vek{\xi})\big|\Big),
   \tag{$\ast\ast$}\label{eq:ch6nollhdef}
\end{equation}
in which the inner maximum is taken over the closed disc
$D_{R}=\{\vek{\xi}\in T:|\vek{\xi}|\le R\}$ and the outer one over the
two-element index set $\{1,2\}$. Both are attained and finite: $D_{R}$
is closed and bounded, hence compact, and $|g_{\alpha}|$ is continuous
on it, so the extreme value theorem applies; the outer maximum is that
of two real numbers. Thus $h(R)$ is a well-defined non-negative number
for each $R$, and by construction
\[
   -h(R)\ \le\ g_{\alpha}(\vek{\xi})\ \le\ h(R)
   \qquad\text{for all }\vek{\xi}\in D_{R}\text{ and }\alpha=1,2 ,
\]
so the disc $e$ of radius $R$ placed at height $-h(R)$ lies below both
caps, both parts $P_{t1}$ and $P_{t2}$ lie above it, and
$\partial P_{t\alpha}=d_{\alpha}\cup f_{\alpha}\cup e$ closes up.

The one property of $h$ that the proof uses is
\begin{equation}
   h(R)=o(R)\qquad\text{as }R\to0^{+},
   \tag{$\ast\ast\ast$}\label{eq:ch6nollh}
\end{equation}
that is, $h(R)/R\to0$, in the sense of Definition~\ref{def:landau}. It
is \eqref{eq:ch6nolltan} read on the disc. Let $\eta>0$ be given.
Tangency supplies $r_{0}>0$ with
\[
   \big|g_{\alpha}(\vek{\xi})\big|\le\eta\,|\vek{\xi}|
   \qquad\text{whenever }|\vek{\xi}|\le r_{0},\ \alpha=1,2 .
\]
For $0<R\le r_{0}$ every $\vek{\xi}\in D_{R}$ has $|\vek{\xi}|\le R$,
so the right-hand side is at most $\eta R$; taking the maximum over
$D_{R}$ and then over $\alpha$ preserves the bound, and
\[
   h(R)\le\eta R ,\qquad\text{that is}\qquad
   \frac{h(R)}{R}\le\eta\quad\text{for }0<R\le r_{0} .
\]
As $\eta$ was arbitrary, $h(R)/R\to0$.

\begin{remark}[the depth of $e$ is not free]\label{rem:ch6nolldepth}
The figure in the book marks the height of the cylinder as $R$, the
same as its radius. That cannot be the intended construction, because
it contradicts \eqr{6.86}$_{2}$.

With the face at a constant depth $h$, the wall of the pillbox reaches
from $-h$ up to $g_{\alpha}$, so its height nowhere exceeds
$h+M_{\alpha}(R)$, where
\[
   M_{\alpha}(R)=\max_{|\vek{\xi}|=R}\big|g_{\alpha}(\vek{\xi})\big|
\]
is the largest departure of $s_{\alpha}$ from $T$ on the rim of the
disc, attained because the rim is compact and $|g_{\alpha}|$
continuous. By \eqref{eq:ch6nolltan}, $M_{\alpha}(R)=o(R)$, so
\[
   A(f_{\alpha})\ \le\ 2\pi R\big[h+o(R)\big],
   \qquad
   \frac{A(f_{\alpha})}{\pi R^{2}}\ \le\ \frac{2h}{R}+o(1) .
\]
Taking $h=R$ leaves the ratio bounded below by $2$: the wall is
\emph{twice} the cap in the limit, not negligible beside it, and
\eqr{6.90} does not follow. Taking $h=h(R)$ as above gives
$2h(R)/R\to0$ by \eqref{eq:ch6nollh} and the wall disappears.

A numerical check on paraboloids $g_{\alpha}=\tfrac12c_{\alpha}
|\vek{\xi}|^{2}$, which are tangent to $T$ at $\bx_{0}$, makes the
contrast plain. With $h=h(R)$ the ratio $A(f_{1})/A(d_{1})$ runs
$0.209,\ 0.105,\ 0.052$ as $R$ runs $0.1,\ 0.05,\ 0.025$, halving with
$R$; with $h=R$ the same ratio runs $2.121,\ 2.063,\ 2.032$ and settles
on $2$.

So the pillbox must be flat, and it must flatten faster than it
narrows. That is the geometric content of tangency, and
Figure~\ref{fig:ch6noll} is drawn to scale in that respect.
\end{remark}

Denote by $A(s)$ the area of a surface $s$ and by $V(P)$ the volume of
a region $P$. Then
\begin{equation}
   A(d_{\alpha})=\pi R^{2}+o(R^{2})
   \qquad\text{and}\qquad
   A(f_{\alpha})=o(R^{2});\qquad\alpha=1,2 ,
   \tag{6.86}\label{eq:6.86}
\end{equation}
and
\begin{equation}
   V(P_{t\alpha})=O(R^{3})=o(R^{2});\qquad\alpha=1,2 .
   \tag{6.87}\label{eq:6.87}
\end{equation}
All three follow from \eqref{eq:ch6nolltan} and
\eqref{eq:ch6nollh}, in the sense of the Landau symbols of
Definition~\ref{def:landau}, and each is a one-line integral over the
disc $|\vek{\xi}|\le R$.

\emph{The caps.} A graph $z=g_{\alpha}$ has area element
$\sqrt{1+|\nabla g_{\alpha}|^{2}}$ times the area element of $T$, so
\[
   A(d_{\alpha})=\int_{|\vek{\xi}|\le R}
   \sqrt{1+|\nabla g_{\alpha}|^{2}}\ \dl a_{\vek{\xi}} .
\]
The gradient vanishes at $\bze$ and is continuous, so
$|\nabla g_{\alpha}|=o(1)$ on the disc and the integrand is $1+o(1)$.
Integrating the constant over a disc of area $\pi R^{2}$,
\[
   A(d_{\alpha})=\pi R^{2}\big[1+o(1)\big]=\pi R^{2}+o(R^{2}),
\]
which is \eqr{6.86}$_{1}$. The cap is a disc to leading order, and the
correction is the curvature, which is second order.

\emph{The lateral strips.} The wall stands on the circle
$|\vek{\xi}|=R$, of circumference $2\pi R$, and reaches from height
$-h(R)$ up to $g_{\alpha}$, so its height nowhere exceeds
$h(R)+|g_{\alpha}|\le 2h(R)$. Therefore
\[
   A(f_{\alpha})\ \le\ 2\pi R\cdot 2h(R)=4\pi R\,h(R)=o(R^{2})
\]
by \eqref{eq:ch6nollh}, which is \eqr{6.86}$_{2}$.

This is the estimate the whole proof turns on. As a ratio rather than an
order,
\[
   \frac{A(f_{\alpha})}{A(d_{\alpha})}\ \le\ \frac{4h(R)}{R}
   \big[1+o(1)\big]\ \longrightarrow\ 0 .
\]
The wall is negligible \emph{beside the cap}, and it is negligible only
because the two surfaces are tangent. Were they to cross at a fixed
angle, so that $g_{\alpha}\approx m_{\alpha}|\vek{\xi}|$ with
$m_{\alpha}\neq0$, then $h(R)$ would be $O(R)$ rather than $o(R)$, the
ratio above would tend to a non-zero constant depending on the
$m_{\alpha}$ but not on $R$, and no amount of shrinking would remove
the wall. Figure~\ref{fig:ch6noll} shows the two cases side by side.

\emph{The volumes.} The region $P_{t\alpha}$ lies between $e$, at height
$-h(R)$, and the cap $d_{\alpha}$, at height $g_{\alpha}$, over the disc
$D_{R}$, so its height above each $\vek{\xi}$ is $g_{\alpha}+h(R)$,
which by \eqref{eq:ch6nollhdef} lies between $0$ and $2h(R)$. Hence
\begin{equation}
   V(P_{t\alpha})=\int_{D_{R}}\big[g_{\alpha}(\vek{\xi})+h(R)\big]
   \,\dl a_{\vek{\xi}}
   \ \le\ 2h(R)\int_{D_{R}}\dl a_{\vek{\xi}}
   =2\pi R^{2}h(R).
   \tag{$\S$}\label{eq:ch6nollvol}
\end{equation}

It remains to convert this into \eqr{6.87}, and the conversion is the
one place where the two Landau symbols must be told apart. Write
$C=2\pi$ and use $h(R)\le R$, which holds for all small $R$ by
\eqref{eq:ch6nollh} with $\eta=1$:
\[
   V(P_{t\alpha})\ \le\ 2\pi R^{2}h(R)\ \le\ 2\pi R^{3}
   =C\,R^{3},
\]
which is the definition of $V(P_{t\alpha})=O(R^{3})$: the ratio
$V/R^{3}$ is bounded by the \emph{single} constant $C$. Lemma
\ref{lem:landau}\textup{(ii)} with $p=3$ and $q=2$ then converts this
into a little-oh statement one power down. Explicitly, given $\eta>0$,
\[
   \frac{V(P_{t\alpha})}{R^{2}}\ \le\ \frac{C R^{3}}{R^{2}}=CR
   \ <\ \eta \qquad\text{whenever } R<\eta/C ,
\]
so the ratio $V/R^{2}$ falls below \emph{every} constant, and
\[
   V(P_{t\alpha})=O(R^{3})=o(R^{2}),
\]
which is \eqr{6.87}. The two symbols differ in exactly one quantifier:
big-oh asks for some constant bound, little-oh for every one, and one
extra power of $R$ is what upgrades the first into the second. The same
reading applies to $A(f_{\alpha})=o(R^{2})$ above, where the extra power
was supplied by $h(R)$ rather than by $R$ itself.

Relative to the cap, $\eqref{eq:ch6nollvol}$ and $A(d_{\alpha})
=\pi R^{2}[1+o(1)]$ give
\[
   \frac{V(P_{t\alpha})}{A(d_{\alpha})}
   \le\frac{2\pi R^{2}h(R)}{\pi R^{2}\big[1+o(1)\big]}
   =2h(R)\big[1+o(1)\big]\ \longrightarrow\ 0 .
\]

The balances of linear momentum for $P_{t1}$ and $P_{t2}$ are given by
\eqr{6.36}:
\begin{equation}
   \int_{P_{t1}}\rho(\dot\bv-\bb)\,\dl v
   =\int_{\partial P_{t1}}\bt_{\partial P_{t1}}\,\dl a
   \qquad\text{and}\qquad
   \int_{P_{t2}}\rho(\dot\bv-\bb)\,\dl v
   =\int_{\partial P_{t2}}\bt_{\partial P_{t2}}\,\dl a ,
   \tag{6.88}\label{eq:6.88}
\end{equation}
which, when subtracted and rearranged so that the two caps stand alone
on the left, give
\begin{equation}
\begin{aligned}
   \int_{d_{1}}\bt_{d_{1}}\,\dl a-\int_{d_{2}}\bt_{d_{2}}\,\dl a
   =\ &\int_{P_{t1}}\rho(\dot\bv-\bb)\,\dl v
       -\int_{P_{t2}}\rho(\dot\bv-\bb)\,\dl v\\
     &+\int_{f_{2}}\bt_{f_{2}}\,\dl a-\int_{f_{1}}\bt_{f_{1}}\,\dl a .
\end{aligned}
   \tag{6.89}\label{eq:6.89}
\end{equation}
The integrals over $e$ have cancelled, and the reason is the step the
construction was designed around.
Both parts lie \emph{above} $e$, by the choice of the depth $h(R)$, so
$e$ enters $\partial P_{t1}$ and $\partial P_{t2}$ as the same surface
with the same exterior normal, namely $-\bn$. The traction on it is
therefore literally the same field in the two balances, whatever that
field may be, and subtracting one balance from the other removes it.

Nothing has been assumed about its value, and in particular Cauchy's
hypothesis has not been used: at this stage the traction is still
allowed to depend on the whole surface, and the argument needs only
that $e$ is one and the same surface seen twice. Had the two parts been
placed on opposite sides of $e$, the exterior normals would have been
opposite, the two tractions would have been opposite by Cauchy's lemma
\eqr{6.83}, and subtracting would have \emph{doubled} the term instead
of cancelling it. That is the reason $e$ is pushed below both caps
rather than laid in the tangent plane.

Assuming all the integrands to be continuous, and hence bounded on
their respective compact domains of integration, the three terms on the
right of \eqr{6.89} are bounded in magnitude by a constant times
$V(P_{t1})$, $V(P_{t2})$, $A(f_{2})$ and $A(f_{1})$ respectively, all
of which are $o(R^{2})$ by \eqr{6.86}$_{2}$ and \eqr{6.87}. Hence
\begin{equation}
   \int_{d_{1}}\bt_{d_{1}}\,\dl a=\int_{d_{2}}\bt_{d_{2}}\,\dl a
   +\vek{o}(R^{2}),
   \tag{6.90}\label{eq:6.90}
\end{equation}
where $\vek{o}(\eps)$ denotes a vector whose magnitude is $o(\eps)$.
Dividing by $\pi R^{2}$ and invoking \eqr{6.86}$_{1}$, which says that
$A(d_{\alpha})$ and $\pi R^{2}$ may be interchanged at the cost of a
further $o(1)$, we obtain
\begin{equation}
   \frac{1}{A(d_{1})}\int_{d_{1}}\bt_{d_{1}}\,\dl a
   =\frac{1}{A(d_{2})}\int_{d_{2}}\bt_{d_{2}}\,\dl a
   +\frac{\vek{o}(R^{2})}{R^{2}} .
   \tag{6.91}\label{eq:6.91}
\end{equation}

Under our continuity hypothesis the mean-value theorem, essentially the
same thing as the localization theorem, is applicable and implies that
\begin{equation}
   \lim_{R\to0}\frac{1}{A(d_{\alpha})}
   \int_{d_{\alpha}}\bt_{d_{\alpha}}(\bx,t)\,\dl a
   =\bt_{d_{\alpha}}(\bx_{0},t);\qquad\alpha=1,2 .
   \tag{6.92}\label{eq:6.92}
\end{equation}
The average of a continuous function over a set shrinking to a point
tends to its value at that point. On passing to the limit in
\eqr{6.91}, the last term disappearing because
$|\vek{o}(R^{2})|/R^{2}\to0$ by the definition of the symbol, we
finally reach
\begin{equation}
   \bt_{d_{1}}(\bx_{0},t)=\bt_{d_{2}}(\bx_{0},t),
   \tag{6.93}\label{eq:6.93}
\end{equation}
and thus conclude that the traction at a point takes the same value for
all surfaces having the same unit normal at that point. Accordingly the
traction is determined by $\bn$, and \eqr{6.84} is proved. It follows
that the traction acting on a surface at a particular point is
identical to the traction acting on the tangent plane to the surface at
the point in question. Here, for definiteness, $\bn$ is taken to be the
exterior unit normal to $\partial P_{t1}$ and $\partial P_{t2}$ at
$\bx_{0}$. Rigorous proofs may be found in the references; the argument
above is a sketch in the sense that the error estimates have been
carried informally.

\bookfig{7}\begin{figure}[htbp]\centering
\begin{tikzpicture}[scale=1.0,
   ar/.style={-{Stealth[length=2.2mm]},thick},
   dim/.style={{Stealth[length=1.5mm]}-{Stealth[length=1.5mm]},
               semithick,gray!75},
   lead/.style={-{Stealth[length=1.4mm]},thin,gray!70},
   inner/.style={thick,orange!75!black}]
%============ panel A: tangent.  s1 : z = c xi^2, drawn with the ======
%============ vertical stretched by three.  h(R) = c R^2, so =========
%============ halving R quarters h and halves h/R. ===================
\begin{scope}[xshift=0cm]
  \fill[blue!7]
     plot[domain=-2.2:2.2,samples=60,smooth] (\x,{0.161*\x*\x})
     -- (2.2,-0.78) -- (-2.2,-0.78) -- cycle;
  \draw[guide] (-2.65,0) -- (2.65,0);
  \node[font=\small,text=gray!85,anchor=east] at (-2.70,0) {$T$};
  % the outer pillbox
  \draw[thick] (-2.2,-0.78) -- (2.2,-0.78);
  \draw[very thick,gray!55] (-2.2,-0.78) -- (-2.2,0.78);
  \draw[very thick,gray!55] ( 2.2,-0.78) -- ( 2.2,0.78);
  % the pillbox at half the radius
  \draw[inner] (-1.1,-0.195) -- (1.1,-0.195);
  \draw[inner] (-1.1,-0.195) -- (-1.1,0.195);
  \draw[inner] ( 1.1,-0.195) -- ( 1.1,0.195);
  % the surface
  \draw[very thick,blue!55!black,domain=-2.2:2.2,samples=60,smooth]
     plot (\x,{0.161*\x*\x});
  \fill (0,0) circle (1.3pt);
  \node[font=\small,anchor=north east] at (-0.04,-0.02) {$\bx_{0}$};
  \draw[ar] (0,0) -- (0,1.15) node[right,font=\small] {$\bn$};
  % labels
  \node[font=\small,text=blue!55!black,anchor=south] at (1.55,0.52)
     {$d_{1}$};
  \draw[lead] (1.55,0.50) -- (1.60,0.42);
  \node[font=\small,anchor=west] at (2.30,0.10) {$f_{1}$};
  \node[font=\small,anchor=north] at (-1.45,-0.80) {$e$};
  \node[font=\small,text=blue!55!black] at (0.02,-0.50) {$P_{t1}$};
  \draw[dim] (0,-1.10) -- (2.2,-1.10);
  \node[font=\footnotesize,fill=white,inner sep=1pt] at (1.1,-1.10)
     {$R$};
  \node[font=\footnotesize,text=orange!75!black,anchor=west]
     at (1.16,-0.30) {$R/2$};
  \node[font=\footnotesize,text=gray!80,anchor=north,align=center]
     at (0,-1.42) {tangent: halving $R$ quarters $h$,\\
     so $h(R)/R\to0$};
\end{scope}
\draw[gray!40] (3.45,-2.05)--(3.45,1.35);
%============ panel B: not tangent.  s1 : z = m xi, so h = mR ========
%============ and every box is a scaled copy of every other. =========
\begin{scope}[xshift=6.55cm]
  \fill[blue!7] (-2.2,-0.79) -- (2.2,-0.79) -- (2.2,0.79) -- cycle;
  \draw[guide] (-2.65,0) -- (2.65,0);
  \draw[thick] (-2.2,-0.79) -- (2.2,-0.79);
  \draw[very thick,gray!55] ( 2.2,-0.79) -- ( 2.2,0.79);
  \draw[very thick,gray!55] (-2.2,-0.79) -- (-2.2,-0.79);
  \draw[inner] (-1.1,-0.40) -- (1.1,-0.40);
  \draw[inner] ( 1.1,-0.40) -- ( 1.1,0.40);
  \draw[very thick,blue!55!black] (-2.2,-0.79) -- (2.2,0.79);
  \fill (0,0) circle (1.3pt);
  \node[font=\small,anchor=north east] at (-0.04,-0.02) {$\bx_{0}$};
  \draw[ar] (0,0) -- (0,1.15) node[right,font=\small] {$\bn$};
  \node[font=\small,anchor=west] at (2.30,0.10) {$f_{1}$};
  \draw[dim] (0,-1.10) -- (2.2,-1.10);
  \node[font=\footnotesize,fill=white,inner sep=1pt] at (1.1,-1.10)
     {$R$};
  \node[font=\footnotesize,text=orange!75!black,anchor=west]
     at (1.16,-0.55) {$R/2$};
  \node[font=\footnotesize,text=gray!80,anchor=north,align=center]
     at (0,-1.42) {not tangent: halving $R$ halves $h$,\\
     so $h(R)/R$ never moves};
\end{scope}
\end{tikzpicture}
\caption{Noll's pillbox in axial cross-section, and the one thing that
decides the argument: how the depth $h(R)$ of the shared face behaves
as the radius $R$ shrinks. The second surface $s_{2}$ is omitted for
clarity; it caps the same box from the other side. In each panel the
grey box has radius $R$ and the orange one radius $R/2$, and the
vertical direction is stretched by the same factor in both, so the
comparison between them is faithful. Left, $s_{1}$ tangent to $T$ at
$\bx_{0}$: the departure from $T$ is quadratic, so $h$ falls by four
when $R$ falls by two, the box flattens as it narrows, and the wall
$f_{1}$ of area $2\pi R\,h$ dies beside the cap $d_{1}$ of area
$\pi R^{2}$. Right, $s_{1}$ crossing $T$ at a fixed angle: the
departure is linear, every box is a scaled copy of every other, and the
ratio $A(f_{1})/A(d_{1})$ is the same at every radius. Tangency is not
a convenience of the drawing; it is the hypothesis, and
Remark~\ref{rem:ch6nolldepth} gives the numbers.}
\label{fig:ch6noll}
\end{figure}

%-----------------------------------------------------------------------
\subsection{Cauchy's stress theorem}\label{sec:ch6cauchythm}
%-----------------------------------------------------------------------

Sections~\ref{sec:ch6action} and~\ref{sec:ch6noll} have reduced the
contact force of \eqr{6.28} as far as the linear momentum balance alone
can take it: by \eqr{6.84} the traction at $\bx$ depends on the cut only
through its unit normal, and by \eqr{6.85} it reverses when the normal
does. The present section extracts the last consequence, and it is the
one that makes the subject computable. A function on the unit sphere
carries infinitely many values; \eqr{6.94} below says that this
particular function is the restriction to the sphere of a linear map,
hence is fixed by nine numbers. Linearity is nowhere assumed. It is
produced by \eqr{6.36}, and the mechanism is the elementary fact that a
tetrahedron has volume of one order higher in its size than its faces
have area.

\begin{theorem}[Cauchy's stress theorem]\label{thm:ch6cauchy}
Let $\bx_{0}\in\kappa_{t}$ and suppose that
\begin{itemize}\itemsep2pt
\item[\textup{(H1)}] the traction obeys Cauchy's hypothesis \eqr{6.84}
   and Cauchy's lemma \eqr{6.85};
\item[\textup{(H2)}] for each fixed unit vector $\bn$ the field
   $\bx\mapsto\bt(\bx,t;\bn)$ is continuous on a neighbourhood of
   $\bx_{0}$;
\item[\textup{(H3)}] the field $\bx\mapsto\rho(\dot\bv-\bb)$ is
   continuous on a neighbourhood of $\bx_{0}$.
\end{itemize}
Then there is a tensor $\bT(\bx_{0},t)\in\Lin(\Vt,\Vt)$, unique, such
that
\begin{equation}
   \bt(\bx,t;\bn)=\bT(\bx,t)\bn
   \tag{6.94}\label{eq:6.94}
\end{equation}
for every unit vector $\bn$. It is called the \emph{Cauchy stress
tensor}.
\end{theorem}

The proof occupies the rest of the section. Three lemmas are separated
out first, because each is used twice and none of them is about
mechanics: the first is the geometry of a tetrahedron, the second
replaces the mean value theorem by a statement that is true for
vector-valued integrands, and the third is the extension from one octant
to all of them.

%-----------------------------------------------------------------------
\begin{lemma}[geometry of the Cauchy tetrahedron]\label{lem:ch6cone}
Fix a unit vector $\bm$ with components $m_{i}=\ip{\be_{i}}\bm>0$ on the
orthonormal basis $\{\be_{i}\}$ of $\Vt$, fix $\delta>0$, and let
$\tau=\tau(\bm,\delta)$ be the tetrahedron with vertex at $\bx_{0}$
bounded by the three coordinate planes through $\bx_{0}$ and by the
plane $\ip\bm{(\bx-\bx_{0})}=\delta$. Write $s$ for the oblique face,
$s_{i}$ for the face lying in the plane normal to $\be_{i}$, and
$A(\cdot)$, $V(\cdot)$ for area and volume. Then
\begin{equation}
\begin{aligned}
   h_{i}&=\frac{\delta}{m_{i}},
   &\qquad
   A_{i}&:=A(s_{i})=\frac{\delta^{2}}{2m_{j}m_{k}},\\[2pt]
   A(s)&=\frac{\delta^{2}}{2m_{1}m_{2}m_{3}},
   &\qquad
   V(\tau)&=\frac{\delta^{3}}{6m_{1}m_{2}m_{3}},
\end{aligned}
   \tag{$\ast_{1}$}\label{eq:ch6tetgeom}
\end{equation}
where $h_{i}$ is the intercept of the oblique plane on the $\be_{i}$
axis and $\{i,j,k\}$ are distinct. Consequently
\begin{equation}
   A(s)=c_{1}\delta^{2},\qquad V(\tau)=c_{2}\delta^{3},\qquad
   \frac{V(\tau)}{A(s)}=\frac{\delta}{3},
   \tag{$\ast_{2}$}\label{eq:ch6tetscale}
\end{equation}
with $c_{1}=1/(2m_{1}m_{2}m_{3})$ and $c_{2}=c_{1}/3$ positive and
independent of $\delta$.
\end{lemma}

\begin{proof}
Take $\bx_{0}$ as origin. A point $h\be_{i}$ lies on the oblique plane
when $\ip\bm{h\be_{i}}=hm_{i}=\delta$, so $h_{i}=\delta/m_{i}$, positive
because $m_{i}>0$.

The face $s_{i}$ is the triangle with vertices $\bze$, $h_{j}\be_{j}$
and $h_{k}\be_{k}$, a right triangle with legs $h_{j}$ and $h_{k}$,
whence $A_{i}=\tfrac12h_{j}h_{k}=\delta^{2}/(2m_{j}m_{k})$. The
tetrahedron is the corner solid on the three mutually perpendicular
edges $h_{i}\be_{i}$, so $V(\tau)=\tfrac16h_{1}h_{2}h_{3}$, which is
$\eqref{eq:ch6tetgeom}_{4}$.

For $A(s)$, the oblique face is the triangle with vertices
$h_{i}\be_{i}$, whose area is half the modulus of the cross product of
two of its edges, by Proposition~\ref{prop:crosscomp}. Writing
$d=\big[\be_{1},\be_{2},\be_{3}\big]=\pm1$ for the orientation of the
triad, so that $\be_{2}\times\be_{3}=d\,\be_{1}$ and its two cyclic
partners,
\[
\begin{aligned}
   2A(s)&=\big|(h_{2}\be_{2}-h_{1}\be_{1})
     \times(h_{3}\be_{3}-h_{1}\be_{1})\big|\\
   &=\big|d\big(h_{2}h_{3}\be_{1}+h_{1}h_{3}\be_{2}
     +h_{1}h_{2}\be_{3}\big)\big|
   =\frac{\delta^{2}}{m_{1}m_{2}m_{3}}
     \big|m_{1}\be_{1}+m_{2}\be_{2}+m_{3}\be_{3}\big| ,
\end{aligned}
\]
the second line by expanding the product, the third by
$h_{j}h_{k}=\delta^{2}/(m_{j}m_{k})
=(\delta^{2}/m_{1}m_{2}m_{3})\,m_{i}$ together with $|d|=1$. The last
modulus is $|\bm|=1$, which gives $\eqref{eq:ch6tetgeom}_{3}$. Only the
orthonormality of $\{\be_{i}\}$ has been used, never its handedness:
the factor $d$ leaves a modulus untouched, and $h_{i}>0$ makes the
volume positive whichever sign $d$ carries. The lemma therefore holds
verbatim for any orthonormal triad, a fact Lemma~\ref{lem:ch6octant}
will need.

Both $c_{1}$ and $c_{2}$ are built from $\bm$ alone, and dividing the
last two members of \eqref{eq:ch6tetgeom} gives
$V(\tau)/A(s)=\delta/3$.
\end{proof}

Geometrically the last of \eqref{eq:ch6tetscale} is the statement that
$\tau$ is a cone of apex $\bx_{0}$ over the base $s$, whose apex stands
at height $\delta$; and the reason $c_{1},c_{2}$ do not move with
$\delta$ is that $\bm$ is held fixed throughout, so the tetrahedra
$\tau(\bm,\delta)$ are homothetic copies of one another about $\bx_{0}$,
scaled by $\delta$. Area scales as the square of a length and volume as
the cube, and the whole proof lives in the gap between those two powers.

%-----------------------------------------------------------------------
\begin{lemma}[averages over a shrinking surface]\label{lem:ch6avg}
Let $\bg$ be a $\Vt$-valued field, continuous at $\bx_{0}$, and let
$\{\sigma_{\delta}\}_{\delta>0}$ be surfaces of positive area whose
\emph{radius about $\bx_{0}$},
\[
   \eps(\delta):=\sup_{\bx\in\sigma_{\delta}}|\bx-\bx_{0}|
   =\text{the least }r\text{ with }
   \sigma_{\delta}\subseteq\{\bx:|\bx-\bx_{0}|\le r\} ,
\]
satisfies $\eps(\delta)\to0$ as $\delta\to0$. The supremum is written
rather than a maximum because $\sigma_{\delta}$ need not contain the
point at which the distance is largest; it is finite as soon as
$\sigma_{\delta}$ is bounded. Then
\[
   \lim_{\delta\to0}\frac{1}{A(\sigma_{\delta})}
   \int_{\sigma_{\delta}}\bg(\bx)\,\dl a=\bg(\bx_{0}).
\]
\end{lemma}

\begin{proof}
Since $A(\sigma_{\delta})^{-1}\int_{\sigma_{\delta}}\dl a=1$, the
constant $\bg(\bx_{0})$ may be taken inside the integral, and the
triangle inequality for integrals gives
\[
\begin{aligned}
   \left|\frac{1}{A(\sigma_{\delta})}\int_{\sigma_{\delta}}\bg\,\dl a
   -\bg(\bx_{0})\right|
   &=\left|\frac{1}{A(\sigma_{\delta})}\int_{\sigma_{\delta}}
     \big[\bg(\bx)-\bg(\bx_{0})\big]\dl a\right|\\
   &\leq\frac{1}{A(\sigma_{\delta})}\int_{\sigma_{\delta}}
     \big|\bg(\bx)-\bg(\bx_{0})\big|\,\dl a
   \ \leq\ \omega\big(\eps(\delta)\big),
\end{aligned}
\]
where
\[
   \omega(r)=\sup_{|\bx-\bx_{0}|\leq r}\big|\bg(\bx)-\bg(\bx_{0})\big|
\]
denotes the largest departure of $\bg$ from its value at $\bx_{0}$ over
the closed ball of radius $r$; the last inequality holds because every
$\bx\in\sigma_{\delta}$ satisfies $|\bx-\bx_{0}|\le\eps(\delta)$, so the
integrand is bounded by $\omega(\eps(\delta))$, and the average of a
constant is that constant.

It remains to show $\omega(r)\to0$ as $r\to0^{+}$, which is continuity
of $\bg$ at $\bx_{0}$ restated. Let $\eta>0$. Continuity supplies
$r_{0}>0$ with $|\bg(\bx)-\bg(\bx_{0})|<\eta$ whenever
$|\bx-\bx_{0}|<r_{0}$; taking the supremum over the ball of any radius
$r<r_{0}$ gives $\omega(r)\le\eta$. Since $\eps(\delta)\to0$, we have
$\eps(\delta)<r_{0}$ for all small $\delta$, and then
$\omega(\eps(\delta))\le\eta$.
\end{proof}

This is the step the book performs with the mean value theorem, and the
two are not interchangeable. For a scalar integrand the mean value
theorem does produce a point $\bx^{*}\in\sigma$ at which the integrand
equals its own average; for a vector-valued integrand no such point need
exist. The map $\bff(x)=(\cos2\pi x,\sin2\pi x)$ on $[0,1]$ has average
$\bze$ while $|\bff(x)|=1$ everywhere, so no $x^{*}$ obliges. Applied
componentwise the mean value theorem is of course legitimate, and then
it yields three sample points, one per component, all of which tend to
$\bx_{0}$; the conclusion is the same. Lemma~\ref{lem:ch6avg} obtains it
in one line and for the vector at once.

%-----------------------------------------------------------------------
\begin{lemma}[from one octant to all]\label{lem:ch6octant}
Assume \textup{(H1)} and \textup{(H2)}, and suppose that
\begin{equation}
   \bt(\bx_{0};\bn)=\sum_{i}\big(\ip{\bff_{i}}\bn\big)\,
   \bt(\bx_{0};\bff_{i})
   \tag{$\ddagger$}\label{eq:ch6octanthyp}
\end{equation}
holds for every orthonormal triad $\{\bff_{i}\}$ of $\Vt$ and every unit
$\bn$ with $\ip{\bff_{i}}\bn>0$ for $i=1,2,3$. Then
$\bt(\bx_{0};\bn)=\bT\bn$ for \emph{every} unit $\bn$, with $\bT$ the
single tensor $\bt(\bx_{0};\be_{i})\otimes\be_{i}$ built on any one
fixed orthonormal triad $\{\be_{i}\}$.
\end{lemma}

\begin{proof}
Suppose first that no component of $\bn$ on $\{\be_{i}\}$ vanishes, and
put $\eps_{i}=\operatorname{sgn}(\ip{\be_{i}}\bn)\in\{\pm1\}$. The
triad $\bff_{i}=\eps_{i}\be_{i}$ is orthonormal, and
$\ip{\bff_{i}}\bn=\eps_{i}\ip{\be_{i}}\bn=|\ip{\be_{i}}\bn|>0$, so
\eqref{eq:ch6octanthyp} applies to it and yields
\[
\begin{aligned}
   \bt(\bx_{0};\bn)
   &=\sum_{i}\big(\ip{\eps_{i}\be_{i}}\bn\big)\,
     \bt(\bx_{0};\eps_{i}\be_{i})
   =\sum_{i}\eps_{i}\big(\ip{\be_{i}}\bn\big)\,
     \eps_{i}\bt(\bx_{0};\be_{i})\\
   &=\sum_{i}\big(\ip{\be_{i}}\bn\big)\,\bt(\bx_{0};\be_{i})=\bT\bn ,
\end{aligned}
\]
where the middle equality used \eqr{6.85} on each coordinate face and
the last $\eps^{2}_{i}=1$.

The unit vectors with some vanishing component form the union of three
great circles, whose complement is dense in the unit sphere. Both
$\bn\mapsto\bt(\bx_{0};\bn)$ and $\bn\mapsto\bT\bn$ are continuous
there, the first by \textup{(H2)} and the second because every linear
map on a finite-dimensional space is continuous, and two continuous maps
agreeing on a dense set agree everywhere.
\end{proof}

%-----------------------------------------------------------------------
\subsection*{Proof of Theorem \ref{thm:ch6cauchy}}
%-----------------------------------------------------------------------

Fix $\bm$ with $m_{i}=\ip{\be_{i}}\bm>0$ and let $\tau=\tau(\bm,\delta)$
be the tetrahedron of Lemma~\ref{lem:ch6cone}, so that the exterior unit
normal on $\partial\tau$ is $\bm$ on $s$ and $-\be_{i}$ on $s_{i}$, the
solid lying on the positive side of each coordinate plane. Euler's first
postulate in the reduced form \eqr{6.36}, applied to $\tau$ and with the
two volume terms collected on the right, reads
\begin{equation}
   \int_{\partial\tau}\bt\,\dl a=\int_{\tau}\rho(\dot\bv-\bb)\,\dl v .
   \tag{6.95}\label{eq:6.95}
\end{equation}
By Remark~\ref{rem:symbolic} this is symbolic notation for the vector
equation
\begin{equation}
   \be_{i}\int_{\partial\tau}t_{i}\,\dl a
   =\be_{i}\int_{\tau}\rho(\dot v_{i}-b_{i})\,\dl v ,
   \tag{6.96}\label{eq:6.96}
\end{equation}
and since $\{\be_{i}\}$ is a basis, \eqr{6.96} is equivalent by
Theorem~\ref{thm:comprep} to the three scalar equations
\begin{equation}
   \int_{\partial\tau}t_{i}\,\dl a
   =\int_{\tau}\rho(\dot v_{i}-b_{i})\,\dl v ;\qquad i=1,2,3 .
   \tag{6.97}\label{eq:6.97}
\end{equation}
Every estimate below is carried out on \eqr{6.97}, so that it is an
estimate on numbers and the objection of Remark~\ref{rem:symbolic} does
not arise.

\emph{Step 1: the volume term is negligible beside the faces.} Fix once
and for all a $\delta_{0}>0$ small enough that $\tau(\bm,\delta_{0})$
lies inside the neighbourhood of \textup{(H3)}. For $0<\delta\le\delta_{0}$
one has $\tau(\bm,\delta)\subseteq\tau(\bm,\delta_{0})$, the tetrahedra
being homothetic about $\bx_{0}$, and $\tau(\bm,\delta_{0})$ is compact,
so the continuous functions $\rho(\dot v_{i}-b_{i})$ attain a finite
maximum modulus on it. Setting
\begin{equation}
   k_{i}=\max_{\bx\in\tau(\bm,\delta_{0})}
   \big|\rho(\dot v_{i}-b_{i})\big| ,
   \tag{6.99}\label{eq:6.99}
\end{equation}
where the maximum is taken over the whole of the largest tetrahedron and
denotes the largest value the modulus attains there, we obtain three
numbers $k_{1},k_{2},k_{3}$ that do not depend on $\delta$. They exist
because $\tau(\bm,\delta_{0})$ is compact and the integrand continuous,
and they dominate the corresponding maxima over every smaller
tetrahedron, since
\[
   \tau(\bm,\delta)\subseteq\tau(\bm,\delta_{0})
   \quad\Longrightarrow\quad
   \max_{\tau(\bm,\delta)}\big|\rho(\dot v_{i}-b_{i})\big|
   \ \le\ k_{i}
   \qquad (0<\delta\le\delta_{0}),
\]
a maximum over a subset never exceeding one over the set. This is the
point at which \eqr{6.99} becomes a genuine constant rather than a
function of $\delta$, and it is what makes the order estimate
\eqr{6.100} meaningful. Now \eqr{6.97} and the triangle inequality for
integrals give
\begin{equation}
   \left|\int_{\partial\tau}t_{i}\,\dl a\right|
   =\left|\int_{\tau}\rho(\dot v_{i}-b_{i})\,\dl v\right|
   \leq\int_{\tau}\big|\rho(\dot v_{i}-b_{i})\big|\,\dl v
   \leq k_{i}V(\tau).
   \tag{6.98}\label{eq:6.98}
\end{equation}
Divide by $A(s)$, which is positive, and invoke
$\eqref{eq:ch6tetscale}_{3}$:
\begin{equation}
   \left|\frac{1}{A(s)}\int_{\partial\tau}t_{i}\,\dl a\right|
   \leq k_{i}\,\frac{V(\tau)}{A(s)}=\frac{k_{i}}{3}\,\delta ,
   \qquad\text{that is}\qquad
   \frac{1}{A(s)}\int_{\partial\tau}t_{i}\,\dl a=O(\delta)
   \tag{6.100}\label{eq:6.100}
\end{equation}
in the sense of Definition~\ref{def:landau}: the ratio of the left-hand
side to $\delta$ is bounded by the constant $k_{i}/3$, uniformly in
$\delta$. By Lemma~\ref{lem:landau}\textup{(ii)} with $p=1$ and $q$ any
smaller positive number this is also $o(\delta^{q})$, but nothing so
strong is needed here: $O(\delta)$ already forces the limit, since
$k_{i}\delta/3\to0$. Reverting to symbolic notation, the three
components together give
\begin{equation}
   \lim_{\delta\to0}\frac{1}{A(s)}\int_{\partial\tau}\bt\,\dl a=\bze .
   \tag{6.101}\label{eq:6.101}
\end{equation}
Nothing has been assumed about the size of $\dot\bv$ or $\bb$; they may
be arbitrarily large, provided only that they are bounded near
$\bx_{0}$, which is what \textup{(H3)} guarantees. The conclusion is
therefore independent of the motion and of the material, and that is why
$\bT$ exists in every continuum.

\emph{Step 2: the surface term splits over the four faces.} The boundary
$\partial\tau=s\cup s_{1}\cup s_{2}\cup s_{3}$ is a union with pairwise
intersections of zero area, so the integral splits. On each piece the
normal is constant, and \eqr{6.84} permits the traction to be evaluated
at the normal of the piece:
\begin{equation}
\begin{aligned}
   \int_{\partial\tau}\bt(\bx;\bn)\,\dl a
   &=\int_{s}\bt(\bx;\bm)\,\dl a
    +\sum_{i=1}^{3}\int_{s_{i}}\bt(\bx;-\be_{i})\,\dl a\\
   &=\int_{s}\bt(\bx;\bm)\,\dl a
    -\sum_{i=1}^{3}\int_{s_{i}}\bt(\bx;\be_{i})\,\dl a ,
\end{aligned}
   \tag{6.102}\label{eq:6.102}
\end{equation}
the second line by Cauchy's lemma \eqr{6.85}. The passive argument $t$
is suppressed throughout. Applying the mean value theorem to each
scalar component, as in \eqr{6.97}, gives points at which the average is
attained,
\begin{equation}
\begin{aligned}
   \int_{s}\bt(\bx;\bm)\,\dl a&=A(s)\,\bt(\bx^{*};\bm)
     \quad\text{for some }\bx^{*}\in s,\\
   \int_{s_{i}}\bt(\bx;\be_{i})\,\dl a
     &=A_{i}\,\bt(\bx^{*}_{i};\be_{i});\qquad
     i\in\{1,2,3\}\text{ fixed},\ \bx^{*}_{i}\in s_{i},
\end{aligned}
   \tag{6.103}\label{eq:6.103}
\end{equation}
which is how the argument is usually written; strictly the sample points
carry a component index as well, and Lemma~\ref{lem:ch6avg} dispenses
with them altogether. It is the lemma that will be used at Step 4.

\emph{Step 3: the coordinate faces are the shadows of the oblique face.}
Dividing $\eqref{eq:ch6tetgeom}_{2}$ by $\eqref{eq:ch6tetgeom}_{3}$,
\[
   \frac{A_{i}}{A(s)}
   =\frac{\delta^{2}/(2m_{j}m_{k})}{\delta^{2}/(2m_{1}m_{2}m_{3})}
   =m_{i} ,
\]
so that
\begin{equation}
   A_{i}=A(s)\,\ip{\be_{i}}\bm .
   \tag{6.104}\label{eq:6.104}
\end{equation}
The same identity follows with no computation from the first equality in
\eqr{5.9}, which states that the integral of the exterior unit normal
over a closed surface vanishes. On $\partial\tau$ that normal is
piecewise constant, so the integral is $A(s)\bm-A_{j}\be_{j}$, and
\[
   A(s)\bm-A_{j}\be_{j}=\bze .
\]
Taking the inner product with $\be_{i}$ and using orthonormality returns
\eqr{6.104}. In either reading, $A_{i}$ is the area of the shadow that
$s$ casts on the $i$th coordinate plane, and the shadow of a plane
region is its area times the cosine of the angle between the planes.
This is the only place in the proof where the normal $\bm$ enters
linearly, and it is where the linearity of \eqr{6.94} comes from.

\emph{Step 4: the limit.} Divide \eqr{6.102} by $A(s)$ and substitute
\eqr{6.103} and \eqr{6.104}:
\begin{equation}
   \frac{1}{A(s)}\int_{\partial\tau}\bt(\bx;\bn)\,\dl a
   =\bt(\bx^{*};\bm)
   -\sum_{i=1}^{3}\big(\ip{\be_{i}}\bm\big)\,\bt(\bx^{*}_{i};\be_{i}).
   \tag{6.105}\label{eq:6.105}
\end{equation}
The size of the tetrahedron has disappeared from the right-hand side
except through the location of the sample points. Now let $\delta\to0$
with $\bm$ held fixed. Every point of $\tau$, hence of each of its four
faces, is a convex combination of $\bx_{0}$ and the three intercepts
$h_{i}\be_{i}$, so it lies within
\[
   \max_{i}h_{i}
   =\max_{i}\frac{\delta}{m_{i}}
   =\frac{\delta}{\min_{i}m_{i}}
\]
of $\bx_{0}$, where $\max_{i}$ and $\min_{i}$ denote the largest and the
smallest of the three numbers indexed by $i=1,2,3$, and the second
equality holds because $t\mapsto\delta/t$ is decreasing on $t>0$. Since
$\bm$ is fixed, $\min_{i}m_{i}$ is a positive constant and the bound
tends to zero with $\delta$. The radius $\eps(\delta)$ of
Lemma~\ref{lem:ch6avg} therefore satisfies
$\eps(\delta)\le\delta/\min_{i}m_{i}\to0$ for each of the four faces,
so that lemma applies to each of the four averages in \eqr{6.102}
divided by its own area, with limits $\bt(\bx_{0};\bm)$ and
$\bt(\bx_{0};\be_{i})$ by \textup{(H2)}. The left-hand side of
\eqr{6.105} tends to $\bze$ by \eqr{6.101}. Therefore
\begin{equation}
   \bt(\bx_{0};\bm)=\bT(\bx_{0})\bm ,
   \qquad\text{where}\quad
   \bT(\bx_{0})=\bt(\bx_{0};\be_{i})\otimes\be_{i},
   \tag{6.106}\label{eq:6.106}
\end{equation}
the identification of the sum
$\sum_{i}(\ip{\be_{i}}\bm)\,\bt(\bx_{0};\be_{i})$ as a tensor applied to
$\bm$ being Definition~\ref{def:tp} read backwards. Reinstating the
argument $t$,
\begin{equation}
   \bT(\bx,t)=\bt_{i}\otimes\be_{i},
   \qquad\text{where}\quad \bt_{i}=\bt(\bx,t;\be_{i})
   \tag{6.107}\label{eq:6.107}
\end{equation}
is the traction on the $i$th coordinate plane at $\bx\in\kappa_{t}$.

\emph{Step 5: all directions, and uniqueness.} Steps 1 to 4 used
nothing about $\{\be_{i}\}$ beyond its orthonormality, as the proof of
Lemma~\ref{lem:ch6cone} was careful to record, and nothing about $\bm$
beyond $m_{i}>0$, which is what puts the three intercepts on the
positive axes. What they establish is therefore exactly the hypothesis
\eqref{eq:ch6octanthyp} of Lemma~\ref{lem:ch6octant}: the identity
\[
   \bt(\bx_{0};\bn)=\sum_{i}\big(\ip{\bff_{i}}\bn\big)\,
   \bt(\bx_{0};\bff_{i})
\]
for every orthonormal triad $\{\bff_{i}\}$ and every unit $\bn$ positive
on it. That lemma removes the restriction and delivers \eqr{6.94} for
every unit vector, with the one tensor \eqr{6.106}. For uniqueness, suppose
$\bT_{1}$ and $\bT_{2}$ both satisfy \eqr{6.94}. Then
$(\bT_{1}-\bT_{2})\bn=\bze$ for every unit $\bn$, hence for every
$\bn\in\Vt$ by homogeneity, so $\bT_{1}=\bT_{2}$ by
Theorem~\ref{thm:comprep}. This completes the proof.\hfill$\blacksquare$

%-----------------------------------------------------------------------

Three features of the argument are worth recording before it is used.

The object produced in \eqr{6.106} is a tensor by construction and not
by assertion: it is a sum of tensor products of vectors of $\Vt$, hence
a linear map $\Vt\to\Vt$ in the sense of Definition~\ref{def:tp}. Its
components on the fixed basis are
\begin{equation}
   T_{ij}=\ip{\be_{i}}{\bT\be_{j}}=\ip{\be_{i}}{\bt_{j}} ,
   \tag{$\dagger$}\label{eq:ch6Tcomp}
\end{equation}
by Theorem~\ref{thm:comprep}, so that the $j$th \emph{column} of the
matrix $(T_{ij})$ is the traction on the plane whose normal is
$\be_{j}$, and conversely $\bt_{j}=T_{ij}\be_{i}$. Although \eqr{6.107}
is written on one basis, the tensor it defines does not depend on that
choice, since by Step 5 it is the unique linear map agreeing with
$\bt(\bx_{0};\cdot)$ on the whole unit sphere.

Linearity in $\bn$ was manufactured rather than assumed. The inputs were
\eqr{6.84} and \eqr{6.85} from Sections~\ref{sec:ch6action}
and~\ref{sec:ch6noll}, the boundedness supplied by \textup{(H3)}, and
the geometric identity \eqr{6.104}. Of these only \eqr{6.104} carries
any dependence on $\bm$, and it is linear in $\bm$; the linear structure
of \eqr{6.94} is that identity surviving the limit.

Finally, the theorem is a statement of structure and says nothing about
\emph{which} tensor $\bT$ is. Fixing it is the business of a
constitutive equation, and the balance laws will constrain it only
twice more: the rotational momentum balance forces
$\bT=\bT^{t}$ at \eqr{6.117}, three scalar conditions that
Proposition~\ref{prop:ch6indep} has already shown do not follow from
\eqr{6.36}, and the second law restricts it further in Chapter 8. What
remains free is the material. Equation \eqr{6.94} holds for steel and
for honey alike.

\bookfig{8}\begin{figure}[htbp]\centering
\begin{tikzpicture}[scale=1.0,
   ar/.style={-{Stealth[length=2.2mm]},thick}]
% --- axes from the vertex x0
\coordinate (O)  at (0,0);
\coordinate (E1) at (2.75,-0.95);   % e_1 direction (towards viewer-right)
\coordinate (E2) at (3.15,0.55);    % e_2 direction (to the right, back)
\coordinate (E3) at (0,2.85);       % e_3 direction (up)
\draw[guide] (O) -- (E1);  \draw[guide] (O) -- (E2);
\draw[guide] (O) -- (E3);
% --- the three intercepts
\coordinate (A1) at (1.95,-0.67);
\coordinate (A2) at (2.35,0.41);
\coordinate (A3) at (0,2.10);
% --- the tetrahedron faces
\fill[blue!8]  (O) -- (A1) -- (A3) -- cycle;   % s_2 face (normal -e_2)
\fill[blue!12] (O) -- (A2) -- (A3) -- cycle;   % s_1 face (normal -e_1)
\fill[blue!5]  (O) -- (A1) -- (A2) -- cycle;   % s_3 face (normal -e_3)
\fill[red!8]   (A1) -- (A2) -- (A3) -- cycle;  % oblique face s
\draw[thick] (O) -- (A1); \draw[thick] (O) -- (A2);
\draw[thick] (O) -- (A3);
\draw[very thick,red!65!black] (A1) -- (A2) -- (A3) -- cycle;
% --- vertex and labels
\fill (O) circle (1.3pt);
\node[font=\small,anchor=north east] at (-0.02,0.02) {$\bx_{0}$};
\node[font=\small,anchor=north] at (2.80,-1.00) {$\be_{1}$};
\node[font=\small,anchor=west]  at (3.18,0.58)  {$\be_{2}$};
\node[font=\small,anchor=east]  at (-0.05,2.85) {$\be_{3}$};
% --- the oblique normal m, from the centroid of the oblique face
\coordinate (C) at (1.43,0.61);
\draw[ar,red!65!black] (C) -- ++(1.05,0.62)
   node[right,font=\small,text=red!65!black] {$\bm$};
\fill[red!65!black] (C) circle (1.0pt);
\node[font=\footnotesize,text=red!65!black,anchor=south west]
   at (1.72,0.98) {$s$};
% --- distance delta from x0 to the oblique plane
\draw[{Stealth[length=1.8mm]}-{Stealth[length=1.8mm]},semithick,
      gray!80] (O) -- (0.72,0.42);
\node[font=\footnotesize,text=gray!85,anchor=north west] at (0.34,0.30)
   {$\delta$};
% --- face labels
\node[font=\footnotesize,anchor=north] at (1.30,0.32) {$s_{3}$};
\node[font=\footnotesize,anchor=east]  at (1.05,1.05) {$s_{2}$};
\node[font=\footnotesize,anchor=west]  at (1.22,1.35) {$s_{1}$};
% --- the tractions on two of the coordinate faces
\draw[ar,blue!55!black] (0.95,0.75) -- ++(-0.70,-0.42);
\node[font=\footnotesize,text=blue!55!black,anchor=east] at (0.22,0.34)
   {$-\bt_{1}$};
\end{tikzpicture}
\caption{Cauchy tetrahedron with vertex at $\bx_{0}$, the oblique face
$s$ at distance $\delta$ and unit normal $\bm$ with
$\ip\bm{\be_{i}}>0$. The three coordinate faces have areas
$A_{i}=A(s)\ip{\be_{i}}\bm$, since $s_{i}$ is the shadow of $s$ on the
$i$th coordinate plane. Since $\bx_{0}$ lies at distance $\delta$ from
the base $s$, the volume is $\tfrac13A(s)\delta$, so dividing the
momentum balance by $A(s)$ leaves the inertia and the body force
bounded by $k\delta/3$; they vanish in the limit while all four surface
terms survive. What is left is
$\bt(\bx_{0};\bm)=\sum_{i}(\ip{\be_{i}}\bm)\bt(\bx_{0};\be_{i})$, and
that sum is linear in $\bm$.}
\label{fig:ch6tetra}
\end{figure}

\subsection*{The equations of motion}

The foregoing construction may be repeated with a small
tetrahedron-like region, having a curved oblique surface, adjoined to
the boundary of an arbitrary part $P_{t}$ of a body with piecewise
smooth boundary $\partial P_{t}$. Such a region approaches a
tetrahedron in the limit as it collapses onto a point, the curvature of
the oblique face contributing only at higher order in $\delta$ exactly
as the tangency estimates of \eqr{6.86} did, and we again recover
\eqr{6.94} at the boundary. Every occurrence of $\bt$ in a surface
integral may therefore be replaced by $\bT\bn$, and Euler's linear
momentum balance \eqr{6.36} becomes
\begin{equation}
   \int_{P_{t}}\rho\dot\bv\,\dl v=\int_{P_{t}}\rho\bb\,\dl v
   +\int_{\partial P_{t}}\bT\bn\,\dl a .
   \tag{6.108}\label{eq:6.108}
\end{equation}

This is now of the generic vector form \eqr{6.2}, with the
identifications
\[
   \vek\psi=\bv,\qquad \br=\bb,\qquad \bA=\bT .
\]
Comparison with \eqr{6.23}, which is \eqr{6.2} localized and reduced by
conservation of mass, immediately delivers the local form of the linear
momentum balance:
\begin{equation}
   \boxed{\ \rho\dot\bv=\rho\bb+\dvg\bT ,\ }
   \tag{6.109}\label{eq:6.109}
\end{equation}
holding in a region of smoothness, or
\begin{equation}
   \rho\dot v_{i}=\rho b_{i}+T_{ij,j}
   \tag{6.110}\label{eq:6.110}
\end{equation}
in terms of Cartesian coordinates, where
$T_{ij}=\ip{\be_{i}}{\bT\be_{j}}=\ip{\be_{i}}{\bt(\bx,t;\be_{j})}
=\ip{\be_{i}}{\bt_{j}}$ by Proposition~\ref{prop:divA} and
Theorem~\ref{thm:comprep}. Thus the traction on a surface with normal
$\be_{j}$ is $\bt_{j}=(\ip{\be_{i}}{\bt_{j}})\be_{i}=T_{ij}\be_{i}$:
reading down the $j$th column of the stress matrix gives the traction
on the $j$th coordinate plane, its diagonal entry being the normal
component and the two off-diagonal entries the shears.

The relevant jump condition at a singular surface is given by
\eqr{6.11} with the same identifications:
\begin{equation}
   \jump{\rho\bv\otimes(\bv-\bu)-\bT}\bn=\bze .
   \tag{6.111}\label{eq:6.111}
\end{equation}

\begin{remark}[the second Rankine--Hugoniot condition]\label{rem:ch6rh2}
Written with the mass flux $m^{*}=\rho(\ip\bv\bn-u)$ of
Remark~\ref{rem:ch6rh1}, and using that $m^{*}$ is continuous across
$s$ by \eqr{6.15}, the condition \eqr{6.111} becomes
\[
   m^{*}\jump\bv=\jump\bT\bn=\jump\bt ,
\]
since $\big[\rho\bv\otimes(\bv-\bu)\big]\bn
=\rho\big(\ip{(\bv-\bu)}\bn\big)\bv=m^{*}\bv$ and the scalar $m^{*}$
passes through the jump bracket. So the traction on $s$ jumps by
exactly the momentum being carried across it per unit area per unit
time. If no material crosses, $m^{*}=0$ and the traction is
continuous, which is the familiar statement that a surface of contact
transmits the same force to both sides. Together with \eqr{6.15} and
the energy condition of Problem 11, this is the Rankine--Hugoniot
system for a shock, and Remark~\ref{rem:jumpsign} is now discharged.
\end{remark}

%-----------------------------------------------------------------------
\subsection{Rotational momentum balance}\label{sec:ch6rotbal}
%-----------------------------------------------------------------------

The combination of \eqr{6.94} with $\bx\times\dot\bv
=(\bx\times\bv)^{\textstyle\cdot}$, established after \eqr{6.37},
reduces the rotational momentum balance \eqr{6.37} to
\begin{equation}
   \int_{P_{t}}\rho\,(\bx\times\bv)^{\textstyle\cdot}\,\dl v
   =\int_{P_{t}}\rho\,\bx\times\bb\,\dl v
   +\int_{\partial P_{t}}\bx\times\bT\bn\,\dl a ,
   \tag{6.112}\label{eq:6.112}
\end{equation}
which is symbolic notation for the vector equation
\begin{equation}
   \be_{i}\eps_{ijk}\int_{P_{t}}\rho\,(x_{j}v_{k})^{\textstyle\cdot}
   \,\dl v
   =\be_{i}\eps_{ijk}\int_{P_{t}}\rho\,x_{j}b_{k}\,\dl v
   +\be_{i}\eps_{ijk}\int_{\partial P_{t}}x_{j}T_{kl}n_{l}\,\dl a ,
   \tag{6.113}\label{eq:6.113}
\end{equation}
where $\eps_{ijk}$ is the permutation symbol of
Definition~\ref{def:perm} and $(\ba\times\bc)_{i}=\eps_{ijk}a_{j}c_{k}$
by Proposition~\ref{prop:crosscomp}.

This is of the form \eqr{6.2}. Taking \eqr{6.24}$_{2}$ into account, we
have the identifications
\begin{equation}
   \psi_{i}=\eps_{ijk}x_{j}v_{k},\qquad
   r_{i}=\eps_{ijk}x_{j}b_{k}
   \qquad\text{and}\qquad
   A_{il}=\eps_{ijk}x_{j}T_{kl} .
   \tag{6.114}\label{eq:6.114}
\end{equation}
The third identification is read off the boundary integral: the flux
tensor is whatever $\bA$ makes $A_{il}n_{l}$ the surface integrand, and
$\eps_{ijk}x_{j}T_{kl}n_{l}=(\bx\times\bT\bn)_{i}$ is exactly that.
Inserting \eqr{6.114} into the local balance law \eqr{6.23}, that is
into
\begin{equation}
   \rho\dot\psi_{i}=\rho r_{i}+A_{il,l},
   \tag{6.115}\label{eq:6.115}
\end{equation}
and expanding both sides gives the result we are after. Take the two
sides in turn.

\emph{The left side.} By the product rule, and using $\dot x_{j}=v_{j}$
from Definition~\ref{def:va},
\[
   \dot\psi_{i}=\eps_{ijk}\big(\dot x_{j}v_{k}+x_{j}\dot v_{k}\big)
   =\eps_{ijk}v_{j}v_{k}+\eps_{ijk}x_{j}\dot v_{k}
   =\eps_{ijk}x_{j}\dot v_{k},
\]
since $\eps_{ijk}v_{j}v_{k}=0$: the array $v_{j}v_{k}$ is symmetric in
$j$ and $k$ while $\eps_{ijk}$ is antisymmetric in them, and the
contraction of a symmetric array with an antisymmetric one vanishes, by
the argument given for Problem 5 of Chapter 1.

\emph{The right side.} The divergence term is a product and must be
differentiated as one:
\[
   A_{il,l}=\eps_{ijk}\big(x_{j,l}T_{kl}+x_{j}T_{kl,l}\big)
   =\eps_{ijk}\big(\dd_{jl}T_{kl}+x_{j}T_{kl,l}\big)
   =\eps_{ijk}T_{kj}+\eps_{ijk}x_{j}T_{kl,l},
\]
using $x_{j,l}=\dd_{jl}$, which is Proposition~\ref{prop:posvec}: the
gradient of the position field is the identity. Notice that the first
of the two terms carries no factor of $\bx$ at all. That is the term
which will survive, and its appearance is entirely due to
differentiating the lever arm rather than the stress.

Substituting both expansions into \eqr{6.115} and collecting the terms
that carry $x_{j}$ on the left,
\begin{equation}
   \eps_{ijk}x_{j}\big\{\rho(\dot v_{k}-b_{k})-T_{kl,l}\big\}
   =\eps_{ijk}T_{kj}.
   \tag{6.116}\label{eq:6.116}
\end{equation}
The terms in the braces vanish by \eqr{6.110}, which says precisely
that $\rho\dot v_{k}-\rho b_{k}-T_{kl,l}=0$, leaving the algebraic
restriction $\eps_{ijk}T_{kj}=0$. Recalling Problem 5 of Chapter 1 and
Remark~\ref{rem:p5use}, we conclude that $\bT$ is a symmetric spatial
tensor,
\begin{equation}
   \boxed{\ \bT=\bT^{t}.\ }
   \tag{6.117}\label{eq:6.117}
\end{equation}

\begin{remark}[a sign in the printed \eqr{6.116}]\label{rem:ch6misprint116}
The printed text has $+T_{kl,l}$ inside the braces of \eqr{6.116}. The
sign should be negative, as the derivation above shows and as the
sentence immediately following it in the book requires: the braces are
said to vanish by \eqr{6.110}, and \eqr{6.110} reads
$\rho\dot v_{k}=\rho b_{k}+T_{kl,l}$, that is
$\rho(\dot v_{k}-b_{k})-T_{kl,l}=0$. With a plus sign the braces would
equal $2T_{kl,l}$ on solutions of \eqr{6.110} and nothing would follow.
\end{remark}

\begin{remark}[what rotational momentum balance did and did not
give]\label{rem:ch6symmeans}
This is the one balance law of the chapter that produces no
differential equation. Every term involving $\bx$ cancelled, and what
remained was an algebraic condition on $\bT$ at each point. The reason
is the following.

The moment balance is the force balance multiplied by a lever arm and
integrated. Once the force balance holds pointwise, the moment of the
forces about any point is automatically balanced by the moment of the
inertia, \emph{except} for the contribution of the lever arm's own
variation, which is the term $\eps_{ijk}T_{kj}$ produced by
$x_{j,l}=\dd_{jl}$. So the moment balance carries exactly one piece of
information beyond the force balance, and that piece is
$\eps_{ijk}T_{kj}=0$.

This is where the non-polar assumption of
Remark~\ref{rem:ch6nonpolar} is paid for. Had couple tractions been
admitted, \eqr{6.112} would carry an extra surface integral
$\int_{\partial P_{t}}\bmm\,\dl a$, its divergence would appear on the
right of \eqr{6.116}, and the conclusion would be not
$\eps_{ijk}T_{kj}=0$ but $\eps_{ijk}T_{kj}=-(\text{couple stress
divergence})_{i}$: the skew part of $\bT$ would be determined by the
couple stresses instead of vanishing. Symmetry of the Cauchy stress is
therefore not a theorem of continuum mechanics as such; it is a theorem
of \emph{non-polar} continuum mechanics, and it is the exact
mathematical content of that hypothesis. Chapters 8 to 10 use it
constantly, for instance in the step $\bT\cdot\bL=\bT\cdot\bD$ that
opens every discussion of stress power.
\end{remark}

The jump condition \eqr{6.11} at a singular surface is, with the
identifications \eqr{6.114},
\begin{equation}
\begin{aligned}
   \bze&=\be_{i}\jump{\rho\psi_{i}(v_{l}-u_{l})-A_{il}n_{l}}\\
   &=\be_{i}\eps_{ijk}
     \jump{\rho x_{j}v_{k}(v_{l}-u_{l})-x_{j}T_{kl}n_{l}} .
\end{aligned}
   \tag{6.118}\label{eq:6.118}
\end{equation}
If the position $\bx$ of a material point is continuous across the
singular surface, that is if the motion of the body introduces no gaps
or fissures in the material, then $\jump{x_{j}}=0$ and the $x_{j}$ may
be taken outside the jump brackets, reducing \eqr{6.118} to
\begin{equation}
   \bx\times\big\{\jump{\rho\bv\otimes(\bv-\bu)-\bT}\bn\big\}=\bze ,
   \tag{6.119}\label{eq:6.119}
\end{equation}
in which the term in braces vanishes by \eqr{6.111}. In this case the
jump condition is identically satisfied and adds nothing. In the event
that position is discontinuous, however, \eqr{6.118} yields a
non-trivial condition, which is one way of seeing that a crack or a
cavitating interface is not covered by the smooth theory.

%-----------------------------------------------------------------------
\subsection{Decomposition of the stress into normal and shear
components}\label{sec:ch6normshear}
%-----------------------------------------------------------------------

Consider the traction acting on the tangent plane to a surface
$s\subset\kappa_{t}$ at the point $\bx$ (Figure~\ref{fig:ch6shear}).
By Noll's theorem this is the same as the traction on $s$ itself, and
by Cauchy's theorem it is $\bt=\bT\bn$, where $\bn$ is a unit normal to
the tangent plane. Evidently
\begin{equation}
   \bt=\I\bt=(\bn\otimes\bn)\bt+\Proj\bn\bt ,
   \tag{6.120}\label{eq:6.120}
\end{equation}
where $\I$ is the spatial identity and
\begin{equation}
   \Proj\bn=\I-\bn\otimes\bn
   \tag{6.121}\label{eq:6.121}
\end{equation}
is the projection onto the tangent plane, Problem 10 of Chapter 1. All
\eqr{6.120} does is add and subtract $(\bn\otimes\bn)\bt$; its content
is that the two pieces so produced are respectively parallel and
perpendicular to $\bn$, which is what Problem 10 established.

We write this in the form
\begin{equation}
   \bt=\sigma\bn+\tau\bs ,
   \tag{6.122}\label{eq:6.122}
\end{equation}
where
\begin{equation}
   \sigma=\ip\bn\bt
   \qquad\text{and}\qquad
   \tau\bs=\Proj\bn\bt ,
   \tag{6.123}\label{eq:6.123}
\end{equation}
with $|\bs|=1$ and
\begin{equation}
   \tau=\big|\Proj\bn\bt\big| .
   \tag{6.124}\label{eq:6.124}
\end{equation}
The first equality of \eqr{6.123} is $(\bn\otimes\bn)\bt
=(\ip\bn\bt)\bn$ read off Definition~\ref{def:tp}, and the second
merely names the magnitude and direction of the tangential part; $\bs$
is well defined whenever $\tau\neq0$ and is arbitrary when $\tau=0$,
in which case the term it multiplies vanishes anyway.

We refer to $\sigma$ and $\tau$ respectively as the \emph{normal} and
\emph{shear} components of the traction. These are given in terms of
the stress tensor by
\begin{equation}
   \sigma=\ip\bn{\bT\bn}
   \qquad\text{and}\qquad
   \tau=\ip\bs{\bT\bn} .
   \tag{6.125}\label{eq:6.125}
\end{equation}
The first is \eqr{6.123}$_{1}$ with $\bt=\bT\bn$. For the second, take
the inner product of \eqr{6.122} with $\bs$ and use the two facts that
define $\bs$, namely $\ip\bs\bn=0$ and $\ip\bs\bs=1$:
\[
   \ip\bs\bt
   =\ip\bs{\big(\sigma\bn+\tau\bs\big)}
   =\sigma\underbrace{\ip\bs\bn}_{=\,0}
    +\tau\underbrace{\ip\bs\bs}_{=\,1}
   =\tau ,
\]
while the left side is $\ip\bs{\bT\bn}$ by \eqr{6.94}.

The normal part can be positive, negative or zero, according as the
material on the far side of the plane is pulling, pushing, or doing
neither. The shear part is non-negative, being a magnitude by
\eqr{6.124}; its sign has been absorbed into the direction $\bs$.

We observe that
\[
   \ip\bs{\bT\bn}=\ip\bn{\bT^{t}\bs}=\ip\bn{\bT\bs},
\]
the first equality by the definition of the transpose,
Definition~\ref{def:transpose}, and the second by the symmetry
\eqr{6.117}. This is called the \emph{complementary property of shear}:
the shear traction on the plane with normal $\bn$, resolved along
$\bs$, equals the shear traction on the plane with normal $\bs$,
resolved along $\bn$. It is the symmetry of the stress made visible,
and it is what forces shear stresses on perpendicular faces of a small
block to occur in balancing pairs.

It is also clear that the shear traction is determined entirely by
$\bT$ and $\bn$, since eliminating $\sigma$ between \eqr{6.122} and
\eqr{6.125}$_{1}$ gives
\begin{equation}
   \tau\bs=\bT\bn-\big(\ip\bn{\bT\bn}\big)\bn .
   \tag{6.126}\label{eq:6.126}
\end{equation}
If $\tau$ vanishes on the tangent plane at the point in question, then
\begin{equation}
   \bT\bn=\sigma\bn ,
   \tag{6.127}\label{eq:6.127}
\end{equation}
and $\bn$ is a principal vector of $\bT$, in the sense of
Definition~\ref{def:princ}, associated with the principal stress
$\sigma$. Because of the symmetry of $\bT$ there are, by the spectral
theorem, Theorem~\ref{thm:spectral}, three real-valued principal
stresses and an associated orthonormal triad of principal vectors.
Thus there exist three mutually orthogonal planes intersecting at each
$\bx\in\kappa_{t}$ on which the shear tractions vanish. These are
uniquely determined if all the $\sigma$'s are distinct, and Chapter 3's
discussion of repeated principal values applies verbatim if they are
not.

\bookfig{11}\begin{figure}[htbp]\centering
\begin{tikzpicture}[scale=1.05,
   ar/.style={-{Stealth[length=2.4mm]},thick}]
% the tangent plane, drawn as a parallelogram
\draw[thick,fill=blue!6]
   (-2.6,-0.75) -- (1.4,-0.75) -- (2.6,0.75) -- (-1.4,0.75) -- cycle;
\node[font=\footnotesize,text=gray!85,anchor=west] at (1.55,-0.98)
   {tangent plane at $\bx$};
% the surface s, a curve through x cutting the plane
\draw[very thick,red!65!black]
   (-2.15,-0.05) .. controls (-0.75,-0.55) and (0.75,0.55) .. (2.15,0.05);
\node[font=\small,text=red!65!black,anchor=west] at (2.18,0.10) {$s$};
\coordinate (X) at (0,0);
\fill (X) circle (1.3pt);
\node[font=\small,anchor=north east] at (-0.04,-0.02) {$\bx$};
% the normal
\draw[ar] (X) -- (0,2.15) node[right,font=\small] {$\bn$};
% the traction, tilted
\draw[ar,blue!55!black] (X) -- (1.55,1.70)
   node[right,font=\small,text=blue!55!black] {$\bt=\bT\bn$};
% its normal component
\draw[ar,red!65!black] (X) -- (0,1.70)
   node[left,font=\small,text=red!65!black] {$\sigma\bn$};
% its shear component, completing the rectangle
\draw[guide] (0,1.70) -- (1.55,1.70);
\draw[ar,red!65!black] (X) -- (1.55,0)
   node[below right=-1pt,font=\small,text=red!65!black] {$\tau\bs$};
\draw[guide] (1.55,0) -- (1.55,1.70);
% right angle at the origin between n and s
\draw[semithick] (0.24,0) -- (0.24,0.24) -- (0,0.24);
\end{tikzpicture}
\caption{Resolution of the traction into normal and shear components.
By Noll's theorem the traction on $s$ at $\bx$ is the traction on the
tangent plane there, and by Cauchy's theorem it is $\bT\bn$. Splitting
it with the projection \eqr{6.121} gives $\sigma\bn$ along the normal
and $\tau\bs$ in the plane. The normal part $\sigma=\ip\bn{\bT\bn}$
may have either sign; the shear part $\tau=|\Proj\bn\bT\bn|$ cannot be
negative, its sign having been absorbed into $\bs$. Planes on which
$\tau=0$ are exactly the principal planes of $\bT$, and by
Theorem~\ref{thm:spectral} three mutually perpendicular ones exist at
every point.}
\label{fig:ch6shear}
\end{figure}

\subsection*{Examples of states of stress}

\subsubsection*{1. Pure pressure}

In a state of pure pressure the traction on every plane containing a
point $\bx$ is perpendicular to that plane. Thus
$\bt(\bx,t;\bn)=-p(\bx,t)\bn$ for all unit vectors $\bn$, where $p$ is
the pressure. Suppressing the passive arguments $\bx$ and $t$, we have
$\bT\bn=-p\bn$. On multiplying by an arbitrary scalar $\nu$ we conclude
that
\[
   \bT\bv=-p\bv=(-p\I)\bv ,
   \qquad\text{in which}\quad \bv=\nu\bn
\]
is an arbitrary vector: every non-zero vector is a positive multiple of
a unit vector, and the case $\bv=\bze$ is covered because both sides
are then zero by linearity. Two linear maps agreeing on every vector
are the same map, by Definition~\ref{def:tensor}, so the stress is
\begin{equation}
   \bT=-p\I ,
   \tag{6.128}\label{eq:6.128}
\end{equation}
which is clearly symmetric, as \eqr{6.117} requires.

Alternatively, in components,
$T_{ij}=\ip{\be_{i}}{\bT\be_{j}}=\ip{\be_{i}}{\bt(\be_{j})}
=-p\,\ip{\be_{i}}{\be_{j}}=-p\dd_{ij}$, which is the component form of
the same result. We may also derive it from \eqr{6.107}:
$\bT=\bt(\bx,t;\be_{i})\otimes\be_{i}=-p\,\be_{i}\otimes\be_{i}
=-p\I$, the last step by the resolution of the identity
$\be_{i}\otimes\be_{i}=\I$ of Theorem~\ref{thm:comprep}.

Note the sign convention. A positive $p$ makes $\bt$ point \emph{into}
the region whose exterior normal is $\bn$, that is, the surroundings
push. This is why a fluid at rest carries $\bT=-p\I$ with $p>0$, and
why the hydrostatic example of Chapter 10 finds
$p=p_{a}+\rho g(h-x_{3})$ rather than its negative.

\subsubsection*{2. Uniaxial tension and compression}

In this example the traction is of the form
$\bt(\bx,t;\bn)=\alpha(\bx,t;\bn)\be$ for all unit vectors $\bn$, where
$\alpha$ is a scalar and $\be$ is a fixed unit vector. Whatever plane
one cuts, the force transmitted is along the one direction $\be$; only
its magnitude depends on the cut.

Thus $\bT(\bx,t)\bn=\alpha(\bx,t;\bn)\be$. Let $\alpha^{*}(\bx,t;\cdot)$
be an extension of $\alpha(\bx,t;\cdot)$ from the set of unit vectors
to the set of all vectors, so that $\alpha^{*}(\bx,t;\bv)
=\alpha(\bx,t;\bv)$ for all $\bv$ with $|\bv|=1$. Then
$\bT(\bx,t)\bv=\alpha^{*}(\bx,t;\bv)\be$ for all $\bv$. Clearly
$\bT(\bx,t)\bv$ is a linear function of $\bv$, and $\be$ is a fixed
non-zero vector, so $\alpha^{*}(\bx,t;\bv)$ is a linear function of
$\bv$ as well. It follows from the Riesz representation theorem,
Theorem~\ref{thm:riesz}, that there is a vector function $\ba(\bx,t)$
with
\[
   \alpha^{*}(\bx,t;\bv)=\ip{\ba(\bx,t)}\bv .
\]

The extension is not a technicality to be skipped. Theorem
\ref{thm:riesz} represents a linear functional on a vector space, and
the set of unit vectors is not a vector space: an arbitrary linear
combination of unit vectors is not a unit vector, and the zero vector
is not a unit vector at all. Without extending $\alpha$ to a domain
that is closed under addition and scalar multiplication, the word
``linear'' has no meaning and the theorem has nothing to act on. The
extension is available because $\bT\bv$ is already defined for every
$\bv$ and $\be\neq\bze$, so one may simply define $\alpha^{*}$ by
$\alpha^{*}(\bv)\be=\bT\bv$.

We thus have $\bT\bv=(\ip\ba\bv)\be=(\be\otimes\ba)\bv$ for all $\bv$,
and hence $\bT=\be\otimes\ba$. The symmetry condition \eqr{6.117}
requires $\be\otimes\ba=\ba\otimes\be$, and this in turn yields
\begin{equation}
   \ba=\ba(\ip\be\be)=(\ba\otimes\be)\be=(\be\otimes\ba)\be=\sigma\be ,
   \qquad\text{where}\quad \sigma=\ip\ba\be .
   \tag{6.129}\label{eq:6.129}
\end{equation}
The four equalities are, in order: the first inserts $\ip\be\be=1$, the
second recognises $\ba(\ip\be\be)$ as $(\ba\otimes\be)\be$ by
Definition~\ref{def:tp}, the third applies the symmetry just imposed,
and the fourth is $(\be\otimes\ba)\be=\be(\ip\ba\be)$, again by
Definition~\ref{def:tp}. So symmetry forces $\ba$ to be parallel to
$\be$. Finally,
\begin{equation}
   \bT(\bx,t)=\sigma(\bx,t)\,\be\otimes\be ,
   \tag{6.130}\label{eq:6.130}
\end{equation}
where $\sigma$ is the uniaxial tension or compression according as it
is positive or negative.

\begin{remark}[what symmetry cost the uniaxial
state]\label{rem:ch6uniaxread}
Before symmetry was imposed, the state $\bT=\be\otimes\ba$ had three
free parameters, the components of $\ba$; after it, one, namely
$\sigma$. Two constraints were spent, which is right: the skew part of
a three-dimensional tensor has three components, and
$\Skw(\be\otimes\ba)=\tfrac12(\be\otimes\ba-\ba\otimes\be)$ vanishes
exactly when $\ba\parallel\be$, one condition of codimension two on the
direction of $\ba$.

The state \eqr{6.130} is exactly the reactive stress found for an
inextensible fibre family in \eqr{10.28}, with $\be$ the fibre
direction, and its traction on a plane with normal $\bn$ is
$\sigma(\ip\be\bn)\be$: proportional to the cosine of the angle between
the cut and the fibre, and always directed along the fibre. Cut
parallel to $\be$ and nothing at all is transmitted, which is what one
would expect of a bundle of strings.
\end{remark}

%-----------------------------------------------------------------------
\subsection{Referential equations of motion}\label{sec:ch6refeqs}
%-----------------------------------------------------------------------

The reader will have noticed that the local equations and their jump
conditions have been obtained entirely in the spatial description. In
particular we have not used the referential expressions
\eqr{6.25}$_{2}$, \eqr{6.27}$_{2}$ or \eqr{6.29} in deducing local
equations from the linear momentum balance \eqr{6.34}. We do so here
for the local equations and leave the associated jump conditions to
Problem 1.

Combining \eqr{6.38} with \eqr{6.29} gives the linear momentum balance
for an arbitrary part $P\subset\kappa$:
\begin{equation}
   \int_{\partial P}\bp\,\dl A=\int_{P}\rho_{\kappa}(\dot\bv-\bb)\,\dl V,
   \tag{6.131}\label{eq:6.131}
\end{equation}
where $\bp$ is the Piola traction.

% reader's note: the five substitutions hidden in the word
% ``combining''
\rnote{ch6:6.131-combining}

We could use this in a version of
Noll's theorem, run with regions in $\kappa$ rather than in
$\kappa_{t}$, to conclude that $\bp$ depends on $\partial P$ through
its local unit normal $\bN$, and then adapt Cauchy's theorem, with a
tetrahedron in $\kappa$, to conclude that
\begin{equation}
   \bp(\bX,t;\bN)=\bP(\bX,t)\bN ,
   \tag{6.132}\label{eq:6.132}
\end{equation}
where $\bP$, the \emph{Piola stress}, is a two-point tensor field, with
Cartesian-component representation
\begin{equation}
   \bP=P_{iA}\,\be_{i}\otimes\hbE_{A} .
   \tag{6.133}\label{eq:6.133}
\end{equation}
Neither argument needs any change: both proceed by shrinking a region
and comparing a surface measure with a volume measure, and both are
indifferent to which configuration the region lives in. Nothing in them
used the motion.

% reader's note: the transplant, carried out rather than asserted
\rnote{ch6:6.132-transplant}

We can combine this with the referential form of the divergence
theorem, for smooth and hence continuous fields, to obtain
\begin{equation}
   \int_{P}\rho_{\kappa}(\dot\bv-\bb)\,\dl V
   =\int_{\partial P}\bP\bN\,\dl A=\int_{P}\Dvg\bP\,\dl V ,
   \tag{6.134}\label{eq:6.134}
\end{equation}
where, in terms of Cartesians,
\begin{equation}
   \Dvg\bP=P_{iA,A}\,\be_{i},
   \tag{6.135}\label{eq:6.135}
\end{equation}
the comma now denoting differentiation with respect to $X_{A}$, and
then localize by Theorem~\ref{thm:localization} to conclude that
\begin{equation}
   \boxed{\ \rho_{\kappa}\dot\bv=\rho_{\kappa}\bb+\Dvg\bP ,\ }
   \tag{6.136}\label{eq:6.136}
\end{equation}
that is
\begin{equation}
   \rho_{\kappa}\dot v_{i}=\rho_{\kappa}b_{i}+P_{iA,A},
   \tag{6.137}\label{eq:6.137}
\end{equation}
at all points $\bX\in\kappa$ that do not belong to a singular surface.

% reader's note: the referential divergence theorem, and the
% localization, in full
\rnote{ch6:6.134-divergence}

\begin{remark}[why anyone bothers with \eqr{6.136}]\label{rem:ch6whyref}
Equation \eqr{6.136} looks like \eqr{6.109} with hats on, and one may
reasonably ask what has been gained by rewriting a correct equation.
Two things, and both are practical.

\emph{The density is known.} In \eqr{6.136} the coefficient
$\rho_{\kappa}(\bX)$ is a specified function, part of the description
of the body, by \eqr{6.18}. In \eqr{6.109} the spatial density
$\rho(\bx,t)$ is an unknown field which must satisfy the partial
differential equation \eqr{6.14}, involving the velocity, and is not
known in advance. So the referential formulation solves conservation of
mass for free and reduces the number of unknown fields by one.

\emph{The acceleration is a partial derivative.} In the referential
description $\dot\bv=\partial\hat\bv(\bX,t)/\partial t$ is simply the
partial time derivative at fixed $\bX$: no convective term, since
material points do not move in $\kappa$. In the spatial description
$\dot\bv=\bv'+(\grad\bv)\bv$ by \eqr{2.21}, and the second term is
quadratic in the unknown. The nonlinearity has not disappeared, it has
moved into the constitutive relation between $\bP$ and the
deformation, which is where a solid mechanician wants it.

The price is that the domain $\kappa$ must be known, which is why the
referential form is the natural one for solids and the spatial form the
natural one for fluids.
\end{remark}

\subsection*{The relation between the two stresses}

The connection \eqr{2.59} between the referential and spatial mass
densities, combined with \eqr{6.109}, suggests that if $\bP$ is such
that
\begin{equation}
   J\dvg\bT=\Dvg\bP ,
   \tag{6.138}\label{eq:6.138}
\end{equation}
then the referential and spatial forms of the balance of linear
momentum are equivalent. Multiplying \eqr{6.109} by $J>0$ and
substituting $J\rho=\rho_{\kappa}$ from \eqr{2.59} term by term,
\[
\begin{aligned}
   \rho\dot\bv&=\rho\bb+\dvg\bT
   &&\Big|\times J\\
   J\rho\,\dot\bv&=J\rho\,\bb+J\dvg\bT\\
   \rho_{\kappa}\dot\bv&=\rho_{\kappa}\bb+J\dvg\bT ,
\end{aligned}
\]
and comparing the last line with the referential balance \eqr{6.136},
\[
   \rho_{\kappa}\dot\bv=\rho_{\kappa}\bb+\Dvg\bP ,
\]
the two agree for every motion precisely when their third terms do,
that is, precisely when $J\dvg\bT=\Dvg\bP$. So \eqr{6.138} is not an
extra assumption about $\bP$ but the exact condition under which the
two descriptions carry the same content.

% reader's note: why the factor is $J$ and not something else
\rnote{ch6:6.138-whyJ}

To determine the requisite expression for $\bP$ we combine \eqr{6.94}
with \eqr{2.50} and the Piola--Nanson formula \eqr{2.54}, obtaining
\begin{equation}
\begin{aligned}
   \int_{P}J\dvg\bT\,\dl V
   &=\int_{P_{t}}\dvg\bT\,\dl v
    =\int_{\partial P_{t}}\bT\bn\,\dl a
    =\int_{\partial P}\bT\bF^{*}\bN\,\dl A\\
   &=\int_{P}\Dvg\big(\bT\bF^{*}\big)\,\dl V .
\end{aligned}
   \tag{6.139}\label{eq:6.139}
\end{equation}
Read the chain from left to right. The first equality is the change of
variables $\dl v=J\,\dl V$ of \eqr{2.50}. The second is the divergence
theorem \eqr{5.7}$_{2}$ in $\kappa_{t}$. The third substitutes
$\bn\,\dl a=\bF^{*}\bN\,\dl A$, Piola--Nanson \eqr{2.54}, and moves the
constant-in-the-integration tensor $\bT$ through:
$\bT(\bF^{*}\bN)=(\bT\bF^{*})\bN$ is associativity of composition. The
fourth is the divergence theorem again, this time in $\kappa$. Note
that $\bT\bF^{*}$ is a two-point tensor, $\bF^{*}\colon\hVt\to\Vt$
followed by $\bT\colon\Vt\to\Vt$, so $\Dvg$ is the right operator for
it, in the sense of \eqr{6.135}.

Collecting the first and last integrals of \eqr{6.139} and localizing,
which is legitimate since $P$ is arbitrary and the integrands are
continuous, gives \eqr{6.138} together with
\begin{equation}
   \bP=\bT\bF^{*},
   \tag{6.140}\label{eq:6.140}
\end{equation}
or, in terms of components,
\begin{equation}
   P_{iA}=T_{ij}F^{*}_{jA}.
   \tag{6.141}\label{eq:6.141}
\end{equation}
Problem 8 establishes \eqr{6.138} directly, by differentiating
\eqr{6.141} and invoking the Piola identity \eqr{5.11}; that
calculation is the one that shows where $\Dvg\bF^{*}=\bze$ is needed.

% reader's note: the chain \eqr{6.139} in components, and the one
% step that is not bookkeeping
\rnote{ch6:6.139-chain}

Indeed, the two expressions \eqr{6.28} and \eqr{6.29} for the contact
force imply that
\begin{equation}
   \int_{\partial P}\bp\,\dl A=\int_{\partial P_{t}}\bt\,\dl a
   =\int_{\partial P_{t}}\bT\bn\,\dl a
   =\int_{\partial P}\bT\bF^{*}\bN\,\dl A ,
   \tag{6.142}\label{eq:6.142}
\end{equation}
which is identically satisfied by \eqr{6.132} and \eqr{6.140}. Since
$\partial P$ is arbitrary this also gives the pointwise relation
$\bp\,\dl A=\bt\,\dl a$ promised in Remark~\ref{rem:ch6perunit}: the
Piola traction is the Cauchy traction rescaled by the ratio of the two
area elements, and it is larger or smaller than $\bt$ according as the
surface has contracted or expanded.

% reader's note: $\bp\,\dl A=\bt\,\dl a$, with the arithmetic and
% one worked case
\rnote{ch6:6.142-tractions}

The symmetry of the Cauchy stress imposes a restriction on the Piola
stress and the deformation jointly. This follows on combining
\eqr{6.140} with \eqr{2.55}, $\bF^{*}=J\bF^{-t}$:
\[
   \bP\bF^{t}=\bT\bF^{*}\bF^{t}=\bT\big(J\bF^{-t}\big)\bF^{t}
   =J\bT\big(\bF^{t}\big)^{-1}\bF^{t}=J\bT ,
\]
which is symmetric by \eqr{6.117}. A tensor equal to its own transpose
gives $\bP\bF^{t}=(\bP\bF^{t})^{t}=\bF\bP^{t}$, so
\begin{equation}
   \bP\bF^{t}=\bF\bP^{t},
   \tag{6.143}\label{eq:6.143}
\end{equation}
that is $P_{iA}F_{jA}=F_{iA}P_{jA}$.

% reader's note: \eqr{6.143} in components, and how many conditions
% it really is
\rnote{ch6:6.143-components}

\begin{remark}[why $\bP$ is not symmetric, and what
\eqr{6.143} means]\label{rem:ch6Pnotsym}
Asking whether $\bP$ is symmetric is asking a meaningless question, and
Remark~\ref{rem:nosym} of Chapter 2 explains why: $\bP$ maps $\hVt$ to
$\Vt$, so $\bP^{t}$ maps $\Vt$ to $\hVt$, and two maps between
different pairs of spaces cannot be compared. Symmetry is defined only
for a tensor mapping a space to itself.

What \eqr{6.143} says instead is that the particular \emph{spatial}
tensor $\bP\bF^{t}$, which does map $\Vt$ to $\Vt$, is symmetric. It
is a restriction on $\bP$ \emph{and} $\bF$ together, three scalar
conditions rather than the three that $\bT=\bT^{t}$ imposes, and it is
not a restriction on $\bP$ alone: given any $\bF$, the nine components
of $\bP$ are cut down to six by \eqr{6.143}, but which six depends on
$\bF$. This is the reason that the purely referential stress of
\eqr{6.144} is introduced: it restores an honest symmetry.
\end{remark}

The purely referential \emph{Piola--Kirchhoff stress} $\bS$, defined by
\begin{equation}
   \bP=\bF\bS ,
   \tag{6.144}\label{eq:6.144}
\end{equation}
is often useful in the formulation of constitutive equations,
especially for solids. Its relationship to the Cauchy stress follows
from \eqr{6.140} and \eqr{2.55}: substituting \eqr{6.144} into
$\bP\bF^{t}=J\bT$, computed just above,
\begin{equation}
   J\bT=\bF\bS\bF^{t}.
   \tag{6.145}\label{eq:6.145}
\end{equation}
Note that $\bS=\bF^{-1}\bP$ maps $\hVt$ to $\hVt$, since
$\bP\colon\hVt\to\Vt$ and $\bF^{-1}\colon\Vt\to\hVt$, so asking whether
it is symmetric is a legitimate question. It is, and the proof is three
lines. Transpose \eqr{6.145}, using $(\bA\bB\bC)^{t}
=\bC^{t}\bB^{t}\bA^{t}$ and $J^{t}=J$ for the scalar:
\[
   J\bT^{t}=\big(\bF\bS\bF^{t}\big)^{t}
   =\big(\bF^{t}\big)^{t}\bS^{t}\bF^{t}=\bF\bS^{t}\bF^{t} .
\]
The Cauchy stress is symmetric by \eqr{6.117}, so the left side is
$J\bT=\bF\bS\bF^{t}$, again by \eqr{6.145}. Equating,
\[
   \bF\bS^{t}\bF^{t}=\bF\bS\bF^{t} .
\]
Both $\bF$ and $\bF^{t}$ are invertible, $J=\det\bF\neq0$, so
multiplying on the left by $\bF^{-1}$ and on the right by $\bF^{-t}$
removes them:
\[
   \bF^{-1}\bF\,\bS^{t}\,\bF^{t}\bF^{-t}
   =\bF^{-1}\bF\,\bS\,\bF^{t}\bF^{-t}
   \qquad\Longrightarrow\qquad
   \bS^{t}=\bS .
\]
So the symmetry of $\bT$, which came from Euler's second postulate,
transfers to $\bS$ but not to $\bP$; the difference is that the
congruence $\bA\mapsto\bF\bA\bF^{t}$ preserves symmetry while the
one-sided product $\bA\mapsto\bF\bA$ does not. In terms of
components,
\begin{equation}
   P_{iA}=F_{iB}S_{BA}
   \qquad\text{and}\qquad
   JT_{ij}=F_{iA}F_{jB}S_{AB}.
   \tag{6.146}\label{eq:6.146}
\end{equation}

% reader's note: the components of \eqr{6.146}, and the six
% conversions on one page
\rnote{ch6:6.146-table}

\begin{remark}[three stresses, one state]\label{rem:ch6threestress}
The three tensors $\bT$, $\bP$ and $\bS$ describe the same physical
state and are related by invertible transformations, so no information
is gained or lost in passing between them. What differs is what each is
convenient for.

\emph{$\bT$ is what is measured.} Force per unit \emph{current} area is
what a gauge on the deformed body reads, and $\bT$ is the only one of
the three that is a spatial tensor and symmetric.

\emph{$\bP$ is what balances.} It appears in \eqr{6.136} with the known
coefficient $\rho_{\kappa}$ and a divergence taken over a fixed domain.
It is a two-point tensor and satisfies no symmetry of its own.

\emph{$\bS$ is what a constitutive equation gives.} It is referential
and symmetric, so it has six components and lives in the same space as
$\bC$ and $\bE$. Chapter 8 will write constitutive equations as maps
from $\bC$ to $\bS$ for exactly this reason, and \eqr{6.160} below
shows that $\bS$ is the stress power conjugate to $\bE$.

A useful mnemonic for the factors: $\bP=\bT\bF^{*}$ converts area
elements, since $\bF^{*}$ is the tensor that carries $\bn\,\dl a$ to
$\bN\,\dl A$ by Remark~\ref{rem:whycof}, and $\bS=\bF^{-1}\bP$ then
pulls the remaining spatial leg back to $\kappa$. Every factor of
$\bF$, $\bF^{*}$ or $J$ in this subsection is there to convert a
measure, and nothing else.
\end{remark}

%-----------------------------------------------------------------------
\subsection{Mechanical energy balance}\label{sec:ch6mechen}
%-----------------------------------------------------------------------

In particle mechanics the power generated by the force $\vek{f}$ acting
on a particle is balanced by the rate of change of the kinetic energy
of the particle. This follows from Newton's law $\vek{f}=m\dot\bv$,
where $m$ and $\bv$ are the particle mass and velocity: the power is
\[
   \mathcal{P}=\ip{\vek{f}}\bv=m\,\ip\bv{\dot\bv}
   =\tfrac12m\big(\ip\bv{\dot\bv}+\ip{\dot\bv}\bv\big)
   =\tfrac12m\big(\ip\bv\bv\big)^{\textstyle\cdot},
\]
the middle step by the symmetry of the inner product, the last by the
product rule read backwards. Thus $\mathcal{P}=\dot{\mathcal{K}}$,
where $\mathcal{K}=\tfrac12m|\bv|^{2}$ is the kinetic energy. We derive
an analogous result for continua, and the interest lies entirely in the
extra term that appears.

To achieve this we form the inner product of \eqr{6.109} with the
material velocity $\bv$, obtaining
\begin{equation}
\begin{aligned}
   \tfrac12\rho\big(|\bv|^{2}\big)^{\textstyle\cdot}
   &=\rho\,\ip\bb\bv+\ip{\dvg\bT}\bv\\
   &=\rho\,\ip\bb\bv+\dvg\big(\bT^{t}\bv\big)-\bT\cdot\bL ,
\end{aligned}
   \tag{6.147}\label{eq:6.147}
\end{equation}
where $\bL$ is the spatial velocity gradient defined by \eqr{4.1}, the
tensor inner product is that of Definition~\ref{def:tip}, and use has
been made of the identity of Proposition~\ref{prop:divprod},
$\ip\bu{\dvg\bA}=\dvg(\bA^{t}\bu)-\tr[\bA^{t}\grad\bu]$, together with
$\tr(\bT^{t}\bL)=T_{ij}L_{ij}=\bT\cdot\bL$ from
Theorem~\ref{thm:tip}. The left-hand side is
$\rho\,\ip{\dot\bv}\bv=\tfrac12\rho(\ip\bv\bv)^{\textstyle\cdot}$ by
the same product rule as in the particle case.

A reduction employing Cartesian coordinates is arguably more
enlightening:
\begin{equation}
   \ip{\dvg\bT}\bv=T_{ij,j}v_{i}=(T_{ij}v_{i})_{,j}-T_{ij}v_{i,j}
   =\dvg\big(\bT^{t}\bv\big)-\bT\cdot\bL .
   \tag{6.148}\label{eq:6.148}
\end{equation}
The middle equality is the product rule for a single ordinary
derivative, run backwards; the identification of the two resulting
terms uses $(\bT^{t}\bv)_{j}=T_{ij}v_{i}$ and $v_{i,j}=L_{ij}$ from
\eqr{4.1}. So the whole content of \eqr{6.147} is that
$T_{ij,j}v_{i}$ is not a divergence by itself and becomes one only at
the cost of the term $T_{ij}v_{i,j}$.

% reader's note: \eqr{6.147} assembled from nothing
\rnote{ch6:6.147-assembly}

Integrating \eqr{6.147} over an arbitrary part
$P_{t}=\bchi(S,t)\subset\kappa_{t}$ and invoking conservation of mass
in the form \eqr{6.24}$_{1}$ with $\varphi=\tfrac12|\bv|^{2}$ gives
\begin{equation}
   \frac{\dl}{\dl t}\int_{P_{t}}\tfrac12\rho|\bv|^{2}\,\dl v
   +\int_{P_{t}}\bT\cdot\bL\,\dl v
   =\int_{P_{t}}\rho\,\ip\bb\bv\,\dl v
   +\int_{\partial P_{t}}\ip\bt\bv\,\dl a ,
   \tag{6.149}\label{eq:6.149}
\end{equation}
where the time derivative pertains to the fixed set $S$ of material
points, and we have used the divergence theorem for smooth fields to
convert the integral of $\dvg(\bT^{t}\bv)$ over $P_{t}$ into an
integral over $\partial P_{t}$ of
\[
   \ip{\bT^{t}\bv}\bn=\ip\bv{\bT\bn}=\ip\bt\bv ,
\]
the first equality by the definition of the transpose and the second by
Cauchy's theorem \eqr{6.94}.

% reader's note: integrating \eqr{6.147}: the four terms, one at a
% time
\rnote{ch6:6.149-integrating}

Whence the \emph{mechanical energy balance}:
\begin{equation}
   \frac{\dl}{\dl t}\mathcal{K}(S,t)+\mathcal{S}(S,t)
   =\mathcal{P}(S,t),
   \tag{6.150}\label{eq:6.150}
\end{equation}
where
\begin{equation}
   \mathcal{K}(S,t)=\tfrac12\int_{P_{t}}\rho|\bv|^{2}\,\dl v
   \tag{6.151}\label{eq:6.151}
\end{equation}
is the kinetic energy of $S$,
\begin{equation}
   \mathcal{P}(S,t)=\int_{P_{t}}\rho\,\ip\bb\bv\,\dl v
   +\int_{\partial P_{t}}\ip\bt\bv\,\dl a
   \tag{6.152}\label{eq:6.152}
\end{equation}
is the power of the forces acting on $S$, and
\begin{equation}
   \mathcal{S}(S,t)=\int_{P_{t}}\bT\cdot\bL\,\dl v
   \tag{6.153}\label{eq:6.153}
\end{equation}
is the \emph{stress power} contained in $S$.

We observe that the stress power has no analogue in particle mechanics.
This is due to the fact that ideal particles have no surface area, so
the concepts of traction and stress have no meaning for them. For a
continuum, the applied power \eqr{6.152} does not all go into kinetic
energy: some of it is spent deforming the material, and \eqr{6.153} is
the rate at which that happens.

\begin{remark}[reading the stress power]\label{rem:ch6stresspower}
Two simplifications of $\bT\cdot\bL$ are used so often that they are
worth recording here.

\emph{Only the stretching does work.} Split $\bL=\bD+\bW$ by \eqr{4.7}.
Since $\bT$ is symmetric and $\bW$ is skew, $\bT\cdot\bW=0$ by
Lemma~\ref{lem:symskw}, so
\[
   \bT\cdot\bL=\bT\cdot\bD .
\]
Spin costs nothing. This is the algebraic expression of the fact that
rigidly rotating a stressed body does no work on it, and it is the
first thing Chapter 10 uses, at \eqr{10.15}.

\emph{Two states in which it vanishes.} For a pure pressure
\eqr{6.128}, $\bT\cdot\bL=-p\,\I\cdot\bL=-p\tr\bL=-p\dvg\bv
=-p\dot J/J$, which vanishes in an isochoric motion. For a rigid
motion, $\bD=\Ob$ by Remark~\ref{rem:Drigid} and the stress power
vanishes whatever $\bT$ is. Both facts return in
Section~\ref{sec:ch6energy} as the two cases in which the thermal and
mechanical problems decouple.
\end{remark}

\subsection*{Referential forms, and power conjugacy}

The referential forms of the kinetic energy and power are
\begin{equation}
   \mathcal{K}(S,t)=\tfrac12\int_{P}\rho_{\kappa}|\bv|^{2}\,\dl V
   \qquad\text{and}\qquad
   \mathcal{P}(S,t)=\int_{P}\rho_{\kappa}\,\ip\bb\bv\,\dl V
   +\int_{\partial P}\ip\bp\bv\,\dl A ,
   \tag{6.154}\label{eq:6.154}
\end{equation}
respectively, obtained by the substitutions
$\rho\,\dl v=\rho_{\kappa}\,\dl V$ and $\bt\,\dl a=\bp\,\dl A$ of
\eqr{6.142}. As for the stress power, we use \eqr{6.140} in the form
$\bT=J^{-1}\bP\bF^{t}$, which is \eqr{6.145} divided by $J$:
\begin{equation}
\begin{aligned}
   \int_{P_{t}}\bT\cdot\bL\,\dl v
   &=\int_{P_{t}}J^{-1}\bP\bF^{t}\cdot\bL\,\dl v\\
   &=\int_{P}\bP\bF^{t}\cdot\bL\,\dl V
    =\int_{P}\tr\big(\bP\bF^{t}\bL^{t}\big)\,\dl V .
\end{aligned}
   \tag{6.155}\label{eq:6.155}
\end{equation}
The second line used $\dl v=J\,\dl V$, whose factor $J$ cancels the
$J^{-1}$, and the third is Definition~\ref{def:tip}.

% reader's note: why \eqr{6.155} comes out so clean
\rnote{ch6:6.155-cancellation}

Using $\bF^{t}\bL^{t}=(\bL\bF)^{t}=\dot\bF^{t}$, which is \eqr{4.3},
$\bL=\dot\bF\bF^{-1}$, rearranged as $\dot\bF=\bL\bF$ and transposed,
we then have $\tr(\bP\bF^{t}\bL^{t})=\tr(\bP\dot\bF^{t})
=\bP\cdot\dot\bF$. In terms of components,
$\bP\cdot\dot\bF=P_{iA}\dot F_{iA}$. We note that whereas the trace of
a two-point tensor is not defined, the product $\bP\dot\bF^{t}$ is a
\emph{spatial} tensor, $\dot\bF^{t}$ carrying $\Vt$ to $\hVt$ and $\bP$
carrying $\hVt$ back to $\Vt$, so its trace is meaningful and furnishes
a well-defined inner product of two-point tensors. Accordingly we reach
\begin{equation}
   \mathcal{S}(S,t)=\int_{P}\bP\cdot\dot\bF\,\dl V .
   \tag{6.156}\label{eq:6.156}
\end{equation}

% reader's note: $\bP\cdot\dot\bF$: how a two-point tensor acquires
% an inner product
\rnote{ch6:6.156-innerproduct}
Indeed the energy balance \eqr{6.150}, with kinetic energy, power and
stress power given by \eqr{6.154} and \eqr{6.156}, may be derived
directly by repeating the foregoing procedure with \eqr{6.109} replaced
by \eqr{6.136}.

A stress measure that is juxtaposed with the material derivative of a
deformation measure in an expression for the stress power is said to be
\emph{power conjugate} to that deformation measure. Thus \eqr{6.156}
says that the Piola stress is power conjugate to the deformation
gradient. We show that the Piola--Kirchhoff stress is power conjugate
to the Lagrange strain $\bE$ defined by \eqr{2.97}. This is proved by
combining \eqr{2.83}$_{2}$, $\bC=\bF^{t}\bF$, with \eqr{2.97},
$\bE=\tfrac12(\bC-\hI)$, to obtain
\begin{equation}
\begin{aligned}
   \dot\bE&=\tfrac12\big(\bF^{t}\bF-\hI\big)^{\textstyle\cdot}
    =\tfrac12\big(\bF^{t}\dot\bF+\dot\bF^{t}\bF\big)\\
   &=\tfrac12\big[\bF^{t}\dot\bF+\big(\bF^{t}\dot\bF\big)^{t}\big]
    =\Sym\big(\bF^{t}\dot\bF\big),
\end{aligned}
   \tag{6.157}\label{eq:6.157}
\end{equation}
the constant $\hI$ dropping out on differentiation and
$(\bF^{t}\dot\bF)^{t}=\dot\bF^{t}\bF$ by the rule for the transpose of
a product. Then
\begin{equation}
   \bP\cdot\dot\bF=\bF\bS\cdot\dot\bF
   =\tr\big(\bF\bS\dot\bF^{t}\big)=\tr\big(\bS\dot\bF^{t}\bF\big),
   \tag{6.158}\label{eq:6.158}
\end{equation}
in which the first trace pertains to a spatial tensor, yielding an
inner product of two-point tensors, and the second pertains to a
referential tensor. That the two are equal is the component identity
\[
   \tr\big(\bF\bS\dot\bF^{t}\big)=F_{iA}S_{AB}\dot F_{iB}
   =S_{AB}\big(\dot F_{iB}F_{iA}\big)=\tr\big(\bS\dot\bF^{t}\bF\big),
\]
the same sum of products of the same numbers written in a different
order; it is the manoeuvre already used at \eqr{10.17} of Chapter 10
and justified there in the same way. Thus
$\bP\cdot\dot\bF=\bS\cdot\bF^{t}\dot\bF$, since
$\tr(\bS\bM)=\bS\cdot\bM^{t}$ with $\bM=\dot\bF^{t}\bF$ and
$\bM^{t}=\bF^{t}\dot\bF$, and the symmetry of the Piola--Kirchhoff
stress reduces this by Lemma~\ref{lem:symskw} to
\begin{equation}
   \bP\cdot\dot\bF=\bS\cdot\Sym\big(\bF^{t}\dot\bF\big),
   \tag{6.159}\label{eq:6.159}
\end{equation}
yielding, by \eqr{6.157}, the stress power in the form
\begin{equation}
   \mathcal{S}(S,t)=\int_{P}\bS\cdot\dot\bE\,\dl V .
   \tag{6.160}\label{eq:6.160}
\end{equation}

% reader's note: $\bP\cdot\dot\bF=\bS\cdot\dot\bE$ done entirely in
% components
\rnote{ch6:6.160-components}

\begin{remark}[what power conjugacy is good for]\label{rem:ch6conjugate}
Three expressions for the same stress power have now been obtained:
\[
\begin{aligned}
   &\bT\cdot\bD &&\text{per unit current volume},\\
   &\bP\cdot\dot\bF\ \text{ and }\ \bS\cdot\dot\bE
   &&\text{per unit reference volume}.
\end{aligned}
\]
They are equal, so nothing physical distinguishes them. What
distinguishes them is which pair of variables a constitutive equation
is written in.

Suppose the material stores energy: there is a function $W(\bF)$, per
unit reference volume, whose rate along every motion is the stress
power density. The hypothesis and the chain rule give two expressions
for the same scalar,
\[
   \dot W=W_{\bF}\cdot\dot\bF
   \qquad\text{and}\qquad
   \dot W=\bP\cdot\dot\bF ,
\]
the second being \eqr{6.156}. Subtracting,
\[
   \big(\bP-W_{\bF}\big)\cdot\dot\bF=0
   \qquad\text{for every }\dot\bF ,
\]
and since $\dot\bF$ ranges over all two-point tensors as the motion
varies, while the bracket does not depend on it, the bracket vanishes:
\[
   \bP=W_{\bF} .
\]
The same argument run on \eqr{6.160}, with $W$ regarded as a function
of $\bE$ and $\dot\bE$ ranging over all symmetric referential tensors,
gives $\bS=W_{\bE}$. Neither statement would be available if the pairing were wrong,
and this is why the list of conjugate pairs is extended. Problem
12 adds two more: the Biot stress conjugate to the right stretch $\bU$
and the Bell stress conjugate to the co-rotational rate of the left
stretch $\bV$.
\end{remark}

% reader's note: the cancellation argument in
% Remark~\ref{rem:ch6conjugate}, with its hypotheses exposed
\rnote{ch6:6.160-conjugacy}


%=======================================================================
\section{Energy balance}\label{sec:ch6energy}
%=======================================================================

The mechanical energy balance \eqr{6.150} does not account for the
interconvertibility of heat and mechanical power, a fundamental aspect
of thermomechanical phenomena. It was not derived from an axiom at all:
it is the linear momentum balance dotted with the velocity, so it
contains no information that \eqr{6.109} did not already contain. To
introduce heat we need a genuinely new postulate, the \emph{balance of
energy}. Thus we assume the existence of scalars $\mathcal{U}(S,t)$ and
$\mathcal{H}(S,t)$, the \emph{internal energy} and the \emph{heating}
of $S\subset B$ respectively, such that
\begin{equation}
   \frac{\dl}{\dl t}\big[\mathcal{K}(S,t)+\mathcal{U}(S,t)\big]
   =\mathcal{P}(S,t)+\mathcal{H}(S,t).
   \tag{6.161}\label{eq:6.161}
\end{equation}
This is the first law of thermodynamics for a part of a continuum: the
energy of $S$, kinetic plus internal, changes at the rate at which work
is done on it plus the rate at which it is heated.

We assume that $\mathcal{U}$ can be represented as an integral over
$P_{t}$, that is, that there exists a function $\eps(\bx,t)$, the
internal energy per unit mass, such that
\begin{equation}
   \mathcal{U}(S,t)=\int_{P_{t}}\rho\eps\,\dl v ,
   \tag{6.162}\label{eq:6.162}
\end{equation}
and that the heating is expressible in the form
\begin{equation}
   \mathcal{H}(S,t)=\int_{P_{t}}\rho r\,\dl v
   +\int_{\partial P_{t}}h\,\dl a ,
   \tag{6.163}\label{eq:6.163}
\end{equation}
where $r(\bx,t)$ is the heating per unit mass, through radioactive
decay for example, and $h(\bx,t)$ is the heating influx supplied via
conduction, for example by contact with material outside $P_{t}$. The
two terms of \eqr{6.163} are the bulk supply and the boundary flux of
the generic form \eqr{6.1}, and \eqr{6.162} is the assertion that
internal energy, like mass, is distributed and not concentrated.

% reader's note: where the internal energy came from, and what the
% first law asserts
\rnote{ch6:6.161-history}

Combining the mechanical energy balance with \eqr{6.161} produces the
\emph{thermal energy balance}. The mechanical balance \eqr{6.150} reads
$\dot{\mathcal{K}}+\mathcal{S}=\mathcal{P}$, so
\[
   \mathcal{P}=\dot{\mathcal{K}}+\mathcal{S} .
\]
Substitute this for $\mathcal{P}$ on the right of \eqr{6.161} and
expand the left:
\[
\begin{aligned}
   \dot{\mathcal{K}}+\dot{\mathcal{U}}
   &=\mathcal{P}+\mathcal{H}
   =\big(\dot{\mathcal{K}}+\mathcal{S}\big)+\mathcal{H}\\
   \dot{\mathcal{U}}&=\mathcal{S}+\mathcal{H} ,
\end{aligned}
\]
the second line by cancelling $\dot{\mathcal{K}}$, which occurs once on
each side. Rearranged,
\begin{equation}
   \mathcal{H}(S,t)=\frac{\dl}{\dl t}\mathcal{U}(S,t)-\mathcal{S}(S,t).
   \tag{6.164}\label{eq:6.164}
\end{equation}

% reader's note: the same cancellation, with every integral shown
\rnote{ch6:6.164-cancelK}

\begin{remark}[what the cancellation means]\label{rem:ch6cancelK}
The kinetic energy has dropped out of \eqr{6.164} entirely, and that is
the point of forming it. The mechanical energy balance is a
\emph{consequence} of the equations of motion and therefore adds
nothing to them; subtracting it from \eqr{6.161} removes exactly the
part of the first law that was already known, and leaves the part that
is new. What is new is \eqr{6.164}: internal energy changes through
heating and through the stress power.

The stress power is the only term the two balances share, and it is
therefore the sole channel of thermomechanical coupling. That
observation is illustrated by two cases.

\emph{Inviscid incompressible fluids.} If the stress is a pure
pressure, then by Remark~\ref{rem:ch6stresspower}
$\bT\cdot\bL=-p\tr\bL=-p\dot J/J$, and if the motion is isochoric,
$\dot J=0$, the stress power vanishes. There is then no coupling at
all, and the thermal and mechanical problems may be solved
independently.

\emph{Rigid bodies.} For a rigid motion $\bD=\Ob$ by
Remark~\ref{rem:Drigid}, so $\bT\cdot\bL=\bT\cdot\bW$, which vanishes
by the symmetry \eqr{6.117} and Lemma~\ref{lem:symskw}. Again the two
problems decouple. This is why traditional courses on heat transfer
treat bodies as rigid, with no consideration given to mechanical
response. It can only be valid approximately, since real bodies are not
perfectly rigid, just as real fluids are not perfectly inviscid; but
the approximation is a good one whenever the stress power is small
compared with the heating. Chapter 10
returns to this from the other side: a material that stores no strain
energy can also tell you nothing about its stress.
\end{remark}

Proceeding, we write the thermal energy balance, granted conservation
of mass, as
\begin{equation}
   \int_{\partial P_{t}}h\,\dl a
   =\int_{P_{t}}\big[\rho(\dot\eps-r)-\bT\cdot\bL\big]\,\dl v .
   \tag{6.165}\label{eq:6.165}
\end{equation}
The steps are: $\frac{\dl}{\dl t}\mathcal{U}
=\frac{\dl}{\dl t}\int_{P_{t}}\rho\eps\,\dl v
=\int_{P_{t}}\rho\dot\eps\,\dl v$ by \eqr{6.24}$_{1}$ with
$\varphi=\eps$; then \eqr{6.164} reads
$\int\rho r\,\dl v+\int h\,\dl a=\int\rho\dot\eps\,\dl v
-\int\bT\cdot\bL\,\dl v$, and moving the bulk terms to the right gives
\eqr{6.165}.

% reader's note: \eqr{6.165}, and why it is written with the surface
% integral alone on the left
\rnote{ch6:6.165-shape}

\subsection*{The heat flux vector}

Equation \eqr{6.165} has exactly the structure that Sections
\ref{sec:ch6noll} and \ref{sec:ch6cauchythm} exploited: a surface
integral of an unknown boundary field equal to a volume integral of
bounded quantities. Assuming $\rho$, $\dot\eps$, $r$ and $\bT\cdot\bL$
to be continuous, and hence bounded on compact subsets, we can
therefore adapt Noll's theorem to conclude that $h$ depends on
$\partial P_{t}$ through its local unit normal $\bn$, and then adapt
Cauchy's theorem to conclude that there is a vector field
$\bq(\bx,t)$, the \emph{heat flux vector}, such that
\begin{equation}
   h(\bx,t;\bn)=-\ip{\bq(\bx,t)}\bn ,
   \tag{6.166}\label{eq:6.166}
\end{equation}
in which the minus sign is inserted to conform to the convention that
$\ip\bq\bn>0$ corresponds to heat outflow, $h<0$, and $\ip\bq\bn<0$ to
heat inflow, $h>0$.

\begin{remark}[the adaptation is word for word]\label{rem:ch6heatcauchy}
Neither proof used the vectorial character of the traction, so both
apply verbatim to a scalar. The three places where $\bt$ entered are
these.

Noll's argument, \eqr{6.86} to \eqr{6.93}, used $\bt$ only through the
bound
\[
   \Big|\int_{\sigma}\bt\,\dl a\Big|
   \leq\big(\max_{\sigma}|\bt|\big)A(\sigma),
\]
and $|h|$ obeys the same bound with the same proof. The other input was
boundedness of the volume term in \eqr{6.165}, which is the hypothesis
just made.

Cauchy's argument then runs on the tetrahedron with $h$ in place of
$\bt$: the analogue of \eqr{6.105} is
\[
   \frac{1}{A(s)}\int_{\partial\tau}h\,\dl a
   =h(\bx^{*};\bm)-\sum_{i=1}^{3}\big(\ip{\be_{i}}\bm\big)
     h(\bx^{*}_{i};\be_{i}),
\]
the left side is $O(\delta)$ by the cone estimate of
Lemma~\ref{lem:ch6cone}, and the limit gives
$h(\bx_{0};\bm)=\sum_{i}(\ip{\be_{i}}\bm)h(\bx_{0};\be_{i})
=\ip{\ba}{\bm}$ with $\ba=h(\bx_{0};\be_{i})\be_{i}$. Setting
$\bq=-\ba$ gives \eqr{6.166}. So the heating influx is a linear
function of the normal for the same reason the traction is, and the
whole content of Fourier's law, when it comes in Chapter 8, will be a
constitutive statement about $\bq$ and not about $h$.

The minus sign is pure convention and could have been absorbed into
$\bq$. Its purpose is to make $\bq$ point in the direction heat
\emph{flows}, so that the second law will read $\ip\bq{\grad\theta}\leq0$,
heat flowing down the temperature gradient, rather than the reverse.
\end{remark}

Then
\begin{equation}
   \int_{\partial P_{t}}h\,\dl a=-\int_{\partial P_{t}}\ip\bq\bn\,\dl a
   =-\int_{P_{t}}\dvg\bq\,\dl v ,
   \tag{6.167}\label{eq:6.167}
\end{equation}
for smooth fields, by the divergence theorem \eqr{5.1}. The printed
text has $\partial P_{r}$ in the middle integral here and again in
\eqr{6.170}; $\partial P_{t}$ is meant in both places. Substituting
\eqr{6.167} into \eqr{6.165}, all terms become volume integrals over an
arbitrary $P_{t}$, and we can localize by
Theorem~\ref{thm:localization} to arrive at the local form
\begin{equation}
   \boxed{\ \rho\dot\eps=\bT\cdot\bL-\dvg\bq+\rho r ,\ }
   \tag{6.168}\label{eq:6.168}
\end{equation}
holding at points in $\kappa_{t}$ that do not belong to a singular
surface.

% reader's note: getting \eqr{6.168}, and reading each of its four
% terms
\rnote{ch6:6.168-localization}

Alternatively, this follows directly on making the appropriate
identifications in \eqr{6.1} and \eqr{6.21}, namely
\[
   \varphi=\eps,\qquad \vek\pi=-\bq,\qquad
   \rho\sigma=\rho r+\bT\cdot\bL ,
\]
and \eqr{6.6} then delivers the associated jump condition
\begin{equation}
   \ip{\jump{\rho\eps(\bv-\bu)+\bq}}\bn=0 ,
   \tag{6.169}\label{eq:6.169}
\end{equation}
holding at all points on a singular surface. Note that the stress power
has been counted as part of the bulk supply $\rho\sigma$ and therefore
contributes nothing to the jump condition, a volume term being
invisible on a surface of zero volume. In terms of the mass flux
$m^{*}$ of Remark~\ref{rem:ch6rh1}, \eqr{6.169} reads
$m^{*}\jump\eps+\ip{\jump\bq}\bn=0$: the internal energy convected
across the surface is balanced by the jump in the conducted heat.

% reader's note: the identification table, and the reason volume
% terms never appear in a jump condition
\rnote{ch6:6.169-identifications}

\subsection*{Referential form}

A referential form of the energy balance proceeds by rewriting the
thermal energy balance as
\begin{equation}
\begin{aligned}
   \int_{P}\big[\rho_{\kappa}(\dot\eps-r)-\bP\cdot\dot\bF\big]\,\dl V
   &=-\int_{\partial P_{t}}\ip\bq\bn\,\dl a
    =-\int_{\partial P}\ip\bq{\bF^{*}\bN}\,\dl A\\
   &=-\int_{\partial P}\ip{\bq_{\kappa}}\bN\,\dl A ,
\end{aligned}
   \tag{6.170}\label{eq:6.170}
\end{equation}
where
\begin{equation}
   \bq_{\kappa}(\bX,t)=\big(\bF^{*}\big)^{t}\bq=J\bF^{-1}\bq
   \tag{6.171}\label{eq:6.171}
\end{equation}
is the \emph{referential heat flux vector}, with components
\begin{equation}
   q_{(\kappa)A}=F^{*}_{iA}q_{i},
   \tag{6.172}\label{eq:6.172}
\end{equation}
in which the $q_{i}$ are those of the spatial vector $\bq$. The left
side of \eqr{6.170} is \eqr{6.165} pulled back: the volume term by
$\rho\,\dl v=\rho_{\kappa}\,\dl V$, and the stress power by
\eqr{6.155} and \eqr{6.156}. On the right, the first equality is
\eqr{6.167}$_{1}$, the second is Piola--Nanson \eqr{2.54}, and the
third moves $\bF^{*}$ from one slot of the inner product to the other
by the definition of the transpose, which is what \eqr{6.171} records.
The second form in \eqr{6.171} is \eqr{2.55}: $(\bF^{*})^{t}
=(J\bF^{-t})^{t}=J\bF^{-1}$.

Then, for smooth fields,
\begin{equation}
   \int_{\partial P}\ip{\bq_{\kappa}}\bN\,\dl A
   =\int_{P}\Dvg\bq_{\kappa}\,\dl V ,
   \qquad\text{where}\quad
   \Dvg\bq_{\kappa}=q_{(\kappa)A,A},
   \tag{6.173}\label{eq:6.173}
\end{equation}
and we localize, $P$ being arbitrary, to obtain
\begin{equation}
   \rho_{\kappa}\dot\eps=\bP\cdot\dot\bF-\Dvg\bq_{\kappa}+\rho_{\kappa}r
   \tag{6.174}\label{eq:6.174}
\end{equation}
at points $\bX\in\kappa$ removed from a singular surface. The
derivation of the associated jump condition is Problem 1.

This result may also be derived directly if \eqr{6.168} is multiplied
by $J$ and use is made of the easily proved identity
\begin{equation}
   J\dvg\bq=\Dvg\big(J\bF^{-1}\bq\big),
   \tag{6.175}\label{eq:6.175}
\end{equation}
which is Problem 8. Multiplying \eqr{6.168} by $J$ turns
$\rho\dot\eps$ into $\rho_{\kappa}\dot\eps$ and $\rho r$ into
$\rho_{\kappa}r$ by \eqr{2.59}, turns $J\bT\cdot\bL$ into
$\bP\cdot\dot\bF$ by the computation of \eqr{6.155}, and turns
$J\dvg\bq$ into $\Dvg\bq_{\kappa}$ by \eqr{6.175}. Note the shape of
\eqr{6.175}: it is \eqr{6.138} for a vector instead of a tensor, and
both are consequences of the Piola identity $\Dvg\bF^{*}=\bze$ of
\eqr{5.11}. That identity is the reason every referential balance law
in this chapter has exactly the same form as its spatial counterpart,
with no leftover terms.

% reader's note: the referential energy balance, term by term, both
% ways
\rnote{ch6:6.175-referential}

%=======================================================================
\section{Entropy and the Clausius--Duhem
inequality}\label{sec:ch6entropy}
%=======================================================================

To account for the thermodynamical concept of dissipation attending any
real physical process, we introduce a further hypothesis, the
\emph{entropy balance}. This is assumed to take the form
\begin{equation}
   \frac{\dl}{\dl t}\int_{P_{t}}\rho\eta\,\dl v
   =\int_{P_{t}}\theta^{-1}(\rho r)\,\dl v
   +\int_{\partial P_{t}}\theta^{-1}h\,\dl a
   +\int_{P_{t}}\rho s\,\dl v ,
   \tag{6.176}\label{eq:6.176}
\end{equation}
where $\eta(\bx,t)$ is the entropy per unit mass contained in $S$,
$\theta(\bx,t)>0$ is the absolute temperature, and $s$ is the entropy
production rate per unit mass.

% reader's note: why anyone needed entropy --- Carnot, Clausius,
% Boltzmann
\rnote{ch6:6.176-history}

The first integral on the right represents the entropy supplied by
sources in the bulk and the second the flux of entropy through the
surface. Their forms are suggested by the classical thermostatics of
reversible processes, in which an amount of heat $\dl Q$ delivered at
temperature $\theta$ carries entropy $\dl Q/\theta$; each heating term
of \eqr{6.163} accordingly reappears in \eqr{6.176} divided by
$\theta$. The last integral, which vanishes in a reversible process,
represents the net entropy production attending an irreversible one.
These are, by definition, such that this integral is positive and hence
that $s>0$. All admissible processes are then subject to the
\emph{Clausius--Duhem inequality}
\begin{equation}
   \frac{\dl}{\dl t}\int_{P_{t}}\rho\eta\,\dl v
   \geq\int_{P_{t}}\theta^{-1}(\rho r)\,\dl v
   +\int_{\partial P_{t}}\theta^{-1}h\,\dl a ,
   \tag{6.177}\label{eq:6.177}
\end{equation}
obtained by deleting the non-negative production term from
\eqr{6.176}.

% reader's note: reading \eqr{6.176} against \eqr{6.163}, and the
% one step where the inequality is born
\rnote{ch6:6.177-reading}

\subsection*{Truesdell's motivation for \eqr{6.177}}

The inequality may be motivated by an argument due to Truesdell. We
assume that there is a function $\mathcal{B}(S,t)$, depending on
material constitution, that bounds the heating from above,
$\mathcal{B}(S,t)\geq\mathcal{H}(S,t)$. From the energy balance
\eqr{6.161} we then have that, in accordance with experience, there is
a limit to the rate at which heat can be converted to energy in the
absence of forces, $\mathcal{P}=0$:
\begin{equation}
   \frac{\dl}{\dl t}\big[\mathcal{K}(S,t)+\mathcal{U}(S,t)\big]
   \leq\mathcal{B}(S,t).
   \tag{6.178}\label{eq:6.178}
\end{equation}
For uniformly distributed temperatures, $\grad\theta=\bze$, we define
the entropy
\begin{equation}
   \mathcal{C}(S,t)=\int_{P_{t}}\rho\eta\,\dl v
   \tag{6.179}\label{eq:6.179}
\end{equation}
by
\begin{equation}
   \mathcal{C}(S,t)=\int_{t_{0}}^{t}\theta(\tau)^{-1}
   \mathcal{B}(S,\tau)\,\dl\tau ,
   \tag{6.180}\label{eq:6.180}
\end{equation}
where $t_{0}$ is a fixed time. Then, differentiating \eqr{6.180} by the
fundamental theorem of calculus and multiplying by $\theta(t)$,
\begin{equation}
   \theta(t)\frac{\dl}{\dl t}\mathcal{C}(S,t)=\mathcal{B}(S,t)
   \geq\mathcal{H}(S,t),
   \tag{6.181}\label{eq:6.181}
\end{equation}
which is just the inequality \eqr{6.177} with the uniform reciprocal
temperature factored outside the integrals on the right-hand side:
dividing by $\theta$ and writing $\mathcal{H}$ out by \eqr{6.163},
\[
   \frac{\dl}{\dl t}\int_{P_{t}}\rho\eta\,\dl v
   \geq\theta^{-1}\left[\int_{P_{t}}\rho r\,\dl v
   +\int_{\partial P_{t}}h\,\dl a\right].
\]
Placing $\theta^{-1}$ inside the integrals, which changes nothing while
$\theta$ is uniform, and making the further assumption that the
resulting inequality also applies to non-uniform temperature fields,
then yields the Clausius--Duhem inequality.

\begin{remark}[how much of this is proof]\label{rem:ch6truesdell}
Very little, and it is important to say so. The argument assumes the
existence of $\mathcal{B}$, defines $\mathcal{C}$ by \eqr{6.180} so
that \eqr{6.181} holds identically, asserts without further evidence
that the $\mathcal{C}$ so defined is the integral \eqr{6.179} of a mass
density $\eta$, and then extends an inequality established for uniform
temperature fields to non-uniform ones by fiat. The last step is the
substantive one and it is not deducible from the preceding.

What the argument does is show that \eqr{6.177} is the natural
continuum generalization of a statement one is willing to make about a
body at uniform temperature, namely that there is a ceiling on the rate
at which heat can be absorbed and that the ceiling is proportional to
the temperature. The inequality is a postulate, on the same footing as
\eqr{6.34} and \eqr{6.161}, and its justification is that the theory
built on it agrees with experiment.

The entropy supply and flux written in \eqr{6.176} are also not the
most general that could be contemplated. One could allow an entropy
flux not equal to $\theta^{-1}h$, and theories of that kind exist. The
classical theory is the one that makes the simplest choice.
\end{remark}

\subsection*{Local forms}

Combining \eqr{6.177} with \eqr{6.166} and localizing as usual yields
the local inequality. The identifications in the generic scalar balance
\eqr{6.1}, now read with $\geq$ in place of $=$, are
\[
   \varphi=\eta,\qquad \sigma=\theta^{-1}r,\qquad
   \vek\pi=-\theta^{-1}\bq ,
\]
the last because the surface integrand is
$\theta^{-1}h=-\theta^{-1}\ip\bq\bn=\ip{(-\theta^{-1}\bq)}\bn$. Then
\eqr{6.21} gives
\begin{equation}
   \rho\dot\eta\geq\theta^{-1}(\rho r)-\dvg\big(\theta^{-1}\bq\big),
   \tag{6.182}\label{eq:6.182}
\end{equation}
holding in regions of smoothness, whereas the jump condition
\begin{equation}
   \ip{\jump{\rho\eta(\bv-\bu)+\theta^{-1}\bq}}\bn\geq0 ,
   \tag{6.183}\label{eq:6.183}
\end{equation}
operative on singular surfaces, follows by making the same
identifications in \eqr{6.6}, with the equality replaced by an
inequality.

Here we recall, with reference to Figure~\ref{fig:ch5shock} of Chapter
5, that $\bn$ is directed into the ``$+$'' region, and hence, with the
definition \eqr{5.28} of the jump, that this inequality is independent
of the orientation of $\bn$. The check is one line. Reversing $\bn$
exchanges the labels $+$ and $-$, so every jump changes sign,
$\jump\cdot\mapsto-\jump\cdot$; and $\bn$ itself changes sign. The
product of the two therefore does not change, so the left side of
\eqr{6.183} is unaltered and the inequality means the same thing. Had
the two sign changes not cancelled, \eqr{6.183} would have asserted
that a shock dissipates when read from one side and creates entropy
from the other, which is not a statement about the world but about the
person reading it.

% reader's note: the identifications for the entropy, and the
% orientation check written out
\rnote{ch6:6.183-orientation}

\subsection*{The reduced inequality}

A more useful form of the local inequality \eqr{6.182} is obtained by
eliminating the radiative heating $r$, which is external data and which
no constitutive equation should have to mention. Solve the thermal
energy balance \eqr{6.168} for it,
\[
   \rho r=\rho\dot\eps-\bT\cdot\bL+\dvg\bq ,
\]
and substitute into \eqr{6.182}:
\begin{equation}
   \rho\dot\eta\geq\theta^{-1}\big(\rho\dot\eps-\bT\cdot\bL+\dvg\bq\big)
   -\dvg\big(\theta^{-1}\bq\big),
   \tag{6.184}\label{eq:6.184}
\end{equation}
which, since $\theta>0$, is equivalent to the inequality obtained on
multiplying through by $\theta$. Doing so, and expanding the last term
by Problem 24(c) of Chapter 1,
\[
   \dvg\big(\theta^{-1}\bq\big)
   =\ip{\grad(\theta^{-1})}\bq+\theta^{-1}\dvg\bq
   =-\theta^{-2}\ip{\grad\theta}\bq+\theta^{-1}\dvg\bq ,
\]
where $\grad(\theta^{-1})=-\theta^{-2}\grad\theta$ by the chain rule,
Theorem~\ref{thm:chain}. Hence
\[
   \theta\,\dvg\big(\theta^{-1}\bq\big)
   =-\theta^{-1}\ip{\grad\theta}\bq+\dvg\bq ,
\]
and multiplying \eqr{6.184} by $\theta$ gives
\[
   \rho\theta\dot\eta\geq\rho\dot\eps-\bT\cdot\bL+\dvg\bq
   +\theta^{-1}\ip\bq{\grad\theta}-\dvg\bq .
\]
The two divergences cancel, which is the whole purpose of the
expansion, and rearranging so that everything stands on one side,
\begin{equation}
   \rho\big(\theta\dot\eta-\dot\eps\big)+\bT\cdot\bL
   -\theta^{-1}\ip\bq{\grad\theta}\geq0 .
   \tag{6.185}\label{eq:6.185}
\end{equation}

On introducing the \emph{Helmholtz free energy} per unit mass $\psi$,
defined by
\begin{equation}
   \psi=\eps-\theta\eta ,
   \tag{6.186}\label{eq:6.186}
\end{equation}
we finally reduce the local Clausius--Duhem inequality to a form in
which the entropy appears only through $\psi$ and $\theta$.
Differentiating \eqr{6.186} along a material point,
\[
   \dot\psi=\dot\eps-\dot\theta\eta-\theta\dot\eta ,
   \qquad\text{so}\qquad
   \theta\dot\eta-\dot\eps=-\dot\psi-\eta\dot\theta ,
\]
and substituting into \eqr{6.185},
\begin{equation}
   \boxed{\ -\rho\big(\dot\psi+\eta\dot\theta\big)+\bT\cdot\bL
   -\theta^{-1}\ip\bq{\grad\theta}\geq0 .\ }
   \tag{6.187}\label{eq:6.187}
\end{equation}
This will play a central role in the discussion of constitutive
equations in subsequent chapters.

% reader's note: the reduced inequality as one chain, and the
% referential form written out
\rnote{ch6:6.187-reduction}

The referential form of the Clausius--Duhem inequality, and its
associated local form and jump condition, are obtained by the same
route as \eqr{6.174}; the reader's note just given writes them out.
The only new ingredient is that $\grad\theta$ is replaced through
$\Grad\theta=\bF^{t}\grad\theta$, which is
Theorem~\ref{thm:GradgradF}.

\begin{remark}[why \eqr{6.187} is the form that gets
used]\label{rem:ch6cdread}
Three features of \eqr{6.187} explain why it, and not \eqr{6.177}, is
what Chapters 8 and 9 quote.

\emph{The supplies have gone.} Neither $r$ nor any boundary datum
appears. Everything in \eqr{6.187} is either a response of the material
($\psi$, $\eta$, $\bT$, $\bq$) or a description of the process
($\bL$, $\theta$, $\grad\theta$). An inequality of that shape can be
turned into a restriction on constitutive equations, which is exactly
what Chapter 8 does with it.

\emph{The free energy replaced the internal energy.} The change of
variable \eqr{6.186} is a Legendre transform trading $\eta$ for
$\theta$. It matters because $\theta$ is measurable and controllable
while $\eta$ is neither: a constitutive equation of the form
$\psi=\hat\psi(\bF,\theta)$ can be fitted to experiments, one of the
form $\eps=\hat\eps(\bF,\eta)$ cannot. The pairing $\eta$ with
$-\partial\psi/\partial\theta$ that Chapter 8 derives is the
thermodynamic analogue of the power conjugacy of
Remark~\ref{rem:ch6conjugate}.

\emph{The reaction stresses are invisible to it.} The stress enters
\eqr{6.187} only through $\bT\cdot\bL$, so any part of $\bT$ that is
workless in every admissible motion drops out identically. That is the
observation on which Chapter 10 builds the
whole theory of constraints on, and it is the reason the pressure in an
incompressible fluid is thermodynamically indeterminate.

Finally, note that \eqr{6.187} is a restriction on \emph{materials},
not on processes. Read as a restriction on processes it would say that
certain motions are forbidden, which cannot be right, since one may
impose whatever boundary data one likes. What it says is that a
constitutive equation must be such that \eqr{6.187} holds for every
process the material can undergo, and constitutive equations that fail
that test are discarded.
\end{remark}

% reader's note: the Coleman--Noll theorem, and the dissipation split
\rnote{ch6:6.187-colemannoll}

%=======================================================================
\setcounter{problemenv}{0}
\section{Problems}\label{sec:problems:c6}
%=======================================================================

%-----------------------------------------------------------------------
\begin{problem}[The referential balance laws]
Derive the referential forms of the local balance laws and jump
conditions discussed in Section~\ref{sec:ch6generic}. \emph{Hint:}
recall Problem 3 of Chapter 5.
\end{problem}

\begin{solution}
Every local statement of Section~\ref{sec:ch6generic} was assembled from
two theorems and a quantifier, and none of the three knows which
configuration it is working in. What changes on passing to $\kappa$ is
one thing only: the material stands still. Problem 3 of Chapter 5
carried out the consequence, deleting from \eqr{5.31} and \eqr{5.34}
every term built from the material velocity, and what survives is
\begin{equation}
\begin{aligned}
   \frac{\dl}{\dl t}\int_{P}\hat\Phi\,\dl V
   &=\int_{P}\frac{\partial\hat\Phi}{\partial t}\,\dl V
   -\int_{S}\jump{\hat\Phi}\,U\,\dl A ,\\
   \int_{\partial P}\ip\bA\bN\,\dl A
   &=\int_{P}\Dvg\bA\,\dl V+\int_{S}\ip{\jump\bA}\bN\,\dl A .
\end{aligned}
   \tag{$\ast_{3}$}\label{eq:ch6refthms}
\end{equation}
Nothing crosses $\partial P$, so there is no convective flux there; the
only motion left in the picture belongs to the singular surface $S$,
which sweeps through the material at speed $U$ along $\bN$. Everything
below is these two formulas applied to a balance law, and the interest
lies in exactly one of the three terms.

\medskip
\noindent\emph{Transferring the three terms.}
Two of the three move without thought. For the left-hand side and the
bulk supply of \eqr{6.1} the substitution is
\[
   \rho\,\dl v
   \overset{\eqr{2.50}}{=}\rho\,\big(J\,\dl V\big)
   =\big(\rho J\big)\,\dl V
   \overset{\eqr{2.59}}{=}\rho_{\kappa}\,\dl V ,
\]
so that
\[
   \int_{P_{t}}\rho\varphi\,\dl v=\int_{P}\rho_{\kappa}\varphi\,\dl V ,
   \qquad
   \int_{P_{t}}\rho\sigma\,\dl v=\int_{P}\rho_{\kappa}\sigma\,\dl V .
\]
The fields $\varphi$ and $\sigma$ are untouched. They are reckoned per
unit \emph{mass}, and the mass of a material element does not depend on
which configuration is used to measure its volume; had they been given
per unit volume, each would have acquired a factor $J$ and the final
equations would not have had the shape they do.

The boundary flux is the term that costs something, and it costs two
moves. First Piola--Nanson \eqr{2.54} converts the oriented area
element, $\bn\,\dl a=\bF^{*}\bN\,\dl A$, which in components reads
$n_{i}\,\dl a=F^{*}_{iA}N_{A}\,\dl A$; and then the tensor so produced
stands in the wrong slot of the inner product and has to be moved. In
components,
\[
\begin{aligned}
   \ip{\vek\pi}{\bF^{*}\bN}
   &=\pi_{i}\big(\bF^{*}\bN\big)_{i}
    =\pi_{i}F^{*}_{iA}N_{A}
    =\big(F^{*}_{iA}\pi_{i}\big)N_{A}\\
   &=\Big(\big(\bF^{*}\big)^{t}\vek\pi\Big)_{A}N_{A}
    =\ip{\big(\bF^{*}\big)^{t}\vek\pi}{\bN} ,
\end{aligned}
\]
which is the definition of the transpose written out. Hence
\[
   \int_{\partial P_{t}}\ip{\vek\pi}\bn\,\dl a
   =\int_{\partial P}\ip{\vek\pi}{\bF^{*}\bN}\,\dl A
   =\int_{\partial P}\ip{\big(\bF^{*}\big)^{t}\vek\pi}\bN\,\dl A ,
\]
and the referential flux is forced on us rather than chosen:
\[
   \boxed{\ \vek\pi_{\kappa}=\big(\bF^{*}\big)^{t}\vek\pi
   =J\bF^{-1}\vek\pi ,
   \qquad
   \pi_{(\kappa)A}=F^{*}_{iA}\pi_{i} .\ }
\]
The second form follows from \eqr{2.55}: transposing $\bF^{*}=J\bF^{-t}$
gives $(\bF^{*})^{t}=(J\bF^{-t})^{t}=J\big((\bF^{-1})^{t}\big)^{t}
=J\bF^{-1}$. That this is the same rule \eqr{6.171} met for the heat
flux is no coincidence: the heat flux was a particular $\vek\pi$, and
the calculation there was this one with a letter changed. The subscript
$\kappa$ is used rather than a capital $\Pi$ because $\vek\Pi$ already
denotes referential relative position in \eqr{6.53}. Collecting the
three terms, \eqr{6.1} reads
\begin{equation}
   \frac{\dl}{\dl t}\int_{P}\rho_{\kappa}\varphi\,\dl V
   =\int_{P}\rho_{\kappa}\sigma\,\dl V
   +\int_{\partial P}\ip{\vek\pi_{\kappa}}\bN\,\dl A .
   \tag{$\ast_{4}$}\label{eq:ch6refbal}
\end{equation}

\medskip
\noindent\emph{Localizing.}
Apply \eqref{eq:ch6refthms}$_{1}$ with $\hat\Phi=\rho_{\kappa}\varphi$
to the left of \eqref{eq:ch6refbal}:
\[
   \frac{\dl}{\dl t}\int_{P}\rho_{\kappa}\varphi\,\dl V
   =\int_{P}\frac{\partial(\rho_{\kappa}\varphi)}{\partial t}\,\dl V
   -\int_{S}\jump{\rho_{\kappa}\varphi}\,U\,\dl A ,
\]
and \eqref{eq:ch6refthms}$_{2}$ with $\bA=\vek\pi_{\kappa}$ to the
boundary term on the right:
\[
   \int_{\partial P}\ip{\vek\pi_{\kappa}}\bN\,\dl A
   =\int_{P}\Dvg\vek\pi_{\kappa}\,\dl V
   +\int_{S}\ip{\jump{\vek\pi_{\kappa}}}\bN\,\dl A .
\]
Substituting both into \eqref{eq:ch6refbal} and moving everything to
one side,
\begin{equation}
   \int_{P}\left\{\frac{\partial(\rho_{\kappa}\varphi)}{\partial t}
   -\Dvg\vek\pi_{\kappa}-\rho_{\kappa}\sigma\right\}\dl V
   -\int_{S}\Big\{\jump{\rho_{\kappa}\varphi}U
   +\ip{\jump{\vek\pi_{\kappa}}}\bN\Big\}\dl A=0 .
   \tag{$\ast_{5}$}\label{eq:ch6reflocal}
\end{equation}

The two-stage argument of Section~\ref{sec:ch6generic} now runs
verbatim, and it is worth running rather than citing. Choose $P$
disjoint from $S$. The surface integral is then absent, and
\[
   \int_{P}\left\{\frac{\partial(\rho_{\kappa}\varphi)}{\partial t}
   -\Dvg\vek\pi_{\kappa}-\rho_{\kappa}\sigma\right\}\dl V=0
\]
for every such $P$. Were the integrand non-zero at some
$\bX_{0}\in\kappa$, continuity would keep it of one sign on a ball
about $\bX_{0}$, and integrating over that ball alone would give a
non-zero answer, contradicting the hypothesis. Hence the brace vanishes
pointwise. Now let $P$ meet $S$. The volume integral in
\eqref{eq:ch6reflocal} vanishes by what has just been proved, whatever
$P$ is, so
\[
   \int_{S}\Big\{\jump{\rho_{\kappa}\varphi}U
   +\ip{\jump{\vek\pi_{\kappa}}}\bN\Big\}\dl A=0 ,
\]
and $S\cap P$ is an arbitrary piece of the singular surface, so the same
argument on surfaces kills the brace there too. Finally, since
$\rho_{\kappa}$ is independent of $t$ by \eqr{6.18},
\[
   \frac{\partial(\rho_{\kappa}\varphi)}{\partial t}
   =\rho_{\kappa}\frac{\partial\varphi}{\partial t}
   +\varphi\underbrace{\frac{\partial\rho_{\kappa}}{\partial t}}_{=\,0}
   =\rho_{\kappa}\dot\varphi ,
\]
the material derivative in $\kappa$ being the partial derivative at
fixed $\bX$. Therefore
\[
   \boxed{\ \rho_{\kappa}\dot\varphi=\Dvg\vek\pi_{\kappa}
   +\rho_{\kappa}\sigma ,
   \qquad
   \jump{\rho_{\kappa}\varphi}U
   +\ip{\jump{\vek\pi_{\kappa}}}\bN=0 ,\ }
\]
the counterparts of \eqr{6.21} and \eqr{6.6}.

% reader's note: the sign comparison, done slowly
\rnote{ch6:prob1-signs}

\medskip
\noindent\emph{The vector case.}
Nothing changes except the alphabet. Let $\vek\psi$ have components
$\psi_{i}$ and let the flux be the tensor $\bA$. Piola--Nanson gives,
componentwise,
\[
   \big(\bA\bn\big)_{i}\,\dl a=A_{ij}n_{j}\,\dl a
   =A_{ij}F^{*}_{jA}N_{A}\,\dl A
   =\big(\bA\bF^{*}\big)_{iA}N_{A}\,\dl A ,
\]
so that $\int_{\partial P_{t}}\bA\bn\,\dl a
=\int_{\partial P}\bA\bF^{*}\bN\,\dl A$ and the referential flux is the
two-point tensor
\[
   \bA_{\kappa}=\bA\bF^{*},
   \qquad
   A_{(\kappa)iA}=A_{ij}F^{*}_{jA} .
\]
Applying the scalar result to each component $\psi_{i}$ with flux
$\vek\pi=\bA^{t}\be_{i}$, and reassembling on the fixed basis
$\{\be_{i}\}$ --- the caveat of Remark~\ref{rem:symbolic} being that the
reassembly is componentwise and the basis must not move ---
\[
   \boxed{\ \rho_{\kappa}\dot{\vek\psi}=\Dvg\bA_{\kappa}
   +\rho_{\kappa}\br ,
   \qquad
   \jump{\rho_{\kappa}\vek\psi}U+\jump{\bA_{\kappa}}\bN=\bze .\ }
\]

\medskip
\noindent\emph{The three instances.}
What remains is bookkeeping, and it is worth doing in full because it
shows how little is left to do: each balance law is a matter of naming
three letters.

\emph{Mass.} Here $\varphi=1$, $\sigma=0$, $\vek\pi=\bze$, so
$\vek\pi_{\kappa}=(\bF^{*})^{t}\bze=\bze$ and $\dot\varphi=0$. The
field equation reads $\rho_{\kappa}\cdot0=\Dvg\bze+0$, which is
vacuous; the content sits in the step before it, namely
$\partial(\rho_{\kappa}\cdot1)/\partial t=0$, that is
\[
   \frac{\partial\rho_{\kappa}}{\partial t}=0 ,
\]
which is \eqr{6.18}. The jump condition is
$\jump{\rho_{\kappa}\cdot1}U+\ip{\jump{\bze}}\bN=0$, that is
\[
   \jump{\rho_{\kappa}}\,U=0 .
\]
So either $U=0$, the surface being stationary in $\kappa$ and hence
material, or $\jump{\rho_{\kappa}}=0$ and the referential density is
continuous across it.

\emph{Linear momentum.} Here $\vek\psi=\bv$, $\br=\bb$, $\bA=\bT$, so
that
\[
   \bA_{\kappa}=\bT\bF^{*}\overset{\eqr{6.140}}{=}\bP ,
\]
and the two results read
\[
   \rho_{\kappa}\dot\bv=\Dvg\bP+\rho_{\kappa}\bb ,
   \qquad
   \jump{\rho_{\kappa}\bv}U+\jump\bP\bN=\bze ,
\]
the first being \eqr{6.136}.

\emph{Energy.} Here $\varphi=\eps$ and $\vek\pi=-\bq$, so that
\[
   \vek\pi_{\kappa}=\big(\bF^{*}\big)^{t}(-\bq)
   =-\big(\bF^{*}\big)^{t}\bq
   \overset{\eqr{6.171}}{=}-\bq_{\kappa} ,
\]
while the supply carries the stress power,
$\rho_{\kappa}\sigma=\rho_{\kappa}r+\bP\cdot\dot\bF$. The field
equation is
\[
   \rho_{\kappa}\dot\eps=\Dvg\big(-\bq_{\kappa}\big)
   +\rho_{\kappa}r+\bP\cdot\dot\bF
   =\bP\cdot\dot\bF-\Dvg\bq_{\kappa}+\rho_{\kappa}r ,
\]
which is \eqr{6.174}, and the jump condition is
\[
   \jump{\rho_{\kappa}\eps}U+\ip{\jump{-\bq_{\kappa}}}\bN=0 ,
   \qquad\text{that is}\qquad
   \jump{\rho_{\kappa}\eps}U-\ip{\jump{\bq_{\kappa}}}\bN=0 ,
\]
which is the condition the book leaves to the reader after \eqr{6.174}.
The generic statement was proved once; these are three substitutions
into it.

\medskip
\noindent\emph{A consistency check.}
The spatial and referential jump conditions look different and must say
the same thing. Problem 3 of Chapter 5 recorded the relation between
the two normal speeds,
\[
   u-\ip\bv\bn=\frac{J}{|\bF^{*}\bN|}\,U ,
\]
so the mass flux of Remark~\ref{rem:ch6rh1} is
\[
   m^{*}=\rho\big(\ip\bv\bn-u\big)
   =-\rho\big(u-\ip\bv\bn\big)
   =-\rho\,\frac{J}{|\bF^{*}\bN|}\,U
   =-\frac{\rho J}{|\bF^{*}\bN|}\,U
   =-\frac{\rho_{\kappa}}{|\bF^{*}\bN|}\,U ,
\]
using $\rho J=\rho_{\kappa}$ at the last step. Now $U$ and $\bN$ belong
to the surface and are single-valued on it. So is $|\bF^{*}\bN|$, and
the reason is worth stating: the motion $\bchi_{\kappa}$ is continuous
across $S$, so its derivatives \emph{along} $S$ agree when computed from
either side, and $\bF^{*}\bN\,\dl A=\bn\,\dl a$ is the cross product of
two vectors tangent to $S$, hence built from precisely those tangential
derivatives. Therefore the factor $-1/|\bF^{*}\bN|$ passes through the
jump bracket, and
\[
   \jump{m^{*}}
   =-\frac{U}{|\bF^{*}\bN|}\,\jump{\rho_{\kappa}} ,
\]
so that $\jump{m^{*}}=0$, which is the spatial condition \eqr{6.15},
holds if and only if $\jump{\rho_{\kappa}}U=0$, which is the referential
one. The other two convert by the same substitution.
\end{solution}

%-----------------------------------------------------------------------
\begin{problem}[The referential center of mass]
Verify that $\bar{\bX}$ appearing in \eqr{6.53} is in fact the position
of the center of mass of $P\subset\kappa$.
\end{problem}

\begin{solution}
The center of mass in $\kappa$ can only be defined one way, by copying
\eqr{6.40} with referential quantities throughout:
\begin{equation}
   m(S)\,\bar{\bX}=\int_{P}\rho_{\kappa}\bX\,\dl V ,
   \qquad
   m(S)\,\bar X_{A}=\int_{P}\rho_{\kappa}X_{A}\,\dl V .
   \tag{$\ast_{6}$}\label{eq:ch6refcom}
\end{equation}
What makes this worth writing down is that neither factor on the right
depends on $t$: $\rho_{\kappa}=\rho_{\kappa}(\bX)$ by \eqr{6.18}, and
$P$ is a fixed region of $\kappa$. So $\bar\bX$ is a \emph{constant}
vector, once and for all, and
\[
   \dot{\bar\bX}=\bze .
\]
That is exactly what \eqr{6.53} needs of it, and it is why the
definition is legitimate in a way its spatial counterpart is not:
$\bar\bx(t)$ moves.

The mass $m(S)$ appearing in \eqref{eq:ch6refcom} is the same number in
both configurations. By \eqr{2.56} together with conservation of mass,
\[
   m(S)=\int_{P_{t}}\rho\,\dl v
   =\int_{P}\rho J\,\dl V
   =\int_{P}\rho_{\kappa}\,\dl V ,
\]
so that it may be used on either side without comment.

\medskip
\noindent\emph{The two centers correspond.}
The question is whether $\bar\bX$ maps to $\bar\bx$ under the motion,
and the answer turns on a single structural feature of \eqr{6.52}.
Transfer \eqr{6.40} to $\kappa$:
\[
   m(S)\,\bar{\bx}(t)=\int_{P_{t}}\rho\bx\,\dl v
   =\int_{P}\rho J\,\bx\,\dl V
   =\int_{P}\rho_{\kappa}\,\bx\,\dl V ,
\]
and substitute the rigid motion $\bx=\bQ(t)\bX+\bc(t)$. Because $\bQ$
and $\bc$ depend on $t$ alone, they are constant over the region of
integration, and the integral splits. In components, with
$x_{i}=Q_{iA}X_{A}+c_{i}$,
\[
\begin{aligned}
   \int_{P}\rho_{\kappa}x_{i}\,\dl V
   &=\int_{P}\rho_{\kappa}\big(Q_{iA}X_{A}+c_{i}\big)\,\dl V\\
   &=Q_{iA}\int_{P}\rho_{\kappa}X_{A}\,\dl V
    +c_{i}\int_{P}\rho_{\kappa}\,\dl V\\
   &=Q_{iA}\,m(S)\bar X_{A}+c_{i}\,m(S) ,
\end{aligned}
\]
by \eqref{eq:ch6refcom} and by the expression for $m(S)$ above. In
direct notation this is
\[
   \int_{P}\rho_{\kappa}\big(\bQ\bX+\bc\big)\,\dl V
   =\bQ\big(m(S)\bar{\bX}\big)+\bc\,m(S) ,
\]
$\bQ$ passing outside as a constant linear map and $\bc$ as a constant
vector. Dividing by the positive number $m(S)$,
\[
   \boxed{\ \bar{\bx}(t)=\bQ(t)\bar{\bX}+\bc(t).\ }
\]
The mass center is therefore carried by the very map that carries every
other material point. Subtracting this from \eqr{6.52},
\[
   \vek\pi=\bx-\bar{\bx}=\big(\bQ\bX+\bc\big)
   -\big(\bQ\bar{\bX}+\bc\big)
   =\bQ\bX-\bQ\bar\bX
   =\bQ\big(\bX-\bar{\bX}\big)=\bQ\vek\Pi ,
\]
which is \eqr{6.53}, the third equality by linearity of $\bQ$. The
translation cancels, as it must: where the body happens to be cannot
affect the position of one of its points relative to another.

\medskip
\noindent\emph{The corollary.}
The same computation run on $\vek\Pi$ alone gives the referential
counterpart of \eqr{6.42}. Componentwise,
\[
\begin{aligned}
   \int_{P}\rho_{\kappa}\Pi_{A}\,\dl V
   &=\int_{P}\rho_{\kappa}\big(X_{A}-\bar X_{A}\big)\,\dl V\\
   &=\int_{P}\rho_{\kappa}X_{A}\,\dl V
    -\bar X_{A}\int_{P}\rho_{\kappa}\,\dl V\\
   &=m(S)\bar X_{A}-\bar X_{A}\,m(S)=0 ,
\end{aligned}
\]
so that
\[
   \int_{P}\rho_{\kappa}\vek\Pi\,\dl V=\bze .
\]
This is the identity \eqr{6.67} uses, and it is the reason the inertia
tensor appearing in \eqr{6.64} is the one about the mass center rather
than about some other point: about any other point the integral would
not vanish and \eqr{6.45} would retain its cross terms.

As a check that the two statements agree, apply $\bQ$ to the last
display and use $\vek\pi=\bQ\vek\Pi$ together with
$\rho\,\dl v=\rho_{\kappa}\,\dl V$:
\[
   \int_{P_{t}}\rho\vek\pi\,\dl v
   =\int_{P}\rho_{\kappa}\bQ\vek\Pi\,\dl V
   =\bQ\int_{P}\rho_{\kappa}\vek\Pi\,\dl V=\bQ\bze=\bze ,
\]
which is \eqr{6.42}.

\medskip
\noindent\emph{Where rigidity entered.}
It is worth being explicit, since the calculation looks as though it
might not have needed it. The step
\[
   \int_{P}\rho_{\kappa}\bchi_{\kappa}(\bX,t)\,\dl V
   =\bQ\int_{P}\rho_{\kappa}\bX\,\dl V+\bc\,m(S)
\]
used that $\bchi_{\kappa}$ is \emph{affine} in $\bX$, and nothing
weaker will let the motion through the integral sign: for a general
$\bchi_{\kappa}$ the integral $\int\rho_{\kappa}\bchi_{\kappa}\,\dl V$
cannot be expressed through $\int\rho_{\kappa}\bX\,\dl V$ at all. For a
general motion the image of the referential mass center is therefore not
the spatial mass center --- bending a bar moves its centroid relative
to the material --- so \eqr{6.53} is a statement about rigid bodies and
not a general identity. The present problem is precisely the check that
it is available in the one place Section~\ref{sec:ch6rigid} uses it.
\end{solution}

%-----------------------------------------------------------------------
\begin{problem}[The co-rotational derivative in components]
Show that the co-rotational derivative defined by \eqr{6.68} is
expressible in the form
$\mathring{\vek{J}}=\dot J_{ij}\,\be^{*}_{i}\otimes\be^{*}_{j}$,
where the $J_{ij}$ are the components of $\vek{J}$ on the rotating
basis \eqr{6.58}.
\end{problem}

\begin{solution}
The vector case \eqr{6.61} said that $\mathring\bu$ is what one gets by
differentiating the components of $\bu$ on a basis painted on the body
and reassembling. A tensor differs from a vector in one respect only:
it has two legs, and both of them turn. Everything below is that
observation carried out, and the asymmetry between the two corrections
in \eqr{6.68} is manufactured in one line of it.

\medskip
\noindent\emph{Resolving on the rotating basis.}
By \eqr{6.58} the triad $\{\be^{*}_{i}\}$ is obtained from a fixed
orthonormal triad by the rotation $\bQ(t)$, so it is orthonormal at
every instant by Proposition~\ref{prop:orth}:
\[
   \ip{\be^{*}_{i}}{\be^{*}_{j}}=\dd_{ij}
   \qquad\text{for all }t .
\]
Theorem~\ref{thm:comprep} therefore resolves any spatial tensor on it,
\begin{equation}
   \vek{J}=J_{ij}\,\be^{*}_{i}\otimes\be^{*}_{j},
   \qquad J_{ij}=\ip{\be^{*}_{i}}{\vek{J}\be^{*}_{j}} ,
   \tag{$\ast_{7}$}\label{eq:ch6Jres}
\end{equation}
the second formula being the first tested on $\be^{*}_{j}$ and paired
with $\be^{*}_{i}$:
\[
   \ip{\be^{*}_{i}}{\vek{J}\be^{*}_{j}}
   =J_{kl}\,\ip{\be^{*}_{i}}{\big(\be^{*}_{k}\otimes\be^{*}_{l}\big)
     \be^{*}_{j}}
   =J_{kl}\,\dd_{lj}\,\ip{\be^{*}_{i}}{\be^{*}_{k}}
   =J_{kl}\dd_{lj}\dd_{ik}=J_{ij} .
\]
Both $J_{ij}$ and $\be^{*}_{i}$ depend on $t$, so differentiating
\eqref{eq:ch6Jres} is a product rule on three factors and produces one
term for the components and one for each leg:
\begin{equation}
   \dot{\vek{J}}=\dot J_{ij}\,\be^{*}_{i}\otimes\be^{*}_{j}
   +J_{ij}\,\dot{\be}^{*}_{i}\otimes\be^{*}_{j}
   +J_{ij}\,\be^{*}_{i}\otimes\dot{\be}^{*}_{j} .
   \tag{$\ast_{8}$}\label{eq:ch6Jdot}
\end{equation}
The first term is the answer we want. The other two are the rotation of
the frame, and \eqr{6.59}, written with the tensor rather than its
axial vector, gives $\dot{\be}^{*}_{i}=\bOm\be^{*}_{i}$ for each $i$.
What is instructive is that the two are disposed of by
\emph{different} identities.

\medskip
\noindent\emph{The first leg.}
A tensor acting on a dyad from the left goes straight onto the first
factor:
\[
   \bA(\ba\otimes\bb)=(\bA\ba)\otimes\bb .
\]
To see it, test both sides on an arbitrary $\bw\in\Vt$. The left gives
$\bA\big[(\ba\otimes\bb)\bw\big]=\bA\big[(\ip\bb\bw)\ba\big]
=(\ip\bb\bw)\,\bA\ba$, using linearity of $\bA$ and
Definition~\ref{def:tp}; the right gives
$\big[(\bA\ba)\otimes\bb\big]\bw=(\ip\bb\bw)\,\bA\ba$. They agree for
every $\bw$, so the tensors are equal. This is
Lemma~\ref{lem:dyadprod}. Applying it with $\bA=\bOm$,
$\ba=\be^{*}_{i}$, $\bb=\be^{*}_{j}$, and summing,
\[
   J_{ij}\,\dot{\be}^{*}_{i}\otimes\be^{*}_{j}
   =J_{ij}\,(\bOm\be^{*}_{i})\otimes\be^{*}_{j}
   =\bOm\Big(J_{ij}\,\be^{*}_{i}\otimes\be^{*}_{j}\Big)
   =\bOm\vek{J} .
\]

\medskip
\noindent\emph{The second leg.}
A tensor acting on a dyad from the right reaches the second factor only
through its transpose:
\[
   (\ba\otimes\bb)\bB^{t}=\ba\otimes(\bB\bb) .
\]
Again test on $\bw$. The left gives
$(\ba\otimes\bb)\big(\bB^{t}\bw\big)=\ip\bb{\bB^{t}\bw}\,\ba
=\ip{\bB\bb}{\bw}\,\ba$, the last step by the definition of the
transpose; the right gives $\big[\ba\otimes(\bB\bb)\big]\bw
=\ip{\bB\bb}\bw\,\ba$. They agree. Now take $\bB=\bOm$. Since
$\bOm$ is skew, $\bOm^{t}=-\bOm$, and this is where the sign is made:
\[
   J_{ij}\,\be^{*}_{i}\otimes\dot{\be}^{*}_{j}
   =J_{ij}\,\be^{*}_{i}\otimes(\bOm\be^{*}_{j})
   =\Big(J_{ij}\,\be^{*}_{i}\otimes\be^{*}_{j}\Big)\bOm^{t}
   =\vek{J}\bOm^{t}
   =-\vek{J}\bOm .
\]

\medskip
\noindent\emph{Assembling.}
Substituting both into \eqref{eq:ch6Jdot},
\[
   \dot{\vek{J}}=\dot J_{ij}\,\be^{*}_{i}\otimes\be^{*}_{j}
   +\bOm\vek{J}-\vek{J}\bOm ,
\]
so that, by the definition \eqr{6.68},
\[
   \boxed{\ \mathring{\vek{J}}
   =\dot{\vek{J}}-\bOm\vek{J}+\vek{J}\bOm
   =\dot J_{ij}\,\be^{*}_{i}\otimes\be^{*}_{j} .\ }
\]

Read backwards, this says that \eqr{6.68} was never a choice. The two
correction terms are not selected for elegance or symmetry; they are
exactly what must be removed to strip the rotation from the two legs,
and the opposite signs are the arithmetic of the two paragraphs above
and nothing else. One leg is acted on from the left, where $\bOm$
arrives unchanged; the other from the right, where it arrives
transposed, and $\bOm^{t}=-\bOm$ does the rest. A vector has one leg
and needs one correction, \eqr{6.60}; a tensor of order $n$ would need
$n$, alternating in sign by the same argument.

\medskip
\noindent\emph{What vanishing means.}
By the boxed formula,
\[
   \mathring{\vek{J}}=\Ob
   \iff
   \dot J_{ij}\,\be^{*}_{i}\otimes\be^{*}_{j}=\Ob
   \iff
   \dot J_{ij}=0\ \text{ for all }i,j ,
\]
the last equivalence because $\{\be^{*}_{i}\otimes\be^{*}_{j}\}$ is a
basis of the space of spatial tensors, being nine linearly independent
dyads built from an orthonormal triad. So $\mathring{\vek{J}}$ vanishes
precisely when $\vek{J}$ has constant components on a basis painted on
the body. That is the exact content of \eqr{6.66}, and it is what makes
Euler's equations \eqr{6.70} a system with constant coefficients when
written in the body frame --- the one frame in which the inertia tensor
stops moving. In the inertial frame $\vek{J}$ is emphatically not
constant, and $\dot{\vek{J}}=\bOm\vek{J}-\vek{J}\bOm$ is precisely how
much it moves; Remark~\ref{rem:ch6gyro} traces the gyroscopic term of
\eqr{6.70} to this and nothing else.
\end{solution}

%-----------------------------------------------------------------------
\begin{problem}[The inertia tensor is positive definite]
Recall the inertia tensor $\vek{J}(S,t)=\int_{P_{t}}\rho[(\ip{\vek\pi}
{\vek\pi})\I-\vek\pi\otimes\vek\pi]\,\dl v$ associated with an
arbitrary part $S$ of a body, where $\vek\pi$ is the position of a
point of $P_{t}$ relative to the center of mass.
\textup{(a)} Show that the integrand is positive semi-definite, a
symmetric tensor $\bA$ being positive semi-definite if and only if
$\ip\ba{\bA\ba}\geq0$ for all $\ba\in\Vt$.
\textup{(b)} Show that $\vek{J}$ is positive definite, that is
$\ip\ba{\vek{J}\ba}>0$ for all $\ba\neq\bze$.
\end{problem}

\begin{solution}
Both parts turn on Cauchy--Schwarz, and the whole interest of the
problem is in its equality case: part (a) needs the inequality, part
(b) needs to know exactly when it is tight.

\medskip
\noindent\textbf{(a)}
Write $\bA=(\ip{\vek\pi}{\vek\pi})\I-\vek\pi\otimes\vek\pi$ for the
integrand and take any $\ba\in\Vt$. Evaluating the two pieces on $\ba$
by Definition~\ref{def:tp},
\[
   \big[(\ip{\vek\pi}{\vek\pi})\I\big]\ba
   =(\ip{\vek\pi}{\vek\pi})\,\ba ,
   \qquad
   \big(\vek\pi\otimes\vek\pi\big)\ba
   =(\ip{\vek\pi}\ba)\,\vek\pi ,
\]
so that
\[
   \bA\ba=(\ip{\vek\pi}{\vek\pi})\,\ba-(\ip{\vek\pi}\ba)\,\vek\pi ,
\]
and pairing with $\ba$,
\[
   \ip\ba{\bA\ba}
   =(\ip{\vek\pi}{\vek\pi})\,\ip\ba\ba
    -(\ip{\vek\pi}\ba)\,\ip\ba{\vek\pi}
   =|\vek\pi|^{2}|\ba|^{2}-\big(\ip{\vek\pi}\ba\big)^{2},
\]
using $\ip\ba\ba=|\ba|^{2}$ and the symmetry
$\ip\ba{\vek\pi}=\ip{\vek\pi}\ba$ of the inner product. In components,
with $\pi_{i}$ and $a_{i}$ the components on a fixed basis,
\[
   \ip\ba{\bA\ba}
   =\big(\pi_{k}\pi_{k}\big)\big(a_{l}a_{l}\big)
   -\big(\pi_{k}a_{k}\big)\big(\pi_{l}a_{l}\big) .
\]
By Cauchy--Schwarz, Theorem~\ref{thm:cs},
$|\ip{\vek\pi}\ba|\leq|\vek\pi|\,|\ba|$; both sides are non-negative,
so squaring preserves the inequality and
$(\ip{\vek\pi}\ba)^{2}\leq|\vek\pi|^{2}|\ba|^{2}$. Hence
\[
   \boxed{\ \ip\ba{\bA\ba}\geq0\ \text{ for every }\ba .\ }
\]
Symmetry of $\bA$ is immediate: $\I^{t}=\I$, and
$(\vek\pi\otimes\vek\pi)^{t}=\vek\pi\otimes\vek\pi$ since
$(\ba\otimes\bb)^{t}=\bb\otimes\ba$. So $\bA$ is symmetric and
positive semi-definite in the stated sense.

The geometry behind the algebra is worth naming, because it makes part
(b) transparent. For $|\ba|=1$ the expression is
$|\vek\pi|^{2}-(\ip{\vek\pi}\ba)^{2}$: the squared length of
$\vek\pi$, less the square of its component along $\ba$. Decomposing
$\vek\pi=(\ip{\vek\pi}\ba)\ba+\vek\pi^{\perp}$ with
$\ip{\vek\pi^{\perp}}\ba=0$ and applying Pythagoras,
\[
   |\vek\pi|^{2}=\big(\ip{\vek\pi}\ba\big)^{2}+|\vek\pi^{\perp}|^{2},
   \qquad\text{so}\qquad
   \ip\ba{\bA\ba}=|\vek\pi^{\perp}|^{2} ,
\]
the squared distance from the point $\vek\pi$ to the axis through the
mass center in the direction $\ba$. Hence $\ip\ba{\vek{J}\ba}$ is the
moment of inertia about that axis, and (a) says that a moment of
inertia is never negative --- a fact one would want to be true before
proving it.

The same decomposition locates the equality. Theorem~\ref{thm:cs} is
tight exactly when $\vek\pi$ and $\ba$ are linearly dependent, that is
when $\vek\pi^{\perp}=\bze$, so for fixed $\ba\neq\bze$ the integrand
vanishes at precisely those points where $\vek\pi$ is parallel to
$\ba$, the case $\vek\pi=\bze$ included. Those points form a single
straight line through the mass center, namely
$\{\vek\pi=\lambda\ba:\lambda\in\Rn\}$.

\medskip
\noindent\textbf{(b)}
Integrate. For any fixed $\ba$,
\[
\begin{aligned}
   \ip\ba{\vek{J}\ba}
   &=\ip\ba{\Big(\int_{P_{t}}\rho\bA\,\dl v\Big)\ba}
    =\int_{P_{t}}\rho\,\ip\ba{\bA\ba}\,\dl v\\
   &=\int_{P_{t}}\rho\Big[|\vek\pi|^{2}|\ba|^{2}
     -\big(\ip{\vek\pi}\ba\big)^{2}\Big]\,\dl v ,
\end{aligned}
\]
the interchange of $\ba$ with the integral being legitimate because
$\ba$ is a fixed vector and the integral is of components on a fixed
basis, in the sense of Remark~\ref{rem:symbolic}: explicitly,
$a_{i}J_{ij}a_{j}=a_{i}\big(\int_{P_{t}}\rho A_{ij}\,\dl v\big)a_{j}
=\int_{P_{t}}\rho\,a_{i}A_{ij}a_{j}\,\dl v$, since the $a_{i}$ are
constants. The integrand is non-negative by (a) and by $\rho>0$, so
$\ip\ba{\vek{J}\ba}\geq0$ for every $\ba$. It remains to exclude
equality when $\ba\neq\bze$.

Suppose $\ip\ba{\vek{J}\ba}=0$ for some $\ba\neq\bze$. The integrand
$f=\rho[\,|\vek\pi|^{2}|\ba|^{2}-(\ip{\vek\pi}\ba)^{2}]$ is continuous
on $P_{t}$ and non-negative, and its integral vanishes. A continuous
non-negative function with vanishing integral over a region vanishes
identically on it: were $f(\bx_{0})>0$ at some point, continuity would
give a ball $B\subset P_{t}$ about $\bx_{0}$ on which
$f>\tfrac12f(\bx_{0})>0$, whence
\[
   \int_{P_{t}}f\,\dl v\ \geq\ \int_{B}f\,\dl v
   \ \geq\ \tfrac12f(\bx_{0})\,V(B)\ >\ 0 ,
\]
the first inequality because $f\geq0$ off $B$, contradicting the
hypothesis. This is Theorem~\ref{thm:localization} run with an
inequality in place of an equation.

So $f\equiv0$ on $P_{t}$, and since $\rho>0$ the bracket vanishes
throughout. By the equality case recorded in (a), every point of
$P_{t}$ then has $\vek\pi$ parallel to $\ba$: the whole of $P_{t}$ lies
on one straight line through the mass center. A line has zero volume
and $P_{t}$ has positive volume, $S$ being a genuine part of a body, so
the supposition is untenable and
\[
   \boxed{\ \ip\ba{\vek{J}\ba}>0\ \text{ for every }\ba\neq\bze .\ }
\]

\medskip
\noindent\emph{What positive definiteness buys.}
It is not an ornament. By Theorem~\ref{thm:pdiff} it gives $\vek{J}$
three strictly positive principal values $J_{1},J_{2},J_{3}$, the
principal moments of inertia, and by Theorem~\ref{thm:spectral} an
orthonormal triad of principal axes on which
$\vek{J}=\sum_{i}J_{i}\,\bn_{i}\otimes\bn_{i}$. It also makes $\vek{J}$
invertible, since a tensor with $\ip\ba{\vek{J}\ba}>0$ for all
$\ba\neq\bze$ has trivial kernel: $\vek{J}\ba=\bze$ would give
$\ip\ba{\vek{J}\ba}=0$ and hence $\ba=\bze$. That is what allows
\eqr{6.70} to be solved for the angular acceleration,
\[
   \dot\bom=\vek{J}^{-1}\big[\bm(S,t)-\bom\times\vek{J}\bom\big] ,
   \qquad
   \vek{J}^{-1}=\sum_{i}J_{i}^{-1}\,\bn_{i}\otimes\bn_{i} .
\]
Without it the rotational equation could not be put in normal form and
the initial-value problem for a rigid body would not be well posed. The
degenerate case excluded above, a body concentrated on a line, is
exactly the case of a body with no resistance to rotation about that
line --- the moment of inertia about it would be zero --- and it is
excluded not by an extra hypothesis but by the continuum hypothesis
itself.
\end{solution}

%-----------------------------------------------------------------------
\begin{problem}[Shear stress on an arbitrary plane; the octahedral
shear]
Let $\{\sigma_{i}\}$ and $\{\bn_{i}\}$ be the principal Cauchy stresses
and the associated orthonormal principal stress axes at a position
$\bx$ in a body. Use these axes to define the unit vector
$\bn(\theta,\varphi)=\cos\varphi\,\ber(\theta)+\sin\varphi\,\bn_{3}$,
where $\ber(\theta)=\cos\theta\,\bn_{1}+\sin\theta\,\bn_{2}$, with
$0\leq\theta<2\pi$ and $-\pi\leq\varphi\leq\pi$.
\textup{(a)} Compute the shear stress $\tau(\theta,\varphi)$ on the
plane with unit normal $\bn(\theta,\varphi)$.
\textup{(b)} Compute the octahedral shear stress, that is the shear
stress on a plane whose normal is inclined at equal angles to each of
the principal axes.
\end{problem}

\begin{solution}
\noindent\emph{Setting up.}
By the spectral theorem, Theorem~\ref{thm:spectral}, applied to the
symmetric $\bT$ of \eqr{6.117},
\[
   \bT=\sum_{i=1}^{3}\sigma_{i}\,\bn_{i}\otimes\bn_{i} ,
   \qquad
   \ip{\bn_{i}}{\bn_{j}}=\dd_{ij} .
\]
Write $\bn=\nu_{i}\bn_{i}$. From the parametrisation given,
$\ber(\theta)=\cos\theta\,\bn_{1}+\sin\theta\,\bn_{2}$ and
$\bn=\cos\varphi\,\ber+\sin\varphi\,\bn_{3}$, so
\[
   \nu_{1}=\cos\varphi\cos\theta,\qquad
   \nu_{2}=\cos\varphi\sin\theta,\qquad
   \nu_{3}=\sin\varphi ,
\]
and
\[
   \nu_{i}\nu_{i}
   =\cos^{2}\varphi\cos^{2}\theta+\cos^{2}\varphi\sin^{2}\theta
     +\sin^{2}\varphi
   =\cos^{2}\varphi\big(\cos^{2}\theta+\sin^{2}\theta\big)
     +\sin^{2}\varphi
   =1 ,
\]
confirming that $\bn$ is a unit vector for every $(\theta,\varphi)$.
The traction is diagonal in the principal frame:
\[
   \bt=\bT\bn
   =\sum_{i}\sigma_{i}\big(\bn_{i}\otimes\bn_{i}\big)\bn
   =\sum_{i}\sigma_{i}\,\ip{\bn_{i}}\bn\,\bn_{i}
   =\sum_{i}\sigma_{i}\nu_{i}\,\bn_{i},
\]
and orthonormality gives the two scalars everything will be built from,
\[
   \sigma=\ip\bn\bt
   =\Big(\sum_{j}\nu_{j}\bn_{j}\Big)\cdot
    \Big(\sum_{i}\sigma_{i}\nu_{i}\bn_{i}\Big)
   =\sum_{i}\sigma_{i}\nu_{i}^{2},
   \qquad
   |\bt|^{2}=\ip\bt\bt=\sum_{i}\sigma_{i}^{2}\nu_{i}^{2}.
\]

\medskip
\noindent\textbf{(a)}
The decomposition \eqr{6.122} writes $\bt=\sigma\bn+\tau\bs$ with
$\ip\bn\bs=0$ and $|\bs|=|\bn|=1$, so the two pieces are orthogonal and
\[
   |\bt|^{2}=\ip{\sigma\bn+\tau\bs}{\sigma\bn+\tau\bs}
   =\sigma^{2}|\bn|^{2}+2\sigma\tau\underbrace{\ip\bn\bs}_{=\,0}
     +\tau^{2}|\bs|^{2}
   =\sigma^{2}+\tau^{2} ,
\]
whence
\[
   \tau^{2}=|\bt|^{2}-\sigma^{2}
   =\sum_{i}\sigma_{i}^{2}\nu_{i}^{2}
    -\Big(\sum_{i}\sigma_{i}\nu_{i}^{2}\Big)^{2}.
\]
This answers (a), and conceals what it means. Put $w_{i}=\nu_{i}^{2}$;
then $w_{i}\geq0$ and $\sum_{i}w_{i}=1$, so $w$ is a probability
distribution on the three principal stresses, and the right-hand side
is $\langle\sigma^{2}\rangle-\langle\sigma\rangle^{2}$: the shear
stress squared is the \emph{variance} of the principal stresses under
the weights $w$. A variance is a sum over pairs, and writing it that
way costs four lines:
\[
\begin{aligned}
   \sum_{i<j}w_{i}w_{j}\big(\sigma_{i}-\sigma_{j}\big)^{2}
   &=\sum_{i<j}w_{i}w_{j}\big(\sigma_{i}^{2}-2\sigma_{i}\sigma_{j}
     +\sigma_{j}^{2}\big)\\
   &=\sum_{i<j}w_{i}w_{j}\big(\sigma_{i}^{2}+\sigma_{j}^{2}\big)
     -2\sum_{i<j}w_{i}w_{j}\sigma_{i}\sigma_{j}\\
   &=\sum_{i}\sigma_{i}^{2}w_{i}\Big(\sum_{j\neq i}w_{j}\Big)
     -2\sum_{i<j}w_{i}w_{j}\sigma_{i}\sigma_{j}\\
   &=\sum_{i}\sigma_{i}^{2}w_{i}(1-w_{i})
     -2\sum_{i<j}w_{i}w_{j}\sigma_{i}\sigma_{j}\\
   &=\sum_{i}\sigma_{i}^{2}w_{i}
     -\Big(\sum_{i}\sigma_{i}^{2}w_{i}^{2}
       +2\sum_{i<j}w_{i}w_{j}\sigma_{i}\sigma_{j}\Big)\\
   &=\sum_{i}\sigma_{i}^{2}w_{i}
     -\Big(\sum_{i}\sigma_{i}w_{i}\Big)^{2} .
\end{aligned}
\]
The third line collected, for each fixed $i$, the two places where
$\sigma_{i}^{2}$ occurs in the pair sum; the fourth used
$\sum_{j\neq i}w_{j}=1-w_{i}$; and the last recognised
$\big(\sum_{i}\sigma_{i}w_{i}\big)^{2}
=\sum_{i}\sigma_{i}^{2}w_{i}^{2}
+2\sum_{i<j}w_{i}w_{j}\sigma_{i}\sigma_{j}$. Therefore
\[
   \boxed{\ \tau^{2}(\theta,\varphi)
   =\sum_{i<j}\big(\sigma_{i}-\sigma_{j}\big)^{2}
   \nu_{i}^{2}\nu_{j}^{2}\ }
\]
with the $\nu_{i}$ above. Written out, the three pairs are
\[
\begin{aligned}
   \nu_{1}^{2}\nu_{2}^{2}
     &=\cos^{4}\varphi\cos^{2}\theta\sin^{2}\theta
      =\tfrac14\cos^{4}\varphi\,\sin^{2}2\theta,\\
   \nu_{1}^{2}\nu_{3}^{2}
     &=\cos^{2}\varphi\cos^{2}\theta\sin^{2}\varphi
      =\tfrac14\sin^{2}2\varphi\,\cos^{2}\theta,\\
   \nu_{2}^{2}\nu_{3}^{2}
     &=\cos^{2}\varphi\sin^{2}\theta\sin^{2}\varphi
      =\tfrac14\sin^{2}2\varphi\,\sin^{2}\theta,
\end{aligned}
\]
using $2\cos\theta\sin\theta=\sin2\theta$ and
$2\cos\varphi\sin\varphi=\sin2\varphi$, so that
\[
\begin{aligned}
   \tau^{2}=\tfrac14\Big[&(\sigma_{1}-\sigma_{2})^{2}
     \cos^{4}\varphi\,\sin^{2}2\theta\\
   &+\sin^{2}2\varphi\Big((\sigma_{1}-\sigma_{3})^{2}\cos^{2}\theta
     +(\sigma_{2}-\sigma_{3})^{2}\sin^{2}\theta\Big)\Big].
\end{aligned}
\]
The pair form earns its keep immediately. Shear is driven by
\emph{differences} of principal stresses, so if all three coincide it
vanishes on every plane, which is the state of pure pressure
\eqr{6.128}; and if $\bn$ is a principal direction, two of the $\nu_{i}$
vanish, killing every pair, which is \eqr{6.127} read backwards.

\medskip
\noindent\emph{The largest value.}
This is worth deriving rather than quoting, and the derivation is an
exercise in choosing coordinates. Order the principal stresses
$\sigma_{1}\geq\sigma_{2}\geq\sigma_{3}$ and name the two gaps
\[
   a=\sigma_{1}-\sigma_{2}\geq0,
   \qquad
   b=\sigma_{2}-\sigma_{3}\geq0,
   \qquad
   \sigma_{1}-\sigma_{3}=a+b ,
\]
so that in the weights the boxed formula reads
\begin{equation}
   \tau^{2}=a^{2}w_{1}w_{2}+(a+b)^{2}w_{1}w_{3}+b^{2}w_{2}w_{3},
   \qquad w_{i}\geq0,\quad \textstyle\sum_{i}w_{i}=1 .
   \tag{$\sharp$}\label{eq:ch6shearw}
\end{equation}
Assume $a+b>0$; the case $a+b=0$ is the pure pressure already disposed
of. Maximizing a quadratic form on a simplex invites Lagrange
multipliers, but the constraint can be eliminated instead, and then the
answer is visible. Hold $w_{2}$ fixed, put $s=1-w_{2}=w_{1}+w_{3}$, and
split the two extreme weights symmetrically about their mean,
\[
   w_{1}=\tfrac{s}{2}+u,\qquad w_{3}=\tfrac{s}{2}-u,
   \qquad |u|\leq\tfrac{s}{2},
\]
so that $u$ measures how unevenly the extremes are loaded. Then
$w_{1}w_{2}=w_{2}(\tfrac{s}{2}+u)$,
$w_{2}w_{3}=w_{2}(\tfrac{s}{2}-u)$ and
$w_{1}w_{3}=\tfrac{s^{2}}{4}-u^{2}$, so
\[
\begin{aligned}
   \tau^{2}&=w_{2}\Big[a^{2}\big(\tfrac{s}{2}+u\big)
     +b^{2}\big(\tfrac{s}{2}-u\big)\Big]
     +(a+b)^{2}\Big(\tfrac{s^{2}}{4}-u^{2}\Big)\\
   &=-(a+b)^{2}u^{2}+w_{2}\big(a^{2}-b^{2}\big)u
     +\tfrac{s}{2}w_{2}\big(a^{2}+b^{2}\big)
     +\tfrac{s^{2}}{4}(a+b)^{2} ,
\end{aligned}
\]
a downward parabola in $u$, which one maximizes by completing the
square. With $a^{2}-b^{2}=(a+b)(a-b)$ the vertex sits at
\[
   u^{*}=\frac{w_{2}(a^{2}-b^{2})}{2(a+b)^{2}}
   =\frac{w_{2}(a+b)(a-b)}{2(a+b)^{2}}
   =\frac{(a-b)\,w_{2}}{2(a+b)} ,
\]
and completing the square gives
\begin{equation}
   \tau^{2}=\frac{(a+b)^{2}}{4}
   -ab\,w_{2}
   -(a+b)^{2}\big(u-u^{*}\big)^{2} .
   \tag{$\sharp\sharp$}\label{eq:ch6shearmax}
\end{equation}

% reader's note: the check that (##) is an identity
\rnote{ch6:prob5-identity}

Both corrections in \eqref{eq:ch6shearmax} are non-positive, the first
because $a,b,w_{2}\geq0$ and the second because it is a square times a
positive coefficient, so
\[
   \boxed{\ \tau^{2}\leq\tfrac14\big(\sigma_{1}-\sigma_{3}\big)^{2},
   \qquad
   \tau_{\max}=\tfrac12\big(\sigma_{1}-\sigma_{3}\big),\ }
\]
with equality only if both vanish, that is $ab\,w_{2}=0$ and $u=u^{*}$.
For distinct principal stresses $a>0$ and $b>0$, so $ab>0$ forces
$w_{2}=0$; then $u^{*}=0$ and the second condition gives $u=0$, whence
$s=1$ and
\[
   w=\big(\tfrac12,\,0,\,\tfrac12\big),
   \qquad
   \nu_{1}=\pm\tfrac{\sqrt2}{2},\ \ \nu_{2}=0,\ \
   \nu_{3}=\pm\tfrac{\sqrt2}{2},
   \qquad
   \bn=\tfrac{\sqrt2}{2}\big(\bn_{1}\pm\bn_{3}\big) :
\]
the plane bisecting the directions of greatest and least principal
stress, with the intermediate stress carrying no weight at all. Two
coincident principal stresses make $ab=0$ and leave $w_{2}$ free, so
the maximizer stops being unique --- the degeneracy Chapter 3 met
whenever a principal value repeated.

\medskip
\noindent\textbf{(b)}
A normal equally inclined to the three axes has
$\nu_{1}^{2}=\nu_{2}^{2}=\nu_{3}^{2}$, and these sum to one, so each is
$\tfrac13$ and every pair in the boxed formula enters with weight
$\tfrac13\cdot\tfrac13=\tfrac19$:
\[
   \tau_{\text{oct}}^{2}
   =\tfrac19\Big[(\sigma_{1}-\sigma_{2})^{2}
   +(\sigma_{2}-\sigma_{3})^{2}+(\sigma_{3}-\sigma_{1})^{2}\Big],
\]
that is
\[
   \boxed{\ \tau_{\text{oct}}
   =\tfrac13\sqrt{(\sigma_{1}-\sigma_{2})^{2}
   +(\sigma_{2}-\sigma_{3})^{2}+(\sigma_{3}-\sigma_{1})^{2}} ,\ }
\]
with corresponding normal stress
$\sigma_{\text{oct}}=\sum_{i}\sigma_{i}\cdot\tfrac13
=\tfrac13\tr\bT$. In terms of the parametrisation,
$\nu_{3}^{2}=\sin^{2}\varphi=\tfrac13$ and
$\nu_{1}^{2}=\cos^{2}\varphi\cos^{2}\theta=\tfrac13$; since
$\cos^{2}\varphi=1-\tfrac13=\tfrac23$, the second gives
$\cos^{2}\theta=\tfrac13\cdot\tfrac32=\tfrac12$. So
$\varphi=\pm\arcsin(1/\sqrt3)$ and
$\theta\in\{\pi/4,3\pi/4,5\pi/4,7\pi/4\}$: eight normals in all, one
into each octant, being the face normals of a regular octahedron
centred at $\bx$, which is where the name comes from.

\medskip
\noindent\emph{Connection with the deviatoric stress.}
The octahedral shear turns out to be an invariant in disguise, and this
is why it, rather than the shear on some particular plane, is what
plasticity uses. Let $\bT'=\bT-\tfrac13(\tr\bT)\I$ be the deviatoric
part, the splitting of Problem 15 of Chapter 1, with principal values
$\sigma_{i}-\bar\sigma$ and $\bar\sigma=\tfrac13\tr\bT
=\tfrac13\sum_{i}\sigma_{i}$. By Theorem~\ref{thm:tip} the squared norm
is the sum of squares of the principal values, so
\[
\begin{aligned}
   |\bT'|^{2}&=\sum_{i}(\sigma_{i}-\bar\sigma)^{2}
   =\sum_{i}\sigma_{i}^{2}-2\bar\sigma\sum_{i}\sigma_{i}
     +3\bar\sigma^{2}\\
   &=\sum_{i}\sigma_{i}^{2}-2\bar\sigma(3\bar\sigma)+3\bar\sigma^{2}
   =\sum_{i}\sigma_{i}^{2}-3\bar\sigma^{2}\\
   &=\sum_{i}\sigma_{i}^{2}
     -\tfrac13\Big(\sum_{i}\sigma_{i}\Big)^{2}
   =\tfrac13\Big[3\sum_{i}\sigma_{i}^{2}
     -\Big(\sum_{i}\sigma_{i}\Big)^{2}\Big]
   =\tfrac13\sum_{i<j}(\sigma_{i}-\sigma_{j})^{2},
\end{aligned}
\]
the last equality being the identity
$\sum_{i<j}(\sigma_{i}-\sigma_{j})^{2}
=3\sum_{i}\sigma_{i}^{2}-\big(\sum_{i}\sigma_{i}\big)^{2}$, which one
checks by expanding the left side as
$2\sum_{i}\sigma_{i}^{2}-2\sum_{i<j}\sigma_{i}\sigma_{j}$ and the right
as $3\sum_{i}\sigma_{i}^{2}-\sum_{i}\sigma_{i}^{2}
-2\sum_{i<j}\sigma_{i}\sigma_{j}$. Comparing with the boxed result,
\[
   \tau_{\text{oct}}^{2}=\tfrac19\sum_{i<j}(\sigma_{i}-\sigma_{j})^{2}
   =\tfrac19\cdot3\,|\bT'|^{2}=\tfrac13|\bT'|^{2},
   \qquad\text{so}\qquad
   \tau_{\text{oct}}=\frac{|\bT'|}{\sqrt3} .
\]
A fixed multiple of $|\bT'|$ depends on $\bT$ alone and not on the
frame, and is unchanged by superposed pressure: replacing $\bT$ by
$\bT+p\I$ moves $\tr\bT$ by $3p$ and leaves
$\bT'=\bT+p\I-\tfrac13(\tr\bT+3p)\I$ exactly where it was. Those are
precisely the two properties one demands of a yield criterion, and they
are why von Mises is stated in terms of this quantity.
\end{solution}

%-----------------------------------------------------------------------
\begin{problem}[Piola tractions and the symmetry restriction]
A body undergoes a homogeneous deformation whose deformation gradient
at a given time $t$ has components
\[
   (F_{iA})=\begin{pmatrix}2&3&0\\1&2&0\\0&0&1\end{pmatrix}
\]
with respect to the fixed basis $\{\be_{i}\otimes\hbE_{A}\}$. The Piola
tractions $\bp(\bN)$ acting on the deformed surfaces whose unit normals
in the reference configuration are $\hbE_{A}$ are
\[
   \bp(\hbE_{1})=a\be_{1}+10\be_{2},\qquad
   \bp(\hbE_{2})=(b-6)\be_{3}+3\be_{2},\qquad
   \bp(\hbE_{3})=c\,\be_{2}+3\be_{3},
\]
where $a$, $b$ and $c$ are constants.
\textup{(a)} Use $\bp(\bN)=\bP\bN$ together with
$\hI=\hbE_{A}\otimes\hbE_{A}$ to show that
$\bP=\bp_{A}\otimes\hbE_{A}$, where $\bp_{A}=\bp(\hbE_{A})$.
\textup{(b)} Find the components $P_{iA}$ as functions of $a$, $b$,
$c$.
\textup{(c)} Apply the rotational momentum balance $\bT=\bT^{t}$ to
determine $a$, $b$, $c$.
\textup{(d)} Calculate the components $T_{ij}$ of the Cauchy stress.
\end{problem}

\begin{solution}
The problem gives three tractions and asks for a tensor, and the bridge
between them is the observation that a tensor is determined by what it
does to a basis.

\medskip
\noindent\textbf{(a)}
The referential identity resolves on the referential basis,
$\hI=\hbE_{A}\otimes\hbE_{A}$, which is Theorem~\ref{thm:comprep}
written in $\hVt$; one checks it by testing on $\hbE_{B}$:
\[
   \big(\hbE_{A}\otimes\hbE_{A}\big)\hbE_{B}
   =\ip{\hbE_{A}}{\hbE_{B}}\,\hbE_{A}
   =\dd_{AB}\hbE_{A}=\hbE_{B} .
\]
Post-multiplying $\bP$ by it changes nothing, and the dyad rule
$\bA(\ba\otimes\bb)=(\bA\ba)\otimes\bb$ of Lemma~\ref{lem:dyadprod}
then moves $\bP$ onto the first leg:
\[
   \bP=\bP\hI=\bP\big(\hbE_{A}\otimes\hbE_{A}\big)
   =\big(\bP\hbE_{A}\big)\otimes\hbE_{A}
   =\bp(\hbE_{A})\otimes\hbE_{A}=\bp_{A}\otimes\hbE_{A},
\]
the fourth equality by \eqr{6.132} with $\bN=\hbE_{A}$. Three tractions
therefore determine $\bP$ completely, and the identical argument run in
$\Vt$ on $\bT$, with $\I=\be_{i}\otimes\be_{i}$ and \eqr{6.94} in place
of \eqr{6.132}, gives \eqr{6.107}.

% reader's note: columns, not rows
\rnote{ch6:prob6-columns}

\medskip
\noindent\textbf{(b)}
Extracting components from (a),
\[
   P_{iA}=\ip{\be_{i}}{\bP\hbE_{A}}=\ip{\be_{i}}{\bp_{A}} ,
\]
so the $A$th \emph{column} of the matrix $(P_{iA})$ is the component
triple of $\bp_{A}$. Reading the three given vectors, and writing
$\bp_{2}$ with its terms in index order --- the problem states it as
$(b-6)\be_{3}+3\be_{2}$, and taking the coefficients in the order
printed is the commonest way of losing this problem ---
\[
   \bp_{1}=a\be_{1}+10\be_{2}+0\be_{3},
   \qquad
   \bp_{2}=0\be_{1}+3\be_{2}+(b-6)\be_{3},
   \qquad
   \bp_{3}=0\be_{1}+c\,\be_{2}+3\be_{3},
\]
that is $\bp_{1}=(a,10,0)$, $\bp_{2}=(0,3,b-6)$, $\bp_{3}=(0,c,3)$.
Stacking these as columns,
\[
   (P_{iA})=\begin{pmatrix}
     a & 0 & 0\\
     10 & 3 & c\\
     0 & b-6 & 3
   \end{pmatrix}.
\]

\medskip
\noindent\textbf{(c)}
The constants are fixed by a symmetry that $\bP$ itself cannot possess.
By \eqr{6.117} the Cauchy stress is symmetric, and by the computation
preceding \eqr{6.143},
\[
   \bP\bF^{t}=\bT\bF^{*}\bF^{t}=\bT\big(J\bF^{-t}\big)\bF^{t}
   =J\bT\big(\bF^{t}\big)^{-1}\bF^{t}=J\bT ,
\]
with $J>0$. So requiring $\bT=\bT^{t}$ is requiring
$\bM:=\bP\bF^{t}$ to be symmetric, which is \eqr{6.143} in the form
$P_{iA}F_{jA}=F_{iA}P_{jA}$. The virtue of this route is that it never
needs $\bT$: it works entirely with the data given.

Since $\big(\bF^{t}\big)_{Aj}=F_{jA}$,
\[
   M_{ij}=P_{iA}\big(\bF^{t}\big)_{Aj}=P_{iA}F_{jA} ,
\]
which pairs the $i$th \emph{row} of $(P_{iA})$ with the $j$th
\emph{row} of $(F_{iA})$. The rows are
\[
\begin{aligned}
   P_{1\cdot}&=(a,\,0,\,0), &\qquad P_{2\cdot}&=(10,\,3,\,c),
   &\qquad P_{3\cdot}&=(0,\,b-6,\,3),\\
   F_{1\cdot}&=(2,\,3,\,0), &\qquad F_{2\cdot}&=(1,\,2,\,0),
   &\qquad F_{3\cdot}&=(0,\,0,\,1),
\end{aligned}
\]
and all nine entries are
\[
\begin{aligned}
   M_{11}&=a\cdot2+0\cdot3+0\cdot0=2a, &
   M_{12}&=a\cdot1+0\cdot2+0\cdot0=a,\\
   M_{13}&=a\cdot0+0\cdot0+0\cdot1=0, &
   M_{21}&=10\cdot2+3\cdot3+c\cdot0=29,\\
   M_{22}&=10\cdot1+3\cdot2+c\cdot0=16, &
   M_{23}&=10\cdot0+3\cdot0+c\cdot1=c,\\
   M_{31}&=0\cdot2+(b-6)3=3(b-6), &
   M_{32}&=0\cdot1+(b-6)2=2(b-6),\\
   M_{33}&=(b-6)\cdot0+3\cdot1=3 .
\end{aligned}
\]
Symmetry is three conditions, one per off-diagonal pair, and they
resolve in cascade rather than simultaneously:
\[
\begin{aligned}
   M_{12}=M_{21}&:\quad a=29 ,\\
   M_{13}=M_{31}&:\quad 0=3(b-6)\ \Longrightarrow\ b=6 ,\\
   M_{23}=M_{32}&:\quad c=2(b-6)=2\cdot0=0 .
\end{aligned}
\]
The last two are consistent --- both give $b=6$ --- which is a check
that the data were not over-determined. Thus
\[
   \boxed{\ a=29,\qquad b=6,\qquad c=0 .\ }
\]

\medskip
\noindent\textbf{(d)}
First the Jacobian. Expanding $\det\bF$ along the third row, whose only
non-zero entry is $F_{33}=1$,
\[
   J=\det\bF
   =1\cdot\det\begin{pmatrix}2&3\\1&2\end{pmatrix}
   =1\cdot(2\cdot2-3\cdot1)=4-3=1 .
\]
So $\bT=J^{-1}\bP\bF^{t}=\bM$, and with $a=29$, $b=6$, $c=0$ the nine
entries computed above become
\[
\begin{aligned}
   M_{11}&=2\cdot29=58, & M_{12}&=29, & M_{13}&=0,\\
   M_{21}&=29, & M_{22}&=16, & M_{23}&=0,\\
   M_{31}&=3\cdot0=0, & M_{32}&=2\cdot0=0, & M_{33}&=3 ,
\end{aligned}
\]
that is
\[
   \boxed{\ (T_{ij})=\begin{pmatrix}
     58 & 29 & 0\\
     29 & 16 & 0\\
     0 & 0 & 3
   \end{pmatrix}. \ }
\]
It is symmetric, as it had to be, and its third row and column are
uncoupled from the $1$--$2$ plane, which is what one expects of a
deformation leaving $X_{3}$ alone.

\medskip
\noindent\emph{Check by inversion.}
Inverting the relation is worthwhile and makes one point that is easy
to get wrong. With $J=1$, \eqr{6.140} and \eqr{2.55} give
$\bP=\bT\bF^{*}=J\bT\bF^{-t}=\bT\bF^{-t}$. Computing $\bF^{-1}$ from
the adjugate, or by solving $\bF\bF^{-1}=\I$ directly: the upper-left
$2\times2$ block has determinant $1$, so its inverse is
$\begin{pmatrix}2&-3\\-1&2\end{pmatrix}$, and the third row and column
are trivial. Hence
\[
   \bF^{-1}=\begin{pmatrix}2&-3&0\\-1&2&0\\0&0&1\end{pmatrix},
   \qquad
   \bF^{-t}=\begin{pmatrix}2&-1&0\\-3&2&0\\0&0&1\end{pmatrix},
\]
which one verifies:
$\bF\bF^{-1}$ has $(1,1)$ entry $2\cdot2+3\cdot(-1)=1$, $(1,2)$ entry
$2\cdot(-3)+3\cdot2=0$, $(2,1)$ entry $1\cdot2+2\cdot(-1)=0$ and
$(2,2)$ entry $1\cdot(-3)+2\cdot2=1$. Then, entry by entry,
\[
\begin{aligned}
   \bT\bF^{-t}
   &=\begin{pmatrix}58&29&0\\29&16&0\\0&0&3\end{pmatrix}
     \begin{pmatrix}2&-1&0\\-3&2&0\\0&0&1\end{pmatrix}\\[1mm]
   &=\begin{pmatrix}
     58\cdot2+29\cdot(-3) & 58\cdot(-1)+29\cdot2 & 0\\
     29\cdot2+16\cdot(-3) & 29\cdot(-1)+16\cdot2 & 0\\
     0 & 0 & 3\end{pmatrix}
    =\begin{pmatrix}29&0&0\\10&3&0\\0&0&3\end{pmatrix},
\end{aligned}
\]
since $116-87=29$, $-58+58=0$, $58-48=10$ and $-29+32=3$. This is
$(P_{iA})$ with the values found, as it should be.

Note that $\bT\neq\bP$ even though $J=1$: the two differ by
$\bF^{-t}$, which is not the identity, and a unit Jacobian says only
that volumes are preserved, not areas. That is
Remark~\ref{rem:ch6threestress} in miniature --- the Piola stress
carries a referential leg, and neither its symmetry nor its numerical
values transfer to $\bT$ without the conversion.
\end{solution}


\begin{problem}[Three stresses for one shear state]
Consider the deformation
\[
   x_{1}=a_{1}(X_{1}+bX_{2}),\qquad x_{2}=a_{2}X_{2},\qquad
   x_{3}=a_{3}X_{3} .
\]
Suppose the Cauchy traction is $\bt=\tau\be_{1}$, $\tau\be_{2}$ and
$\bze$ on the planes with unit normals $\bn=\be_{2}$, $\be_{1}$ and
$\be_{3}$ respectively; note the crossed correspondence in the first
two. Find the stress tensors $\bT$, $\bP$ and $\bS$ in terms of $\tau$,
$a_{1}$, $a_{2}$, $a_{3}$ and $b$.
\end{problem}

\begin{solution}
\noindent\emph{The Cauchy stress.}
The data prescribe $\bT$ on a basis, which by \eqr{6.107} is $\bT$
itself: the $j$th column of $(T_{ij})$ is $\bT\be_{j}$, since
$T_{ij}=\ip{\be_{i}}{\bT\be_{j}}$. Reading the three given values,
\[
   \bT\be_{1}=\tau\be_{2}=(0,\tau,0),
   \qquad
   \bT\be_{2}=\tau\be_{1}=(\tau,0,0),
   \qquad
   \bT\be_{3}=\bze=(0,0,0),
\]
and stacking as columns,
\[
   \boxed{\;\bT=\tau\big(\be_{1}\otimes\be_{2}
   +\be_{2}\otimes\be_{1}\big),\qquad
   (T_{ij})=\begin{pmatrix}0&\tau&0\\ \tau&0&0\\ 0&0&0\end{pmatrix}.\;}
\]
That this comes out symmetric is a check on the data rather than a
derivation, and the crossed correspondence the problem warns about is
exactly what makes it work: had the problem prescribed
$\bT\be_{2}=\tau\be_{1}$ together with $\bT\be_{1}=\bze$, the matrix
would have had $T_{12}=\tau$ and $T_{21}=0$, violating \eqr{6.117}, and
no state of stress would exist.

\medskip
\noindent\emph{The deformation gradient and its inverse.}
From $F_{iA}=\partial x_{i}/\partial X_{A}$ applied to
$x_{1}=a_{1}(X_{1}+bX_{2})$, $x_{2}=a_{2}X_{2}$, $x_{3}=a_{3}X_{3}$,
\[
\begin{aligned}
   F_{11}&=a_{1}, & F_{12}&=a_{1}b, & F_{13}&=0,\\
   F_{21}&=0, & F_{22}&=a_{2}, & F_{23}&=0,\\
   F_{31}&=0, & F_{32}&=0, & F_{33}&=a_{3},
\end{aligned}
\qquad
   (F_{iA})=\begin{pmatrix}a_{1}&a_{1}b&0\\ 0&a_{2}&0\\
   0&0&a_{3}\end{pmatrix},
\]
and $J=\det\bF=a_{1}a_{2}a_{3}$, the determinant of an upper
triangular matrix being the product of its diagonal.

Inverting by back substitution: solving $\bF\bY=\I$ column by column,
the third row gives $a_{3}Y_{3j}=\dd_{3j}$, so $Y_{31}=Y_{32}=0$ and
$Y_{33}=1/a_{3}$; the second row gives $a_{2}Y_{2j}=\dd_{2j}$, so
$Y_{21}=0$, $Y_{22}=1/a_{2}$, $Y_{23}=0$; and the first row gives
$a_{1}Y_{1j}+a_{1}bY_{2j}=\dd_{1j}$, whence $Y_{11}=1/a_{1}$,
$Y_{12}=-b/a_{2}$ and $Y_{13}=0$. Thus
\[
   (F^{-1}_{Ai})=\begin{pmatrix}1/a_{1}&-b/a_{2}&0\\
   0&1/a_{2}&0\\ 0&0&1/a_{3}\end{pmatrix},
\]
which one verifies directly: the $(1,1)$ entry of $\bF\bF^{-1}$ is
$a_{1}(1/a_{1})+a_{1}b\cdot0=1$, the $(1,2)$ entry is
$a_{1}(-b/a_{2})+a_{1}b(1/a_{2})=-a_{1}b/a_{2}+a_{1}b/a_{2}=0$, and
the remaining entries are immediate.

\medskip
\noindent\emph{The cofactor.}
By \eqr{2.55}, $\bF^{*}=J\bF^{-t}$, so transposing the inverse and
multiplying every entry by $J=a_{1}a_{2}a_{3}$,
\[
   (F^{*}_{iA})=a_{1}a_{2}a_{3}\begin{pmatrix}1/a_{1}&0&0\\
   -b/a_{2}&1/a_{2}&0\\ 0&0&1/a_{3}\end{pmatrix}
   =\begin{pmatrix}a_{2}a_{3}&0&0\\
   -a_{1}a_{3}b&a_{1}a_{3}&0\\ 0&0&a_{1}a_{2}\end{pmatrix},
\]
entry by entry: $a_{1}a_{2}a_{3}/a_{1}=a_{2}a_{3}$;
$a_{1}a_{2}a_{3}\cdot(-b/a_{2})=-a_{1}a_{3}b$;
$a_{1}a_{2}a_{3}/a_{2}=a_{1}a_{3}$; and
$a_{1}a_{2}a_{3}/a_{3}=a_{1}a_{2}$.

\medskip
\noindent\emph{The Piola stress.}
By \eqr{6.140}, $P_{iA}=T_{ij}F^{*}_{jA}$. Because $\bT$ has only two
non-zero entries, $T_{12}=T_{21}=\tau$, the product merely selects
rows of $\bF^{*}$:
\[
\begin{aligned}
   P_{1A}&=T_{11}F^{*}_{1A}+T_{12}F^{*}_{2A}+T_{13}F^{*}_{3A}
    =\tau F^{*}_{2A},\\
   P_{2A}&=T_{21}F^{*}_{1A}+T_{22}F^{*}_{2A}+T_{23}F^{*}_{3A}
    =\tau F^{*}_{1A},\\
   P_{3A}&=0 .
\end{aligned}
\]
So the first row of $\bP$ is $\tau$ times the second row of $\bF^{*}$,
the second row of $\bP$ is $\tau$ times the first row of $\bF^{*}$, and
the third vanishes:
\[
   \boxed{\;(P_{iA})=\begin{pmatrix}
   -\tau a_{1}a_{3}b&\ \tau a_{1}a_{3}&\ 0\\
   \tau a_{2}a_{3}&0&0\\ 0&0&0\end{pmatrix}.\;}
\]
This is not symmetric and must not be, $\bP$ having one referential and
one spatial leg, so that its transpose is a different kind of object
altogether; that is Remark~\ref{rem:ch6Pnotsym}. What holds instead is
\eqr{6.143}, and since it is the only available check it is worth
carrying out in full. With $M_{ij}=P_{iA}F_{jA}$, pairing rows of $\bP$
with rows of $\bF$,
\[
\begin{aligned}
   M_{11}&=(-\tau a_{1}a_{3}b)a_{1}+(\tau a_{1}a_{3})(a_{1}b)+0
     =-\tau a_{1}^{2}a_{3}b+\tau a_{1}^{2}a_{3}b=0,\\
   M_{12}&=(-\tau a_{1}a_{3}b)\cdot0+(\tau a_{1}a_{3})a_{2}+0
     =\tau a_{1}a_{2}a_{3},\\
   M_{13}&=(-\tau a_{1}a_{3}b)\cdot0+(\tau a_{1}a_{3})\cdot0
     +0\cdot a_{3}=0,\\
   M_{21}&=(\tau a_{2}a_{3})a_{1}+0\cdot(a_{1}b)+0
     =\tau a_{1}a_{2}a_{3},\\
   M_{22}&=(\tau a_{2}a_{3})\cdot0+0\cdot a_{2}+0=0,\\
   M_{23}&=0,\qquad M_{3j}=0 .
\end{aligned}
\]
Hence $\bM=\tau a_{1}a_{2}a_{3}\big(\be_{1}\otimes\be_{2}
+\be_{2}\otimes\be_{1}\big)=J\bT$, symmetric, as \eqr{6.140} and the
symmetry of $\bT$ together require.

\medskip
\noindent\emph{The Piola--Kirchhoff stress.}
By \eqr{6.144}, $\bS=\bF^{-1}\bP$, that is
$S_{AB}=F^{-1}_{Ai}P_{iB}$, pairing the $A$th row of $\bF^{-1}$ with
the $B$th column of $\bP$:
\[
   \boxed{\;(S_{AB})=\begin{pmatrix}
   -2\tau a_{3}b&\ \tau a_{3}&\ 0\\
   \tau a_{3}&0&0\\ 0&0&0\end{pmatrix},\;}
\]
and this one is symmetric, as \eqr{6.145} guarantees. That $S_{12}$ and
$S_{21}$ arrive by entirely different routes and agree is the
arithmetic check on the whole calculation.

% reader's note: the product F^{-1}P, all nine entries
\rnote{ch6:prob7-product}

\medskip
\noindent\emph{Reading the answer.}
In $\kappa_{t}$ the state is a pure shear: no normal component on any
coordinate plane, $T_{11}=T_{22}=T_{33}=0$. Neither referential measure
inherits that. Tracking the $11$ entry through the three tensors,
\[
   T_{11}=0
   \quad\longrightarrow\quad
   P_{11}=-\tau a_{1}a_{3}b
   \quad\longrightarrow\quad
   S_{11}=-2\tau a_{3}b ,
\]
so $\bP$ picks up a normal entry from the tilt of the area element, and
$\bS$ carries twice what $\bP$ does, because $\bF^{*}$ tilts the area
element once and $\bF^{-1}$ tilts the surviving leg the same way again.
The stress did not change; the surfaces on which one reckons it did,
and the shear $b$ is what tilts them. Setting $b=0$ makes
$F^{*}_{21}=0$ and returns both to pure shear form; setting
$a_{1}=a_{2}=a_{3}=1$ with $b=0$ makes $\bF=\I$, hence
$\bF^{*}=\I$ and $J=1$, and collapses all three tensors onto each
other, which is Remark~\ref{rem:ch6threestress} at its trivial end.
\end{solution}

\begin{problem}[The two Piola transforms, by direct calculation]
Establish \eqr{6.138}, $J\dvg\bT=\Dvg\bP$, and \eqr{6.175},
$J\dvg\bq=\Dvg(J\bF^{-1}\bq)$, by direct calculation with the aid of the
Piola identity.
\end{problem}

\begin{solution}
Both identities are the same statement, once for a tensor and once for
a vector, and both rest on two facts about $\bF^{*}$.

The first is the Piola identity $\Dvg\bF^{*}=\bze$, that is
$F^{*}_{iA,A}=0$, proved as Problem 2(b) of Chapter 2 and revisited in
Section~\ref{sec:piolaid}. The second is the contraction
\begin{equation}
   F_{kA}F^{*}_{jA}=J\,\dd_{kj} ,
   \tag{$\ast_{9}$}\label{eq:ch6FFstar}
\end{equation}
which deserves a line of its own. Transposing \eqr{2.55},
$\bF^{*}=J\bF^{-t}$, gives
\[
   \big(\bF^{*}\big)^{t}=\big(J\bF^{-t}\big)^{t}
   =J\Big(\big(\bF^{-1}\big)^{t}\Big)^{t}=J\bF^{-1} ,
\]
so that $\bF\big(\bF^{*}\big)^{t}=\bF\big(J\bF^{-1}\big)
=J\bF\bF^{-1}=J\I$. In components,
$\big(\bF(\bF^{*})^{t}\big)_{kj}
=F_{kA}\big((\bF^{*})^{t}\big)_{Aj}=F_{kA}F^{*}_{jA}$, which is
\eqref{eq:ch6FFstar}. The first fact kills a term; the second collapses
the two-point machinery to a Kronecker delta. Nothing else is needed.

\medskip
\noindent\emph{The stress.}
With $P_{iA}=T_{ij}F^{*}_{jA}$ from \eqr{6.141}, the referential
divergence differentiates a product of two $\bX$-dependent factors:
\[
   P_{iA,A}=\frac{\partial}{\partial X_{A}}
   \big(T_{ij}F^{*}_{jA}\big)
   =\frac{\partial T_{ij}}{\partial X_{A}}F^{*}_{jA}
   +T_{ij}\,\frac{\partial F^{*}_{jA}}{\partial X_{A}}
   =\frac{\partial T_{ij}}{\partial X_{A}}F^{*}_{jA}
   +T_{ij}\underbrace{F^{*}_{jA,A}}_{=\,0} ,
\]
the second term vanishing by the Piola identity. In the first, $\bT$ is
a \emph{spatial} field, $T_{ij}=T_{ij}(\bx,t)$, and a referential
coordinate reaches it only through the motion
$\bx=\bchi_{\kappa}(\bX,t)$. The chain rule, exactly as at \eqr{2.37},
therefore supplies a factor of $\bF$:
\[
   \frac{\partial T_{ij}}{\partial X_{A}}
   =\frac{\partial T_{ij}}{\partial x_{k}}
    \frac{\partial\chi_{k}}{\partial X_{A}}
   =T_{ij,k}F_{kA} ,
\]
summed on $k$. That factor is exactly what \eqref{eq:ch6FFstar} is
waiting for:
\[
   P_{iA,A}=T_{ij,k}\,F_{kA}F^{*}_{jA}
   =T_{ij,k}\,J\dd_{kj}
   =J\,T_{ij,j} ,
\]
the last step because the delta identifies $k$ with $j$. In direct
notation this is $\Dvg\bP=J\dvg\bT$, which is \eqr{6.138}.

% reader's note: the three moves in one line
\rnote{ch6:prob8-moves}

\medskip
\noindent\emph{The heat flux.}
The same argument runs with one index deleted. By \eqr{6.172} the
referential flux has components $q_{(\kappa)A}=F^{*}_{iA}q_{i}$, so
\[
\begin{aligned}
   q_{(\kappa)A,A}
   &=\frac{\partial}{\partial X_{A}}\big(F^{*}_{iA}q_{i}\big)
    =\underbrace{F^{*}_{iA,A}}_{=\,0}q_{i}
     +F^{*}_{iA}\frac{\partial q_{i}}{\partial X_{A}}\\
   &=F^{*}_{iA}\,q_{i,k}F_{kA}
    =q_{i,k}\,F_{kA}F^{*}_{iA}\\
   &=q_{i,k}\,J\dd_{ki}
    =J\,q_{i,i} ,
\end{aligned}
\]
using the Piola identity, then the chain rule, then
\eqref{eq:ch6FFstar} with $j$ replaced by $i$. That is
$\Dvg\bq_{\kappa}=J\dvg\bq$, which is \eqr{6.175} once
$\bq_{\kappa}=(\bF^{*})^{t}\bq=J\bF^{-1}\bq$ is recognised from
\eqr{6.171}.

\medskip
\noindent\emph{What has been proved.}
A single fact, applied twice: multiplying a spatial field by $\bF^{*}$
and taking the referential divergence reproduces $J$ times the spatial
divergence, whatever the rank of the field. The Piola identity is what
makes it work, and Remark~\ref{rem:piolauses} of Chapter 2 predicted
this use.

Geometrically it is Piola--Nanson differentiated. The identity
$\bn\,\dl a=\bF^{*}\bN\,\dl A$ says $\bF^{*}$ converts area elements,
hence flux integrals; the divergence theorem converts those into volume
integrals on either side; and $\dl v=J\,\dl V$ supplies the factor
relating the two volume integrals. Chasing that around,
\[
   \int_{P}J\dvg\bq\,\dl V=\int_{P_{t}}\dvg\bq\,\dl v
   =\int_{\partial P_{t}}\ip\bq\bn\,\dl a
   =\int_{\partial P}\ip{\bq_{\kappa}}\bN\,\dl A
   =\int_{P}\Dvg\bq_{\kappa}\,\dl V ,
\]
and localizing on the arbitrary $P$ gives \eqr{6.175} again, this time
without differentiating anything.

It is worth naming what would go wrong without $\Dvg\bF^{*}=\bze$. The
dropped terms were $T_{ij}F^{*}_{jA,A}$ and $q_{i}F^{*}_{iA,A}$, so in
general one would have
\[
   \Dvg\bP=J\dvg\bT+\bT\,\Dvg\bF^{*} ,
   \qquad
   \Dvg\bq_{\kappa}=J\dvg\bq+\ip{\bq}{\Dvg\bF^{*}} ,
\]
and since $F^{*}_{iA}$ is built from first derivatives of the motion,
$F^{*}_{iA,A}$ involves second derivatives. The spatial and referential
balances would then differ by a term carrying second derivatives of
$\bchi_{\kappa}$ --- a quantity with no physical name and no place in
either formulation. That the term vanishes identically, for every
motion, is what makes the correspondence between the two descriptions
term-by-term rather than merely overall.
\end{solution}

\begin{problem}[The virtual power equation, or weak form]
Suppose a body is in equilibrium in the configuration $\kappa_{t}$.
Derive the virtual power equation
\[
   \int_{P_{t}}T_{ij}u_{i,j}\,\dl v
   =\int_{P_{t}}\rho b_{i}u_{i}\,\dl v
   +\int_{\partial P_{t}}t_{i}u_{i}\,\dl a ,
\]
where $P_{t}\subset\kappa_{t}$ is any part and $\bu$ is any
differentiable vector field whatsoever.
\end{problem}

\begin{solution}
Equilibrium means $\dot\bv=\bze$, so \eqr{6.110} reduces to
\begin{equation}
   T_{ij,j}+\rho b_{i}=0 .
   \tag{$\ast_{10}$}\label{eq:ch6equil}
\end{equation}
The identity to be proved has a boundary integral on the right and none
on the left, so the natural direction of attack is to start at the
boundary and work inward, converting the one term that can be
converted.

By \eqr{6.94}, $t_{i}=T_{ij}n_{j}$, so for each fixed $i$ the surface
integrand $t_{i}u_{i}$ --- summed on $i$ --- is
\[
   t_{i}u_{i}=T_{ij}n_{j}u_{i}=\big(T_{ij}u_{i}\big)n_{j} ,
\]
the normal component of the vector field with components
$v_{j}:=T_{ij}u_{i}$, summed on $i$. The divergence theorem,
Theorem~\ref{thm:div}, applies to it --- both $\bT$ and $\bu$ being
differentiable, so $v_{j}$ is --- and gives
\[
   \int_{\partial P_{t}}t_{i}u_{i}\,\dl a
   =\int_{\partial P_{t}}\big(T_{ij}u_{i}\big)n_{j}\,\dl a
   =\int_{P_{t}}\big(T_{ij}u_{i}\big)_{,j}\,\dl v .
\]
Expanding the derivative by the product rule produces exactly two
terms, one of which \eqref{eq:ch6equil} removes:
\[
   \big(T_{ij}u_{i}\big)_{,j}=T_{ij,j}u_{i}+T_{ij}u_{i,j}
   =\big(-\rho b_{i}\big)u_{i}+T_{ij}u_{i,j}
   =-\rho b_{i}u_{i}+T_{ij}u_{i,j} .
\]
Integrating over $P_{t}$,
\[
   \int_{\partial P_{t}}t_{i}u_{i}\,\dl a
   =-\int_{P_{t}}\rho b_{i}u_{i}\,\dl v
   +\int_{P_{t}}T_{ij}u_{i,j}\,\dl v ,
\]
and carrying the body-force term across,
\[
   \boxed{\;\int_{P_{t}}T_{ij}u_{i,j}\,\dl v
   =\int_{P_{t}}\rho b_{i}u_{i}\,\dl v
   +\int_{\partial P_{t}}t_{i}u_{i}\,\dl a\;}
\]
for every differentiable $\bu$ and every part $P_{t}$. The field $\bu$
was never required to be a velocity, or to satisfy any equation at all,
which is what makes the identity useful: one is free to choose it.

\medskip
\noindent\emph{Choosing $\bu$ constant.}
Take $\bu=\bc$ with $\bc$ a fixed vector. Then $u_{i,j}=c_{i,j}=0$, so
the left side vanishes and
\[
   0=\int_{P_{t}}\rho b_{i}c_{i}\,\dl v
   +\int_{\partial P_{t}}t_{i}c_{i}\,\dl a
   =c_{i}\left[\int_{P_{t}}\rho b_{i}\,\dl v
   +\int_{\partial P_{t}}t_{i}\,\dl a\right],
\]
the constants $c_{i}$ coming out of both integrals. In direct notation,
\[
   \ip{\bc}{\Big(\int_{P_{t}}\rho\bb\,\dl v
   +\int_{\partial P_{t}}\bt\,\dl a\Big)}=0
   \quad\text{for every }\bc ,
\]
so the bracket vanishes by Lemma~\ref{lem:cvec}, and by \eqr{6.30} that
bracket is $\vek{f}(S,t)$. This is \eqr{6.34} in equilibrium.

\medskip
\noindent\emph{Choosing $\bu$ a rigid rotation field.}
Take $\bu=\bw\times\bx$ with $\bw$ fixed. By
Proposition~\ref{prop:axial}, $\grad\bu$ is the skew tensor with axial
vector $\bw$; explicitly $u_{i}=\eps_{ipq}w_{p}x_{q}$, so
\[
   u_{i,j}=\eps_{ipq}w_{p}x_{q,j}=\eps_{ipq}w_{p}\dd_{qj}
   =\eps_{ipj}w_{p} ,
\]
which is antisymmetric in $i$ and $j$ because $\eps_{ipj}$ is. The left
side vanishes again, but for a different reason than before: not
because the gradient is zero but because
\[
   T_{ij}u_{i,j}=T_{ij}\,\eps_{ipj}w_{p}=0 ,
\]
the contraction of the symmetric $T_{ij}$ with the skew $\eps_{ipj}$
over the pair $(i,j)$, which vanishes by Lemma~\ref{lem:symskw}. What
remains is
\[
   0=\int_{P_{t}}\rho b_{i}u_{i}\,\dl v
   +\int_{\partial P_{t}}t_{i}u_{i}\,\dl a
   =\int_{P_{t}}\rho\,\ip\bb{\bw\times\bx}\,\dl v
   +\int_{\partial P_{t}}\ip\bt{\bw\times\bx}\,\dl a ,
\]
and the scalar triple product identity
$\ip\bb{\bw\times\bx}=\ip\bw{\bx\times\bb}$ --- both sides being the
determinant $[\bw,\bx,\bb]$ up to the same cyclic permutation ---
converts each term:
\[
   \ip{\bw}{\Big(\int_{P_{t}}\rho\,\bx\times\bb\,\dl v
   +\int_{\partial P_{t}}\bx\times\bt\,\dl a\Big)}=0
   \quad\text{for every }\bw ,
\]
so the bracket vanishes by Lemma~\ref{lem:cvec} again, and by
\eqr{6.31}--\eqr{6.33} it is $\bm_{\bo}(S,t)$. This is \eqr{6.35} in
equilibrium. One scalar equation, quantified over all $\bu$, therefore
carries both of Euler's postulates, and the two are recovered by the two
simplest choices of $\bu$ available.

\medskip
\noindent\emph{What the left side sees.}
The second computation generalises. In direct notation the left side is
$\int_{P_{t}}\bT\cdot\grad\bu\,\dl v$, and splitting $\grad\bu$ into
symmetric and skew parts as in \eqr{4.7},
\[
   \bT\cdot\grad\bu
   =\bT\cdot\Sym\grad\bu+\bT\cdot\Skw\grad\bu
   =\bT\cdot\Sym\grad\bu ,
\]
the second term vanishing by Lemma~\ref{lem:symskw} since $\bT$ is
symmetric. So a virtual field with skew gradient contributes nothing,
and by \eqr{4.31} those are exactly the rigid virtual velocities
$\bu=\bw\times\bx+\bd$. The weak form does not see them, which is the
statement that a body in equilibrium does no work in a rigid virtual
motion.

\medskip
\noindent\emph{Why it is called weak.}
The derivation used $T_{ij,j}$, but the result does not: the boxed
identity contains only $\bT$ itself and the first derivatives of $\bu$.
It therefore continues to make sense for stress fields that are merely
integrable, whereas the local equation \eqr{6.110} demands that $\bT$
be differentiable and is correspondingly called the strong form. The
logical relation is that \eqr{6.110} implies the boxed identity, as
just shown, and conversely the boxed identity for all $\bu$ implies
\eqr{6.110} wherever $\bT$ is differentiable --- run the argument
backwards, integrate by parts the other way, and localize. Every finite
element computation in solid mechanics solves the boxed equation
restricted to a finite-dimensional space of fields $\bu$, and it does so
precisely because the weaker demand on $\bT$ admits approximations that
are not differentiable across element boundaries.
\end{solution}

\bookfig{13}\begin{figure}[htbp]\centering
\begin{tikzpicture}[scale=1.0,
   ar/.style={-{Stealth[length=2.2mm]},thick},
   dim/.style={{Stealth[length=1.7mm]}-{Stealth[length=1.7mm]},
               semithick,gray!75}]
% the block, clamped on the left face
\fill[gray!8] (0,0) rectangle (5.2,1.6);
\draw[thick] (0,0) rectangle (5.2,1.6);
\fill[pattern=north east lines,pattern color=gray!65]
   (-0.36,-0.12) rectangle (-0.02,1.72);
\draw[thick] (0,-0.12)--(0,1.72);
% the pressure on l1 < x1 < l2
\foreach \x in {1.80,2.16,2.52,2.88,3.24,3.60}{
   \draw[ar,gray!85] (\x,2.16)--(\x,1.64);}
\draw[semithick,gray!85] (1.80,2.16)--(3.60,2.16);
\node[font=\small,anchor=south] at (2.70,2.22) {$p$};
% origin
\fill (0,0) circle (1.4pt);
\node[font=\small,anchor=north east] at (-0.04,-0.04) {$O$};
% dimensions
\draw[dim] (0,-0.62)--(5.2,-0.62);
\node[font=\small,fill=white,inner sep=1pt] at (2.6,-0.62) {$w$};
\draw[dim] (5.62,0)--(5.62,1.6);
\node[font=\small,fill=white,inner sep=1pt] at (5.62,0.8) {$h$};
\draw[guide] (1.80,1.64)--(1.80,2.72);
\draw[guide] (3.60,1.64)--(3.60,3.16);
\draw[guide] (0,1.72)--(0,3.16);
\draw[dim] (0,2.62)--(1.80,2.62);
\node[font=\small,fill=white,inner sep=1pt] at (0.90,2.62) {$l_{1}$};
\draw[dim] (0,3.06)--(3.60,3.06);
\node[font=\small,fill=white,inner sep=1pt] at (1.80,3.06) {$l_{2}$};
\end{tikzpicture}
\caption{The block of Problem 10(d), of width $w$ and height $h$,
clamped on the face $x_{1}=0$ and carrying a uniform pressure $p$ on the
part $l_{1}<x_{1}<l_{2}$ of its top surface. The origin is the lower
corner of the clamped face, $x_{1}$ runs along the length and $x_{2}$
upward, and the thickness in the $x_{3}$ direction is taken as unity.
The reactions on the clamped face are unknown, which is why the
calculation goes through $\bar T_{21}$ and uses the symmetry of the mean
stress to recover $\bar T_{12}$.}
\label{fig:ch6block}
\end{figure}

\begin{problem}[The mean stress in a body at rest]
Let $V(\kappa_{t})$ be the volume of a configuration $\kappa_{t}$ and
define the mean Cauchy stress
\[
   \bar\bT=\big[V(\kappa_{t})\big]^{-1}\int_{\kappa_{t}}\bT\,\dl v .
\]
\textup{(a)} If the body is in equilibrium, show that
\[
   V(\kappa_{t})\,\bar T_{ij}=\int_{\partial\kappa_{t}}t_{i}x_{j}\,\dl a
   +\int_{\kappa_{t}}\rho b_{i}x_{j}\,\dl v .
\]
\textup{(b)} Apply the global rotational momentum balance to show that
$e_{kij}\bar T_{ij}=0$ in equilibrium.
\textup{(c)} Suppose the body rests on a rigid plane under gravity
$\bb=-g\be_{3}$, with no applied tractions other than those due to
contact. Show that $\bar T_{33}=-\bar\rho g\bar x_{3}$, where $\bar\rho$
is the mean density and $\bar x_{3}=\ip{\bar\bx}{\be_{3}}$ with the
origin on the plane. The result appears to depend on the choice of
origin; show that it does not.
\textup{(d)} For the clamped block of Figure~\ref{fig:ch6block},
carrying a uniform pressure $p$ on the interval $l_{1}<x_{1}<l_{2}$ of
its top surface, evaluate $\bar T_{12}$.
\end{problem}

\begin{solution}
The mean stress is a volume average and the data are boundary data, so
the problem is to manufacture a volume integral of $\bT$ out of surface
integrals. The device is that the position field differentiates to a
constant, $x_{j,k}=\dd_{jk}$ by Proposition~\ref{prop:posvec}, so
inserting $x_{j}$ under a divergence produces $\bT$ itself as one of
the two terms of the product rule.

\medskip
\noindent\textbf{(a)}
Fix $j$ and apply the divergence theorem to the vector field with
components $T_{ik}x_{j}$, $k$ being the index summed against the
normal:
\[
\begin{aligned}
   \int_{\partial\kappa_{t}}t_{i}x_{j}\,\dl a
   &=\int_{\partial\kappa_{t}}T_{ik}n_{k}x_{j}\,\dl a
    =\int_{\partial\kappa_{t}}\big(T_{ik}x_{j}\big)n_{k}\,\dl a
    =\int_{\kappa_{t}}\big(T_{ik}x_{j}\big)_{,k}\,\dl v\\
   &=\int_{\kappa_{t}}\big(T_{ik,k}x_{j}+T_{ik}x_{j,k}\big)\,\dl v
    =\int_{\kappa_{t}}\big(T_{ik,k}x_{j}+T_{ik}\dd_{jk}\big)\,\dl v ,
\end{aligned}
\]
the first equality by \eqr{6.94}. Now $T_{ik}\dd_{jk}=T_{ij}$, and in
equilibrium $T_{ik,k}=-\rho b_{i}$ by \eqref{eq:ch6equil}, so
\[
   \int_{\partial\kappa_{t}}t_{i}x_{j}\,\dl a
   =-\int_{\kappa_{t}}\rho b_{i}x_{j}\,\dl v
   +\int_{\kappa_{t}}T_{ij}\,\dl v .
\]
The last integral is $V(\kappa_{t})\bar T_{ij}$ by the definition of
the mean, so moving the body-force term across,
\[
   \boxed{\ V(\kappa_{t})\,\bar T_{ij}
   =\int_{\partial\kappa_{t}}t_{i}x_{j}\,\dl a
   +\int_{\kappa_{t}}\rho b_{i}x_{j}\,\dl v .\ }
\]
This is Problem 9 with the particular choice $u_{i}=c_{i}x_{j}$, one
index $j$ at a time: then $u_{i,k}=c_{i}\dd_{jk}$ and the left side of
the virtual power equation is $c_{i}\int T_{ij}\,\dl v$. The weak form
was general enough to contain this.

\medskip
\noindent\textbf{(b)}
In equilibrium the global rotational balance \eqr{6.37} reads
\[
   \int_{\kappa_{t}}\rho\,\bx\times\bb\,\dl v
   +\int_{\partial\kappa_{t}}\bx\times\bt\,\dl a=\bze .
\]
Its $k$th component uses $(\ba\times\bc)_{k}=e_{kji}a_{j}c_{i}$, so
\[
   e_{kji}\int_{\kappa_{t}}\rho\,x_{j}b_{i}\,\dl v
   +e_{kji}\int_{\partial\kappa_{t}}x_{j}t_{i}\,\dl a=0 ,
\]
that is, the constants $e_{kji}$ coming outside,
\[
   e_{kji}\Big[\int_{\kappa_{t}}\rho b_{i}x_{j}\,\dl v
   +\int_{\partial\kappa_{t}}t_{i}x_{j}\,\dl a\Big]=0 .
\]
The bracket is $V\bar T_{ij}$ by (a), and $V>0$, so
\[
   e_{kji}\bar T_{ij}=0,
   \qquad\text{equivalently}\qquad
   e_{kij}\bar T_{ij}=0 ,
\]
the two differing only by the sign $e_{kji}=-e_{kij}$. By Problem 5 of
Chapter 1 a tensor annihilated in this way by the permutation symbol is
symmetric --- writing the three components out,
$e_{1ji}\bar T_{ij}=\bar T_{32}-\bar T_{23}$ and cyclically --- so
$\bar\bT=\bar\bT^{t}$. This is of course the average of the pointwise
symmetry \eqr{6.117}, but the content of the exercise is that the
global balance alone delivers it, without ever passing to a point.

\medskip
\noindent\textbf{(c)}
Put $i=j=3$ in (a). The plane is $x_{3}=0$, and by hypothesis the only
tractions act where the body touches it, so $x_{3}$ vanishes on the
whole support of $\bt$ and the surface integral vanishes term by term.
With $b_{3}=-g$ constant,
\[
   V\bar T_{33}=\int_{\kappa_{t}}\rho(-g)x_{3}\,\dl v
   =-g\int_{\kappa_{t}}\rho x_{3}\,\dl v
   =-g\,m(S)\,\bar x_{3} ,
\]
the last step by \eqr{6.40}, whose third component reads
$m(S)\bar x_{3}=\int_{\kappa_{t}}\rho x_{3}\,\dl v$. Writing
$\bar\rho=m(S)/V$ for the mean density,
\[
   \boxed{\;\bar T_{33}=-\bar\rho\,g\,\bar x_{3} .\;}
\]

The apparent dependence on the origin is not real. The statement fixes
the origin on the plane but not where on it, so the admissible origins
differ by translations \emph{within} the plane, $\bo\mapsto\bo+\bc$
with $\ip\bc{\be_{3}}=0$. Such a translation sends
$x_{3}=\ip\bx{\be_{3}}$ to $\ip{(\bx-\bc)}{\be_{3}}
=x_{3}-\ip\bc{\be_{3}}=x_{3}$, so $x_{3}$ and hence $\bar x_{3}$ are
untouched. Thus $\bar x_{3}$ is the height of the mass center above the
plane, a geometric quantity, and the formula is unambiguous. Moving the
origin off the plane would change $\bar x_{3}$ and would also
invalidate the derivation, since the contact tractions would no longer
sit at $x_{3}=0$ and the surface integral would not vanish.

\medskip
\noindent\textbf{(d)}
The block occupies $0\le x_{1}\le w$, $0\le x_{2}\le h$ with unit
thickness in $x_{3}$, so $V=wh$, and $\bb=\bze$.

Computing $\bar T_{12}$ from (a) would need $\int t_{1}x_{2}\,\dl a$
over the clamped face $x_{1}=0$, where the reactions $t_{1}$ are
unknown; that route is blocked. But $\bar\bT$ is symmetric by (b), so
compute $\bar T_{21}$ instead:
\[
   V\bar T_{21}=\int_{\partial\kappa_{t}}t_{2}x_{1}\,\dl a .
\]
Now every contribution is known. Going round the boundary:
on the clamped face $x_{1}=0$, so the integrand carries the factor
$x_{1}=0$ and the unknown reactions drop out --- the whole point of
choosing this component; on the bottom $x_{2}=0$, on the free end
$x_{1}=w$, and on the unloaded part of the top the surface is
traction-free, so $\bt=\bze$ there and each contributes nothing; and on
the loaded part of the top the outward normal is $\be_{2}$ and the
pressure acts inward, so $\bt=-p\be_{2}$, giving $t_{2}=-p$, with
$x_{1}$ running from $l_{1}$ to $l_{2}$ and $\dl a=\dl x_{1}$ for unit
thickness. Hence
\[
   V\bar T_{21}=\int_{l_{1}}^{l_{2}}(-p)\,x_{1}\,\dl x_{1}
   =-p\left[\frac{x_{1}^{2}}{2}\right]_{l_{1}}^{l_{2}}
   =-\frac{p}{2}\big(l_{2}^{2}-l_{1}^{2}\big) ,
\]
and dividing by $V=wh$,
\[
   \boxed{\;\bar T_{12}=\bar T_{21}
   =-\frac{p\big(l_{2}^{2}-l_{1}^{2}\big)}{2wh} .\;}
\]
The negative sign says the mean shear opposes the applied load, and the
factor $l_{2}^{2}-l_{1}^{2}$, growing with the distance from the clamp,
says that pressure applied further out counts for more --- both as one
expects of a cantilever. The computation is a small demonstration of
what \eqr{6.117} is worth: it converted an integral over an unknown
boundary into one over a known boundary.
\end{solution}

\begin{problem}[The energy flux and the jump conditions it carries]
All the scalar balance laws considered have the generic form \eqr{6.1}.
Take $\varphi$ to be the total energy per unit mass,
$\varphi=\varepsilon+\tfrac12|\bv|^{2}$. What is the energy flux
$\vek\pi$? Obtain the relevant jump condition and compare it with
\eqr{6.169}. Thus derive the jump condition for \eqr{6.150}.
\end{problem}

\begin{solution}
\noindent\emph{The flux.}
The balance of total energy is \eqr{6.161}, and to read a flux off it
one has to write every term as an integral and then force the surface
integrand into the shape $\ip{\vek\pi}\bn$ that \eqr{6.1} demands.
Using \eqr{6.151} and \eqr{6.162} on the left and \eqr{6.152} and
\eqr{6.163} on the right,
\[
   \frac{\dl}{\dl t}\int_{P_{t}}\rho
   \big(\varepsilon+\tfrac12|\bv|^{2}\big)\,\dl v
   =\int_{P_{t}}\rho\big(\ip\bb\bv+r\big)\,\dl v
   +\int_{\partial P_{t}}\big(\ip\bt\bv+h\big)\,\dl a ,
\]
the bulk terms having been collected over the common factor $\rho$ and
the surface terms over the common measure.

Of the two boundary contributions one is already in the required form:
\eqr{6.166} reads $h=-\ip\bq\bn=\ip{(-\bq)}\bn$. The other is not, and
making it so is the only step in the problem that uses anything.
Inserting $\bt=\bT\bn$ and moving $\bT$ across the inner product by the
definition of the transpose, then using the symmetry \eqr{6.117} to put
it back onto the other factor,
\[
   \ip\bt\bv=\ip{\bT\bn}\bv=\ip\bn{\bT^{t}\bv}
   =\ip{\bT^{t}\bv}\bn
   =\ip{\bT\bv}\bn ,
\]
the third equality by symmetry of the inner product and the fourth by
$\bT^{t}=\bT$. Hence the whole surface integrand is
\[
   \ip\bt\bv+h=\ip{\bT\bv}\bn+\ip{(-\bq)}\bn
   =\ip{\big(\bT\bv-\bq\big)}\bn ,
\]
and comparing with \eqr{6.1} term by term,
\[
   \varphi=\varepsilon+\tfrac12|\bv|^{2},\qquad
   \sigma=\ip\bb\bv+r,\qquad
   \boxed{\;\vek\pi=\bT\bv-\bq\;}
\]
The flux of total energy is the working of the traction less the heat
conducted out, which is what one would guess. What the derivation adds
is that the guess needs $\bT=\bT^{t}$: without symmetry the working is
$\ip\bn{\bT^{t}\bv}$, the normal component of $\bT^{t}\bv$ rather than
of $\bT\bv$, and while that is still a vector, the identification of
the flux would carry the transpose along with it into every later
formula.

\medskip
\noindent\emph{The jump condition.}
Substituting the three letters into \eqr{6.6},
$\ip{\jump{\rho\varphi(\bv-\bu)-\vek\pi}}\bn=0$, gives
\[
   \boxed{\;\jump{\rho\big(\varepsilon+\tfrac12|\bv|^{2}\big)(\bv-\bu)
   -\bT\bv+\bq}\cdot\bn=0\;}
\]
at every point of a singular surface, the two minus signs in
$-\vek\pi=-\bT\bv+\bq$ being those of the definition just made.

Set this beside \eqr{6.169}, which is what the thermal balance alone
carries:
\[
   \jump{\rho\varepsilon(\bv-\bu)+\bq}\cdot\bn=0 .
\]
The total-energy condition has two terms the thermal one lacks, namely
the convected kinetic energy $\rho\tfrac12|\bv|^{2}(\bv-\bu)$ and the
stress working $-\bT\bv$, and no others: the terms
$\rho\varepsilon(\bv-\bu)$ and $+\bq$ occur in both, with the same
signs.

\medskip
\noindent\emph{The mechanical jump condition.}
Subtracting is therefore legitimate --- both are equalities at the same
point of the same surface --- and everything common cancels:
\[
\begin{aligned}
   &\jump{\rho\big(\varepsilon+\tfrac12|\bv|^{2}\big)(\bv-\bu)
     -\bT\bv+\bq}\cdot\bn
    -\jump{\rho\varepsilon(\bv-\bu)+\bq}\cdot\bn\\
   &\qquad=\jump{\rho\varepsilon(\bv-\bu)
     +\rho\tfrac12|\bv|^{2}(\bv-\bu)-\bT\bv+\bq
     -\rho\varepsilon(\bv-\bu)-\bq}\cdot\bn\\
   &\qquad=\jump{\rho\tfrac12|\bv|^{2}(\bv-\bu)-\bT\bv}\cdot\bn ,
\end{aligned}
\]
the jump bracket being linear, so that
\[
   \boxed{\;\jump{\rho\tfrac12|\bv|^{2}(\bv-\bu)-\bT\bv}\cdot\bn=0 ,\;}
\]
which is the condition belonging to the mechanical energy balance
\eqr{6.150}.

The result may be checked head-on rather than by subtraction. Equation
\eqr{6.150} is of the form \eqr{6.1} with
\[
   \varphi=\tfrac12|\bv|^{2},
   \qquad
   \vek\pi=\bT\bv,
   \qquad
   \sigma=\ip\bb\bv-\rho^{-1}\bT\cdot\bL ,
\]
the last because the stress power appears on the left of \eqr{6.150}
and must be moved to the right as a supply with a sign change and a
factor $\rho^{-1}$, the generic $\sigma$ being per unit mass. Feeding
these into \eqr{6.6} reproduces the display.

\medskip
\noindent\emph{One observation.}
That last check makes the real point of the problem visible. The stress
power $\bT\cdot\bL$, which is the entire difference between the
mechanical and thermal balances in the bulk, appears in none of the
three jump conditions. It cannot: it entered as part of $\sigma$, and
$\sigma$ is a volume term. Collapsing a part onto a surface kills
volume terms and spares surface terms --- a pillbox of thickness $\eta$
and face area $A$ has volume $\eta A$ and surface $2A+O(\eta)$, so
dividing by $A$ and letting $\eta\to0$ removes everything carrying the
factor $\eta$ --- so only $\varphi$ and $\vek\pi$ survive into
\eqr{6.6}, however large $\sigma$ may be. This is the shrinking
argument used at \eqr{6.81} and quantified by
Lemma~\ref{lem:ch6cone}, and it is why the three balances --- total,
thermal, mechanical --- differ in the bulk by the stress power but
differ on a shock only by convected and worked terms.
\end{solution}

\begin{problem}[Cofactors, and two more power-conjugate pairs]
Let $\bA$ and $\bB$ be invertible tensors.
\textup{(a)} Show that $(\bA\bB)^{*}=\bA^{*}\bB^{*}$.
\textup{(b)} For invertible $\bQ$, show that $\bQ^{*}=\bQ$ if and only
if $\bQ$ is a rotation.
\textup{(c)} The referential \emph{Biot stress} is
$\vek\sigma=\Sym(\bP^{t}\bR)$, with $\bR$ the rotation of the polar
decomposition of $\bF$. Show that $\vek\sigma$ is power conjugate to the
right stretch $\bU$, that is, that the stress power per unit reference
volume is $\vek\sigma\cdot\dot\bU$.
\textup{(d)} The spatial \emph{Bell stress} is
$\vek\Sigma=\Sym(\bP\bR^{t})$. Derive an expression for the tensor
$\bA$ such that $\vek\Sigma\cdot\bA$ is the stress power per unit
reference volume.
\end{problem}

\begin{solution}
\noindent\textbf{(a)}
The cofactor is characterised by what it does to cross products,
\[
   \bA^{*}(\bu\times\bv)=\bA\bu\times\bA\bv
   \qquad\text{for all }\bu,\bv ,
\]
which is Definition~\ref{def:cof}. Multiplicativity follows by using
that characterisation three times in succession:
\[
\begin{aligned}
   (\bA\bB)^{*}(\bu\times\bv)
   &=(\bA\bB\bu)\times(\bA\bB\bv)
    =\bA\big(\bB\bu\big)\times\bA\big(\bB\bv\big)\\
   &=\bA^{*}\big[(\bB\bu)\times(\bB\bv)\big]
    =\bA^{*}\bB^{*}(\bu\times\bv) ,
\end{aligned}
\]
the third equality reading the defining property for $\bA$ with the
vectors $\bB\bu$ and $\bB\bv$, and the fourth reading it for $\bB$.
Since every vector of $\Vt$ is a cross product of two others --- given
$\bw\neq\bze$, pick any $\bp$ with $\ip\bp\bw=0$ and $|\bp|=1$ and set
$\bu=\bp$, $\bv=\bw\times\bp/|\bw|$, whence
$\bu\times\bv=\bw$ --- the two tensors $(\bA\bB)^{*}$ and
$\bA^{*}\bB^{*}$ agree on all of $\Vt$ and are equal.

\medskip
\noindent\textbf{(b)}
One direction is a computation. If $\bQ$ is a rotation then
$\det\bQ=1$ and $\bQ^{t}=\bQ^{-1}$, so $\bQ^{-t}=(\bQ^{-1})^{t}
=(\bQ^{t})^{t}=\bQ$, and
\[
   \bQ^{*}=(\det\bQ)\bQ^{-t}=1\cdot\bQ=\bQ .
\]

The other is a trick with determinants. Suppose $\bQ^{*}=\bQ$. Since
the cofactor of an invertible tensor is $(\det\bQ)\bQ^{-t}$, this reads
$(\det\bQ)\bQ^{-t}=\bQ$, and post-multiplying by $\bQ^{t}$,
\[
   (\det\bQ)\,\bQ^{-t}\bQ^{t}=\bQ\bQ^{t},
   \qquad\text{that is}\qquad
   \bQ\bQ^{t}=(\det\bQ)\,\I ,
\]
since $\bQ^{-t}\bQ^{t}=\I$. Take determinants of both sides. On the
left, $\det(\bQ\bQ^{t})=\det\bQ\,\det\bQ^{t}=(\det\bQ)^{2}$; on the
right, $\det\big((\det\bQ)\I\big)=(\det\bQ)^{3}$, since
$\det(\alpha\I)=\alpha^{3}$ in three dimensions. So
\[
   (\det\bQ)^{2}=(\det\bQ)^{3} .
\]
Invertibility gives $\det\bQ\neq0$, so dividing by $(\det\bQ)^{2}$
leaves $\det\bQ=1$, and substituting back into
$\bQ\bQ^{t}=(\det\bQ)\I$ gives $\bQ\bQ^{t}=\I$. Thus $\bQ$ is
orthogonal with determinant $+1$: a rotation. The dimension enters only
in that exponent, which is worth noticing --- the statement is
three-dimensional, and in $n$ dimensions the same argument gives
$(\det\bQ)^{2}=(\det\bQ)^{n}$.

\medskip
\noindent\textbf{(c)}
Parts (c) and (d) follow one pattern: differentiate a polar
decomposition, and dispose of the term carrying $\dot\bR$ by a symmetry
argument. By \eqr{6.156} the stress power per unit reference volume is
$\bP\cdot\dot\bF$, and differentiating $\bF=\bR\bU$ from \eqr{3.32},
\[
   \dot\bF=\dot\bR\bU+\bR\dot\bU ,
   \qquad\text{so}\qquad
   \bP\cdot\dot\bF=\bP\cdot(\dot\bR\bU)+\bP\cdot(\bR\dot\bU) .
\]

The second term is the one that survives. Using
$\bA\cdot\bB=\tr(\bA\bB^{t})$ from Definition~\ref{def:tip}, then
$(\bR\dot\bU)^{t}=\dot\bU^{t}\bR^{t}$, then cyclicity of the trace,
\[
   \bP\cdot(\bR\dot\bU)=\tr\big(\bP(\bR\dot\bU)^{t}\big)
   =\tr\big(\bP\dot\bU^{t}\bR^{t}\big)
   =\tr\big(\bR^{t}\bP\dot\bU^{t}\big)
   =(\bR^{t}\bP)\cdot\dot\bU ,
\]
in which $\bR^{t}\bP$ maps $\hVt$ to $\hVt$, both legs now referential:
$\bP\colon\hVt\to\Vt$ and $\bR^{t}\colon\Vt\to\hVt$. So the inner
product is one of referential tensors. Since $\bU$ is symmetric at
every instant, so is $\dot\bU$, and Lemma~\ref{lem:symskw} discards the
skew part of $\bR^{t}\bP$:
\[
\begin{aligned}
   (\bR^{t}\bP)\cdot\dot\bU
   &=\Big[\Sym\big(\bR^{t}\bP\big)+\Skw\big(\bR^{t}\bP\big)\Big]
     \cdot\dot\bU
    =\Sym\big(\bR^{t}\bP\big)\cdot\dot\bU\\
   &=\Sym\big(\bP^{t}\bR\big)\cdot\dot\bU
    =\vek\sigma\cdot\dot\bU ,
\end{aligned}
\]
the penultimate equality because $\Sym\bX=\Sym\bX^{t}$ and
$(\bR^{t}\bP)^{t}=\bP^{t}(\bR^{t})^{t}=\bP^{t}\bR$.

The first term must vanish, and the reason is the symmetry restriction
\eqr{6.143} in disguise. Put $\bG=\bR^{t}\bP$, so that $\bP=\bR\bG$;
then $\bF^{t}=(\bR\bU)^{t}=\bU\bR^{t}$ and $\bP^{t}=\bG^{t}\bR^{t}$,
so $\bP\bF^{t}=\bF\bP^{t}$ becomes
\[
   \bR\bG\bU\bR^{t}=\bR\bU\bG^{t}\bR^{t} .
\]
Cancelling the invertible $\bR$ on the left and $\bR^{t}$ on the right
--- multiply by $\bR^{t}$ in front and $\bR$ behind, using
$\bR^{t}\bR=\hI$ --- gives
\[
   \bG\bU=\bU\bG^{t} ,
\]
and the right-hand form says $\bG\bU$ is symmetric, since
$(\bG\bU)^{t}=\bU^{t}\bG^{t}=\bU\bG^{t}=\bG\bU$. Meanwhile
\[
\begin{aligned}
   \bP\cdot(\dot\bR\bU)
   &=\tr\big(\bP(\dot\bR\bU)^{t}\big)
    =\tr\big(\bP\bU^{t}\dot\bR^{t}\big)
    =\tr\big(\bP\bU\dot\bR^{t}\big)\\
   &=\tr\big(\bR\bG\bU\dot\bR^{t}\big)
    =\tr\big(\bG\bU\,\dot\bR^{t}\bR\big),
\end{aligned}
\]
using $\bU^{t}=\bU$ and then cyclicity. Now $\dot\bR^{t}\bR$ is skew:
differentiating $\bR^{t}\bR=\hI$ gives
\[
   \dot\bR^{t}\bR+\bR^{t}\dot\bR=\hOb ,
\]
and $\bR^{t}\dot\bR=(\dot\bR^{t}\bR)^{t}$, so
$\dot\bR^{t}\bR+(\dot\bR^{t}\bR)^{t}=\hOb$, which is skewness. A
symmetric tensor contracted with a skew one vanishes, by
Lemma~\ref{lem:symskw} once more, so the whole term is zero and
\[
   \boxed{\;\bP\cdot\dot\bF=\vek\sigma\cdot\dot\bU\;}
\]
The Biot stress is power conjugate to $\bU$, and no correction was
needed.

\medskip
\noindent\textbf{(d)}
The same pattern on the other factorisation $\bF=\bV\bR$, for which
$\dot\bF=\dot\bV\bR+\bV\dot\bR$. The first term gives, by the trace
manipulation above with $\bR$ now on the right,
\[
   \bP\cdot(\dot\bV\bR)=\tr\big(\bP(\dot\bV\bR)^{t}\big)
   =\tr\big(\bP\bR^{t}\dot\bV^{t}\big)
   =(\bP\bR^{t})\cdot\dot\bV
   =\vek\Sigma\cdot\dot\bV ,
\]
$\bP\bR^{t}$ being spatial --- $\bR^{t}\colon\Vt\to\hVt$ followed by
$\bP\colon\hVt\to\Vt$ --- and $\dot\bV$ symmetric, so that as in (c)
only $\Sym(\bP\bR^{t})=\vek\Sigma$ survives.

This time the remaining term does \emph{not} vanish. Put
$\bH=\bP\bR^{t}$, so $\bP=\bH\bR$, and let $\bOm_{R}=\dot\bR\bR^{t}$,
spatial and skew by the argument of \eqr{6.55}. Differentiating
$\bR\bR^{t}=\I$ gives $\dot\bR\bR^{t}+\bR\dot\bR^{t}=\Ob$, so
$\bR\dot\bR^{t}=-\bOm_{R}$, and
\[
\begin{aligned}
   \bP\cdot(\bV\dot\bR)
   &=\tr\big(\bP(\bV\dot\bR)^{t}\big)
    =\tr\big(\bP\dot\bR^{t}\bV^{t}\big)
    =\tr\big(\bP\dot\bR^{t}\bV\big)\\
   &=\tr\big(\bH\bR\dot\bR^{t}\bV\big)
    =-\tr\big(\bH\bOm_{R}\bV\big)
    =-\tr\big(\bV\bH\,\bOm_{R}\big) ,
\end{aligned}
\]
using $\bV^{t}=\bV$ and cyclicity at the last step: a leftover
proportional to the spin of the polar rotation.

The claim is that restoring it is exactly what converts $\dot\bV$ into
its co-rotational rate
\[
   \mathring\bV=\dot\bV-\bOm_{R}\bV+\bV\bOm_{R} ,
\]
formed as in \eqr{6.68} but about $\bR$ rather than the rigid $\bQ$. To
see it, compute the correction against $\vek\Sigma$. Since
$\bOm_{R}^{t}=-\bOm_{R}$ and $\bV^{t}=\bV$,
\[
   \big(\bOm_{R}\bV-\bV\bOm_{R}\big)^{t}
   =\bV^{t}\bOm_{R}^{t}-\bOm_{R}^{t}\bV^{t}
   =-\bV\bOm_{R}+\bOm_{R}\bV
   =\bOm_{R}\bV-\bV\bOm_{R} ,
\]
so the commutator is symmetric and
\[
\begin{aligned}
   \vek\Sigma\cdot\big(\bOm_{R}\bV-\bV\bOm_{R}\big)
   &=\tr\Big(\vek\Sigma\big(\bOm_{R}\bV-\bV\bOm_{R}\big)^{t}\Big)
    =\tr\big(\vek\Sigma\bOm_{R}\bV\big)
     -\tr\big(\vek\Sigma\bV\bOm_{R}\big)\\
   &=\tr\big(\bV\vek\Sigma\bOm_{R}\big)
     -\tr\big(\vek\Sigma\bV\bOm_{R}\big)
    =\tr\big[\big(\bV\vek\Sigma-\vek\Sigma\bV\big)\bOm_{R}\big],
\end{aligned}
\]
the third equality by cyclicity applied to the first trace alone.
Write $\bH=\vek\Sigma+\bW_{H}$ with $\bW_{H}=\Skw\bH$. The symmetry
restriction in the form $\bH\bV=\bV\bH^{t}$, obtained as in (c) but
with $\bF=\bV\bR$, reads
\[
   \big(\vek\Sigma+\bW_{H}\big)\bV=\bV\big(\vek\Sigma-\bW_{H}\big),
   \qquad\text{that is}\qquad
   \vek\Sigma\bV-\bV\vek\Sigma+\bV\bW_{H}+\bW_{H}\bV=\Ob ,
\]
using $\bH^{t}=\vek\Sigma-\bW_{H}$. This says precisely that
$\vek\Sigma\bV+\bV\bW_{H}$ is symmetric: its transpose is
$\bV\vek\Sigma+\bW_{H}^{t}\bV=\bV\vek\Sigma-\bW_{H}\bV$, and the
displayed identity is exactly the statement that the two agree. Its
contraction with the skew $\bOm_{R}$ therefore vanishes, leaving
\[
   \tr\big[\big(\bV\vek\Sigma-\vek\Sigma\bV\big)\bOm_{R}\big]
   =\tr\big(\bV\bH\,\bOm_{R}\big) ,
\]
which is precisely the leftover with its sign reversed. Assembling,
\[
\begin{aligned}
   \vek\Sigma\cdot\mathring\bV
   &=\vek\Sigma\cdot\dot\bV
     -\vek\Sigma\cdot\big(\bOm_{R}\bV-\bV\bOm_{R}\big)
    =\vek\Sigma\cdot\dot\bV-\tr\big(\bV\bH\bOm_{R}\big)\\
   &=\bP\cdot(\dot\bV\bR)+\bP\cdot(\bV\dot\bR)
    =\bP\cdot\dot\bF ,
\end{aligned}
\]
so that
\[
   \boxed{\;\bA=\mathring\bV=\dot\bV-\bOm_{R}\bV+\bV\bOm_{R},
   \qquad\bOm_{R}=\dot\bR\bR^{t} .\;}
\]

\medskip
\noindent\emph{Reading the pair.}
The two halves differ in one structural respect, and it accounts for
everything. Part (c) pairs a referential stress with the rate of a
referential stretch, and both live in a space that does not move; no
correction arises. Part (d) pairs a spatial stress with a spatial
stretch, and $\bV$ sits in a frame the material carries around with it,
so its bare rate $\dot\bV$ measures the turning of that frame along
with the stretching. The correction is the co-rotational derivative of
\eqr{6.68}, met first for the inertia tensor and promised in
Remark~\ref{rem:ch6corot}, and the rotation is $\bR$ rather than $\bQ$
because in a general deformation $\bR$ is what plays the part the rigid
rotation played there. Using $\dot\bV$ in place of $\mathring\bV$ gives
the wrong answer whenever $\bOm_{R}\bV\neq\bV\bOm_{R}$, which is the
generic case; the two commute only when $\bV$ and the axis of
$\bOm_{R}$ are compatibly aligned.
\end{solution}
